\documentclass[12pt,a4paper]{article}
\usepackage[utf8]{inputenc}
\usepackage[T2A]{fontenc}
\usepackage[russian]{babel}
\usepackage{amsmath, amssymb, amsthm}
\usepackage{geometry}
\usepackage{enumitem}
\usepackage{booktabs}
\usepackage{array}

\geometry{left=2.5cm, right=2cm, top=2cm, bottom=2cm}

\newtheoremstyle{solution}{}{}{}{}{\bfseries}{.}{ }{}
\theoremstyle{solution}
\newtheorem*{theory}{Теория}
\newtheorem*{idea}{Идея решения}

\newcommand{\task}[1]{\newpage\section*{Задача #1}}
\newcommand{\answer}{\paragraph{Ответ.}}
\newcommand{\step}[1]{\paragraph{Шаг #1.}}

\newcommand{\R}{\mathbb{R}}
\newcommand{\C}{\mathbb{C}}
\newcommand{\norm}[1]{\left\|#1\right\|}
\newcommand{\inner}[2]{\langle #1,\, #2 \rangle}
\newcommand{\tr}{\mathrm{tr}}

\title{\textbf{Подробные решения задач по линейной алгебре}\\[0.5em]
\large Скалярные произведения, нормы, операторы, собственные значения, квадратичные формы, ОДУ}
\date{}

\begin{document}
\maketitle
\tableofcontents

%=============================================================
\task{1}
%=============================================================
\subsection*{Условие}
В пространстве $\R[x]_2$ со скалярным произведением
$\displaystyle (p,q)=\int_{-1}^{1}p(x)q(x)\,dx$
построить ортонормированный базис из многочленов степеней $0$, $1$ и $2$.

\begin{theory}
\textbf{Процесс Грама--Шмидта.} Дана линейно независимая система $\{f_1,\ldots,f_n\}$.
Строим ортогональную систему:
\[
g_1=f_1,\qquad g_k = f_k - \sum_{j=1}^{k-1}\frac{(f_k,g_j)}{(g_j,g_j)}\,g_j.
\]
Нормируем: $e_k = g_k/\norm{g_k}$.
\end{theory}

\step{1} Берём исходный базис $f_1=1$, $f_2=x$, $f_3=x^2$.

\step{2} Первый вектор:
\[
g_1=1,\qquad \norm{g_1}^2=\int_{-1}^{1}1\,dx=2,\qquad e_1=\frac{1}{\sqrt{2}}.
\]

\step{3} Второй вектор. Вычислим проекцию:
\[
(f_2,g_1)=\int_{-1}^{1}x\cdot 1\,dx=0\quad\Rightarrow\quad g_2=x.
\]
\[
\norm{g_2}^2=\int_{-1}^{1}x^2\,dx=\frac{2}{3},\qquad e_2=\sqrt{\frac{3}{2}}\,x.
\]

\step{4} Третий вектор:
\[
(f_3,g_1)=\int_{-1}^{1}x^2\,dx=\frac{2}{3},\quad
(f_3,g_2)=\int_{-1}^{1}x^3\,dx=0.
\]
\[
g_3=x^2-\frac{(f_3,g_1)}{(g_1,g_1)}\,g_1=x^2-\frac{2/3}{2}\cdot 1=x^2-\frac{1}{3}.
\]
\[
\norm{g_3}^2=\int_{-1}^{1}\!\left(x^2-\tfrac{1}{3}\right)^2dx
=\int_{-1}^{1}\!\left(x^4-\tfrac{2}{3}x^2+\tfrac{1}{9}\right)dx
=\frac{2}{5}-\frac{4}{9}+\frac{2}{9}=\frac{8}{45}.
\]
\[
e_3=\frac{g_3}{\norm{g_3}}=\sqrt{\frac{45}{8}}\left(x^2-\frac{1}{3}\right)=\frac{3\sqrt{5}}{2\sqrt{2}}\left(x^2-\frac{1}{3}\right).
\]

\answer
\[
e_1=\frac{1}{\sqrt{2}},\qquad e_2=\sqrt{\frac{3}{2}}\,x,\qquad
e_3=\frac{3\sqrt{5}}{2\sqrt{2}}\!\left(x^2-\frac{1}{3}\right).
\]
Это многочлены Лежандра (с точностью до нормировки).

%=============================================================
\task{2}
%=============================================================
\subsection*{Условие}
В пространстве $C[0,1]$ рассматривается подмножество $V=\{f\in C[0,1]: f(0)=0\}$.
\begin{enumerate}
\item Доказать, что $\langle f,g\rangle=\int_0^1 f'(x)g'(x)\,dx$ — скалярное произведение на $V$.
\item Проверить неравенство Коши--Буняковского--Шварца для $f(x)=x$, $g(x)=x^2$.
\end{enumerate}

\begin{theory}
Для скалярного произведения требуется: \textbf{(i)} билинейность, \textbf{(ii)} симметричность, \textbf{(iii)} положительная определённость. Ключевое: $\langle f,f\rangle=0 \Rightarrow f=0$ — нетривиально и использует условие $f(0)=0$.
\end{theory}

\subsubsection*{Часть 1. Аксиомы скалярного произведения}

\textbf{Симметричность} и \textbf{билинейность} очевидны из свойств интеграла.

\textbf{Положительная определённость.}
\[
\langle f,f\rangle=\int_0^1[f'(x)]^2\,dx\geq 0.
\]
Если $\langle f,f\rangle=0$, то $f'(x)=0$ на $[0,1]$, значит $f=\mathrm{const}$. Из $f(0)=0$ получаем $f\equiv 0$. Аксиома выполнена.

\subsubsection*{Часть 2. Неравенство КБШ}

Для $f(x)=x$ и $g(x)=x^2$ (оба удовлетворяют $f(0)=g(0)=0$):
\[
\langle f,g\rangle=\int_0^1 1\cdot 2x\,dx=1,\quad
\langle f,f\rangle=\int_0^1 1\,dx=1,\quad
\langle g,g\rangle=\int_0^1 4x^2\,dx=\frac{4}{3}.
\]
Неравенство КБШ: $|\langle f,g\rangle|^2\leq\langle f,f\rangle\cdot\langle g,g\rangle$:
\[
1^2\leq 1\cdot\frac{4}{3}\quad\Leftrightarrow\quad 1\leq\frac{4}{3}.\quad\checkmark
\]

\answer Операция задаёт скалярное произведение на $V$; неравенство КБШ выполнено (равенство не достигается, так как $f$ и $g$ не пропорциональны).

%=============================================================
\task{3}
%=============================================================
\subsection*{Условие}
Пусть $V$ --- унитарное пространство. Доказать тождество поляризации для комплексного случая:
\[
\langle x,y\rangle=\frac{1}{4}\sum_{k=0}^{3}i^k\|x+i^ky\|^2.
\]

\begin{theory}
В унитарном пространстве $\langle x,y\rangle$ является полуторалинейным (линейным по первому аргументу и сопряжённо-линейным по второму), и $\langle y,x\rangle=\overline{\langle x,y\rangle}$.
\end{theory}

\step{1} Раскроем $\|x+i^k y\|^2 = \langle x+i^k y,\, x+i^k y\rangle$:
\[
\|x+i^k y\|^2 = \|x\|^2 + i^k\langle y,x\rangle + \overline{i^k}\langle x,y\rangle + \|y\|^2
= \|x\|^2+\|y\|^2 + \bar{i}^k\langle x,y\rangle + i^k\overline{\langle x,y\rangle}.
\]
Обозначим $w=\langle x,y\rangle$. Тогда:
\[
\|x+i^k y\|^2 = \|x\|^2+\|y\|^2 + \bar{i}^k w + i^k\bar w.
\]

\step{2} Умножаем на $i^k$ и суммируем:
\[
\sum_{k=0}^{3}i^k\|x+i^k y\|^2
=(\|x\|^2+\|y\|^2)\underbrace{\sum_{k=0}^3 i^k}_{=0}
+w\sum_{k=0}^3 i^k\bar{i}^k + \bar w\sum_{k=0}^3 i^{2k}.
\]

\step{3} Вычислим суммы ($i^k\bar i^k=|i^k|^2=1$):
\[
\sum_{k=0}^3 i^k\bar i^k = 4,\qquad
\sum_{k=0}^3 i^{2k}=\sum_{k=0}^3(-1)^k=0.
\]

\step{4} Итого:
\[
\sum_{k=0}^3 i^k\|x+i^k y\|^2 = 4w = 4\langle x,y\rangle.
\]
Делим на $4$ --- тождество доказано. $\blacksquare$

%=============================================================
\task{4}
%=============================================================
\subsection*{Условие}
Доказать, что для любых $a_1,\ldots,a_n>0$:
\[
(a_1+\cdots+a_n)\!\left(\frac{1}{a_1}+\cdots+\frac{1}{a_n}\right)\geq n^2.
\]
Указать условие равенства.

\begin{theory}
Неравенство КБШ в $\R^n$: $|\langle u,v\rangle|^2\leq\langle u,u\rangle\cdot\langle v,v\rangle$. Со стандартным скалярным произведением и векторами $u=(\sqrt{a_1},\ldots,\sqrt{a_n})$, $v=(1/\sqrt{a_1},\ldots,1/\sqrt{a_n})$.
\end{theory}

\step{1} Положим $u_i=\sqrt{a_i}$, $v_i=1/\sqrt{a_i}$.

\step{2} По КБШ:
\[
\left(\sum_{i=1}^n u_iv_i\right)^2\leq\left(\sum_{i=1}^n u_i^2\right)\!\left(\sum_{i=1}^n v_i^2\right).
\]

\step{3} Подставляем:
\[
\left(\sum_{i=1}^n 1\right)^2 = n^2\leq \left(\sum_{i=1}^n a_i\right)\!\left(\sum_{i=1}^n\frac{1}{a_i}\right).\quad\checkmark
\]

\step{4} \textbf{Условие равенства:} в КБШ равенство достигается тогда и только тогда, когда $u$ и $v$ коллинеарны, то есть $\sqrt{a_i}=c/\sqrt{a_i}$ для всех $i$, что означает $a_i=\mathrm{const}$, то есть $a_1=a_2=\cdots=a_n$.

\answer Неравенство доказано. Равенство при $a_1=a_2=\cdots=a_n$.

%=============================================================
\task{5}
%=============================================================
\subsection*{Условие}
Пусть $\langle f,g\rangle=\int_0^1 f(x)g(x)\,dx$ на $C[0,1]$.
\begin{enumerate}
\item Доказать: $\displaystyle\left(\int_0^1 f(x)\,dx\right)^2\leq\int_0^1 f^2(x)\,dx$.
\item Для каких функций достигается равенство?
\end{enumerate}

\begin{theory}
КБШ: $|\langle f,g\rangle|^2\leq\langle f,f\rangle\cdot\langle g,g\rangle$. Берём $g\equiv 1$.
\end{theory}

\step{1} Применим КБШ с $g\equiv 1$:
\[
\left|\int_0^1 f(x)\cdot 1\,dx\right|^2\leq\int_0^1 f^2(x)\,dx\cdot\int_0^1 1\,dx
=\int_0^1 f^2(x)\,dx.\quad\checkmark
\]

\step{2} \textbf{Равенство} в КБШ достигается тогда и только тогда, когда $f$ и $g=1$ коллинеарны в смысле данного скалярного произведения, то есть $f(x)=c\cdot 1$ --- $f$ является константой на $[0,1]$.

\answer Неравенство доказано. Равенство достигается для постоянных функций $f(x)\equiv c$.

%=============================================================
\task{6}
%=============================================================
\subsection*{Условие}
В $\R^2$ задана норма $\|x\|_p=(|x_1|^p+|x_2|^p)^{1/p}$. Доказать аналитически, что при $p=4$ эта норма не порождается никаким скалярным произведением.

\begin{theory}
\textbf{Критерий Иордана--фон Неймана:} норма $\|\cdot\|$ порождена скалярным произведением тогда и только тогда, когда выполняется тождество параллелограмма:
\[
\|u+v\|^2+\|u-v\|^2=2(\|u\|^2+\|v\|^2)\quad\text{для всех }u,v.
\]
\end{theory}

\step{1} Возьмём $u=(1,0)$, $v=(0,1)$.
\[
\|u\|_4^4=1,\quad\|v\|_4^4=1\quad\Rightarrow\quad\|u\|_4=\|v\|_4=1.
\]

\step{2}
\[
u+v=(1,1),\quad\|u+v\|_4=(1+1)^{1/4}=2^{1/4}.
\]
\[
u-v=(1,-1),\quad\|u-v\|_4=(1+1)^{1/4}=2^{1/4}.
\]

\step{3} Левая часть тождества параллелограмма:
\[
\|u+v\|_4^2+\|u-v\|_4^2=2\cdot(2^{1/4})^2=2\cdot 2^{1/2}=2\sqrt{2}.
\]

\step{4} Правая часть:
\[
2(\|u\|_4^2+\|v\|_4^2)=2(1+1)=4.
\]

\step{5} Поскольку $2\sqrt{2}\approx 2{,}83\neq 4$, тождество параллелограмма нарушается.

\answer По критерию Иордана--фон Неймана норма $\|\cdot\|_4$ на $\R^2$ не может быть порождена никаким скалярным произведением.

%=============================================================
\task{7}
%=============================================================
\subsection*{Условие}
В $C[0,1]$ задана норма $\|f\|_1=\int_0^1|f(x)|\,dx$.
\begin{enumerate}
\item Доказать три аксиомы нормы.
\item Можно ли порождать её скалярным произведением?
\end{enumerate}

\begin{theory}
Аксиомы нормы: \textbf{(N1)} $\|f\|\geq0$, $\|f\|=0\Leftrightarrow f=0$; \textbf{(N2)} $\|\alpha f\|=|\alpha|\|f\|$; \textbf{(N3)} $\|f+g\|\leq\|f\|+\|g\|$.
\end{theory}

\subsubsection*{Часть 1. Аксиомы нормы}

\textbf{(N1):} $|f(x)|\geq0$, поэтому $\|f\|_1\geq0$. Если $\|f\|_1=0$ и $f$ непрерывна, то $f\equiv0$.

\textbf{(N2):} $\|\alpha f\|_1=\int_0^1|\alpha f|dx=|\alpha|\int_0^1|f|\,dx=|\alpha|\|f\|_1$.

\textbf{(N3):} $\|f+g\|_1=\int_0^1|f+g|\,dx\leq\int_0^1(|f|+|g|)\,dx=\|f\|_1+\|g\|_1$.

\subsubsection*{Часть 2. Связь со скалярным произведением}

Проверим тождество параллелограмма для $f(x)=1$, $g(x)=x$:
\[
\|f\|_1=1,\quad\|g\|_1=\frac{1}{2}.
\]
\[
\|f+g\|_1=\int_0^1(1+x)\,dx=\frac{3}{2},\quad
\|f-g\|_1=\int_0^1(1-x)\,dx=\frac{1}{2}.
\]
Левая часть: $\left(\frac{3}{2}\right)^2+\left(\frac{1}{2}\right)^2=\frac{9}{4}+\frac{1}{4}=\frac{10}{4}=2{,}5$.

Правая часть: $2\left(1^2+\frac{1}{4}\right)=\frac{5}{2}=2{,}5$.

В данном примере тождество выполняется. Однако это не доказывает существование скалярного произведения --- нужна проверка на всех парах. Возьмём $f(x)=1$, $g(x)=\mathrm{sign}(x-1/2)$ (приближение ступенькой). В пространстве $C[0,1]$ норма $\|\cdot\|_1$ не порождает скалярное произведение, так как не является евклидовой нормой (единичный шар не является эллипсоидом).

\answer Все три аксиомы нормы выполнены. Данная норма \textbf{не} порождается скалярным произведением (единичный шар в $\|\cdot\|_1$ имеет «рёбра», что нарушает тождество параллелограмма для некоторых пар функций).

%=============================================================
\task{8}
%=============================================================
\subsection*{Условие}
В $\R^2$ задана норма $\|x\|_1=|x_1|+|x_2|$.
\begin{enumerate}
\item Изобразить и описать единичную сферу.
\item Доказать, что для $x=(1,0)$, $y=(0,1)$ тождество параллелограмма не выполняется.
\end{enumerate}

\step{1} \textbf{Единичная сфера:} $\{(x_1,x_2):|x_1|+|x_2|=1\}$ --- это квадрат (ромб) с вершинами $(\pm1,0)$, $(0,\pm1)$, стороны которого наклонены под углом $45°$ к осям.

\step{2} Проверим тождество параллелограмма для $x=(1,0)$, $y=(0,1)$:
\[
\|x\|_1=1,\quad\|y\|_1=1,\quad\|x+y\|_1=|1|+|1|=2,\quad\|x-y\|_1=1+1=2.
\]
Левая часть: $\|x+y\|_1^2+\|x-y\|_1^2=4+4=8$.

Правая часть: $2(\|x\|_1^2+\|y\|_1^2)=2(1+1)=4$.

$8\neq4$ --- тождество \textbf{нарушено}.

\answer Единичная сфера --- ромб (квадрат, повёрнутый на $45°$). Норма $\|\cdot\|_1$ не порождается скалярным произведением.

%=============================================================
\task{9}
%=============================================================
\subsection*{Условие}
В $\R^2$ задана норма $\|x\|_4=(x_1^4+x_2^4)^{1/4}$.
\begin{enumerate}
\item Вычислить обе части тождества параллелограмма для $x=(1,0)$, $y=(0,1)$.
\item Можно ли ввести скалярное произведение, порождающее эту норму?
\end{enumerate}

\step{1} $\|x\|_4=1$, $\|y\|_4=1$, $x+y=(1,1)$, $x-y=(1,-1)$.
\[
\|x+y\|_4=(1^4+1^4)^{1/4}=2^{1/4},\quad\|x-y\|_4=2^{1/4}.
\]

\step{2} Левая часть тождества параллелограмма:
\[
\|x+y\|_4^2+\|x-y\|_4^2=2\cdot(2^{1/4})^2=2\cdot\sqrt{2}=2\sqrt{2}\approx2{,}83.
\]

\step{3} Правая часть:
\[
2\bigl(\|x\|_4^2+\|y\|_4^2\bigr)=2(1+1)=4.
\]

\step{4} $2\sqrt{2}\neq4$ --- тождество параллелограмма \textbf{нарушено}.

\answer По теореме Иордана--фон Неймана норма $\|\cdot\|_4$ \textbf{не может} быть порождена никаким скалярным произведением на $\R^2$.

%=============================================================
\task{10}
%=============================================================
\subsection*{Условие}
$f(x)=1-x$, $g(x)=x$ на $[0,1]$.
\begin{enumerate}
\item Вычислить $\|f\|_\infty$, $\|g\|_\infty$, $\|f+g\|_\infty$, $\|f-g\|_\infty$.
\item Проверить критерий Иордана--фон Неймана.
\end{enumerate}

\step{1} Вычислим равномерные нормы:
\[
\|f\|_\infty=\max_{x\in[0,1]}|1-x|=1,\quad
\|g\|_\infty=\max_{x\in[0,1]}|x|=1.
\]
\[
f(x)+g(x)=1,\quad\|f+g\|_\infty=1.
\]
\[
f(x)-g(x)=1-2x,\quad\|f-g\|_\infty=\max_{x\in[0,1]}|1-2x|=1.
\]

\step{2} Левая часть тождества параллелограмма:
\[
\|f+g\|_\infty^2+\|f-g\|_\infty^2=1+1=2.
\]

\step{3} Правая часть:
\[
2\bigl(\|f\|_\infty^2+\|g\|_\infty^2\bigr)=2(1+1)=4.
\]

\step{4} $2\neq4$ --- тождество нарушено.

\answer Равномерная норма на $C[0,1]$ \textbf{не порождается} никаким скалярным произведением.

%=============================================================
\task{11}
%=============================================================
\subsection*{Условие}
В $\mathrm{Mat}_{2\times2}(\R)$ со скалярным произведением $(A,B)=\tr(A^TB)$ найти расстояние от матрицы $M=\begin{pmatrix}1&1\\1&1\end{pmatrix}$ до подпространства кососимметричных матриц.

\begin{theory}
Расстояние от вектора до подпространства равно норме его ортогональной составляющей. Для кососимметричных матриц $\mathrm{Skew}=\{A:A^T=-A\}$. Ортогональное дополнение $\mathrm{Skew}^\perp$ --- подпространство симметричных матриц $\mathrm{Sym}=\{A:A^T=A\}$. Разложение: $M=M_{\mathrm{sym}}+M_{\mathrm{skew}}$.
\end{theory}

\step{1} Разложим $M$ на симметричную и кососимметричную части:
\[
M_{\mathrm{sym}}=\frac{M+M^T}{2}=M=\begin{pmatrix}1&1\\1&1\end{pmatrix},\quad
M_{\mathrm{skew}}=\frac{M-M^T}{2}=\begin{pmatrix}0&0\\0&0\end{pmatrix}.
\]
(Матрица $M$ уже симметрична.)

\step{2} Ортогональная проекция $M$ на $\mathrm{Skew}$ равна $M_{\mathrm{skew}}=0$.

\step{3} Расстояние:
\[
d(M,\mathrm{Skew})=\|M-M_{\mathrm{skew}}\|=\|M_{\mathrm{sym}}\|=\sqrt{\tr(M^TM)}.
\]
\[
M^TM=M^2=\begin{pmatrix}2&2\\2&2\end{pmatrix},\quad\tr(M^2)=4.
\]
\[
d=\sqrt{4}=2.
\]

\answer $d(M,\mathrm{Skew})=2$.

%=============================================================
\task{12}
%=============================================================
\subsection*{Условие}
Найти многочлен $p(x)=ax+b$, минимизирующий $\displaystyle\int_0^1(x^2-p(x))^2\,dx$.

\begin{theory}
Задача эквивалентна нахождению ортогональной проекции функции $f(x)=x^2$ на подпространство $L=\mathrm{span}\{1,x\}$ в пространстве $C[0,1]$ со скалярным произведением $\langle f,g\rangle=\int_0^1 fg\,dx$. Проекция находится из условий $\langle x^2-p,1\rangle=0$ и $\langle x^2-p,x\rangle=0$.
\end{theory}

\step{1} Условие ортогональности:
\[
\int_0^1(x^2-ax-b)\,dx=0\quad\Rightarrow\quad\frac{1}{3}-\frac{a}{2}-b=0.
\]
\[
\int_0^1 x(x^2-ax-b)\,dx=0\quad\Rightarrow\quad\frac{1}{4}-\frac{a}{3}-\frac{b}{2}=0.
\]

\step{2} Система уравнений:
\[
\begin{cases}\dfrac{a}{2}+b=\dfrac{1}{3},\\[6pt]\dfrac{a}{3}+\dfrac{b}{2}=\dfrac{1}{4}.\end{cases}
\]

\step{3} Из первого уравнения $b=\frac{1}{3}-\frac{a}{2}$. Подставляем:
\[
\frac{a}{3}+\frac{1}{2}\!\left(\frac{1}{3}-\frac{a}{2}\right)=\frac{1}{4}
\;\Rightarrow\;\frac{a}{3}+\frac{1}{6}-\frac{a}{4}=\frac{1}{4}
\;\Rightarrow\;\frac{a}{12}=\frac{1}{12}\;\Rightarrow\;a=1.
\]
\[
b=\frac{1}{3}-\frac{1}{2}=-\frac{1}{6}.
\]

\answer $p(x)=x-\dfrac{1}{6}$ является ортогональной проекцией $x^2$ на $\mathrm{span}\{1,x\}$ и минимизирует заданный интеграл.

%=============================================================
\task{13}
%=============================================================
\subsection*{Условие}
$V=C[0,1]$. Проверить, является ли скалярным произведением формула $\langle f,g\rangle=f(0)g(0)+f(1)g(1)$.

\begin{theory}
Для скалярного произведения необходима \textbf{положительная определённость}: $\langle f,f\rangle=0\Rightarrow f=0$.
\end{theory}

\step{1} Симметричность и билинейность очевидны.

\step{2} Проверим положительную определённость. Рассмотрим функцию $h(t)=t(1-t)\in C[0,1]$. Тогда $h(0)=0$, $h(1)=0$, значит:
\[
\langle h,h\rangle=h(0)^2+h(1)^2=0+0=0.
\]
Но $h\not\equiv0$ (например, $h(1/2)=1/4\neq0$).

\step{3} Аксиома положительной определённости \textbf{нарушена}.

\answer Формула не задаёт скалярного произведения. Нарушена аксиома положительной определённости. Пример ненулевой функции с $\langle h,h\rangle=0$: $h(t)=t(1-t)$.

%=============================================================
\task{14}
%=============================================================
\subsection*{Условие}
В $\mathrm{Mat}_{2\times2}$ с $\langle A,B\rangle=\tr(A^TB)$.
Даны $A=\begin{pmatrix}1&-1\\0&2\end{pmatrix}$, $B=\begin{pmatrix}2&2\\3&0\end{pmatrix}$.
\begin{enumerate}
\item Доказать ортогональность $A$ и $B$.
\item Вычислить $\|A\|$, $\|B\|$, $\|A+B\|$; проверить теорему Пифагора.
\end{enumerate}

\step{1} \textbf{Ортогональность.} Вычислим $\tr(A^TB)$:
\[
A^T=\begin{pmatrix}1&0\\-1&2\end{pmatrix},\quad
A^TB=\begin{pmatrix}1&0\\-1&2\end{pmatrix}\begin{pmatrix}2&2\\3&0\end{pmatrix}
=\begin{pmatrix}2&2\\4&-2\end{pmatrix}.
\]
\[
\tr(A^TB)=2+(-2)=0.\quad\checkmark
\]

\step{2} \textbf{Нормы.}
\[
\|A\|^2=\tr(A^TA)=1^2+0^2+(-1)^2+2^2=1+0+1+4=6,\quad\|A\|=\sqrt{6}.
\]
\[
\|B\|^2=\tr(B^TB)=2^2+3^2+2^2+0^2=4+9+4+0=17,\quad\|B\|=\sqrt{17}.
\]
\[
A+B=\begin{pmatrix}3&1\\3&2\end{pmatrix},\quad
\|A+B\|^2=9+9+1+4=23,\quad\|A+B\|=\sqrt{23}.
\]

\step{3} \textbf{Теорема Пифагора:} $\|A+B\|^2=\|A\|^2+\|B\|^2$ (т.к. $\langle A,B\rangle=0$):
\[
6+17=23.\quad\checkmark
\]

\answer $A\perp B$; $\|A\|=\sqrt{6}$, $\|B\|=\sqrt{17}$, $\|A+B\|=\sqrt{23}=\sqrt{6+17}$ --- теорема Пифагора подтверждена.

%=============================================================
\task{15}
%=============================================================
\subsection*{Условие}
В $C[-\pi,\pi]$ с $\langle f,g\rangle=\int_{-\pi}^{\pi}f(x)g(x)\,dx$. $f(x)=\sin x$, $g(x)=\cos x$.
\begin{enumerate}
\item Доказать ортогональность.
\item Связать с линейной независимостью; вычислить вронскиан в $x=0$.
\end{enumerate}

\step{1} \textbf{Ортогональность:}
\[
\langle\sin x,\cos x\rangle=\int_{-\pi}^{\pi}\sin x\cos x\,dx
=\frac{1}{2}\int_{-\pi}^{\pi}\sin 2x\,dx
=\frac{1}{2}\left[-\frac{\cos 2x}{2}\right]_{-\pi}^{\pi}=0.\quad\checkmark
\]

\step{2} \textbf{Линейная независимость} следует из ортогональности: ненулевые ортогональные векторы всегда линейно независимы.

\step{3} \textbf{Вронскиан:}
\[
W(x)=\begin{vmatrix}\sin x&\cos x\\\cos x&-\sin x\end{vmatrix}=-\sin^2x-\cos^2x=-1.
\]
\[
W(0)=-1\neq0.
\]
Невырожденный вронскиан подтверждает линейную независимость.

\answer $\sin x\perp\cos x$ на $[-\pi,\pi]$; функции линейно независимы (подтверждено вронскианом $W(0)=-1\neq0$).

%=============================================================
\task{16}
%=============================================================
\subsection*{Условие}
В $P_1[-1,1]$ с $\langle f,g\rangle=\int_{-1}^1 fg\,dx$ задана система $e_1=\frac{1}{\sqrt{2}}$, $e_2=\sqrt{\frac{3}{2}}x$.
\begin{enumerate}
\item Доказать, что система ортонормированная и образует ОНБ.
\item Разложить $f(x)=2x+1$ по этому базису.
\end{enumerate}

\step{1} \textbf{Ортонормированность.}
\[
\langle e_1,e_1\rangle=\frac{1}{2}\int_{-1}^1 1\,dx=1,\quad
\langle e_2,e_2\rangle=\frac{3}{2}\int_{-1}^1 x^2\,dx=\frac{3}{2}\cdot\frac{2}{3}=1.
\]
\[
\langle e_1,e_2\rangle=\frac{\sqrt{3}}{2}\int_{-1}^1 x\,dx=0.
\]
Система ортонормирована. Так как $\dim P_1=2$ и система содержит 2 вектора --- это ОНБ.

\step{2} \textbf{Коэффициенты Фурье} $f(x)=2x+1$:
\[
c_1=\langle f,e_1\rangle=\frac{1}{\sqrt{2}}\int_{-1}^1(2x+1)\,dx
=\frac{1}{\sqrt{2}}\left[x^2+x\right]_{-1}^1
=\frac{1}{\sqrt{2}}[(1+1)-(1-1)]=\frac{2}{\sqrt{2}}=\sqrt{2}.
\]
\[
c_2=\langle f,e_2\rangle=\sqrt{\frac{3}{2}}\int_{-1}^1 x(2x+1)\,dx
=\sqrt{\frac{3}{2}}\int_{-1}^1(2x^2+x)\,dx
=\sqrt{\frac{3}{2}}\cdot\frac{4}{3}=\frac{4}{3}\sqrt{\frac{3}{2}}=\frac{2\sqrt{2}}{\sqrt{3}}.
\]

\step{3} Разложение: $f=c_1 e_1+c_2 e_2=\sqrt{2}\cdot\frac{1}{\sqrt{2}}+\frac{2\sqrt{2}}{\sqrt{3}}\cdot\sqrt{\frac{3}{2}}\,x=1+2x$.

\answer $f(x)=2x+1=\sqrt{2}\,e_1+\dfrac{2\sqrt{2}}{\sqrt{3}}\,e_2$.

%=============================================================
\task{17}
%=============================================================
\subsection*{Условие}
В $\R^3$ с ОНБ $\{e_1,e_2,e_3\}$: $x=(1,-2,2)$, $y=(3,0,-4)$.
\begin{enumerate}
\item Вычислить $\|x\|$, $\|y\|$, $\langle x,y\rangle$.
\item Проверить равенство Парсеваля для $z=x+y$ двумя способами.
\end{enumerate}

\step{1} В ОНБ: $\|x\|^2=1+4+4=9$, $\|x\|=3$; $\|y\|^2=9+0+16=25$, $\|y\|=5$.
\[
\langle x,y\rangle=1\cdot3+(-2)\cdot0+2\cdot(-4)=3-8=-5.
\]

\step{2} $z=x+y=(4,-2,-2)$.

\textbf{Способ 1 (прямой):} $\|z\|^2=16+4+4=24$.

\textbf{Способ 2 (через $x$ и $y$):}
\[
\|z\|^2=\|x+y\|^2=\|x\|^2+2\langle x,y\rangle+\|y\|^2=9+2(-5)+25=24.\quad\checkmark
\]

\answer $\|x\|=3$, $\|y\|=5$, $\langle x,y\rangle=-5$; $\|z\|^2=24$ (оба способа совпадают).

%=============================================================
\task{18}
%=============================================================
\subsection*{Условие}
Проверить, является ли нормой в $\mathrm{Mat}_{n\times n}(\R)$ величина $\|A\|_*=\max_{i,j}|a_{ij}|$. Доказать неравенство треугольника. Является ли норма евклидовой при $n=2$?

\step{1} \textbf{(N1)} $\|A\|_*\geq0$; $\|A\|_*=0\Leftrightarrow$ все $a_{ij}=0\Leftrightarrow A=0$.

\step{2} \textbf{(N2)} $\|\alpha A\|_*=\max_{i,j}|\alpha a_{ij}|=|\alpha|\max_{i,j}|a_{ij}|=|\alpha|\|A\|_*$.

\step{3} \textbf{(N3) Неравенство треугольника:}
\[
\|A+B\|_*=\max_{i,j}|a_{ij}+b_{ij}|\leq\max_{i,j}(|a_{ij}|+|b_{ij}|)\leq\max_{i,j}|a_{ij}|+\max_{i,j}|b_{ij}|=\|A\|_*+\|B\|_*.
\]

\step{4} \textbf{Евклидовость при $n=2$.} Проверим тождество параллелограмма для $A=\begin{pmatrix}1&0\\0&0\end{pmatrix}$, $B=\begin{pmatrix}0&1\\0&0\end{pmatrix}$:
\[
\|A\|_*=1,\;\|B\|_*=1,\;\|A+B\|_*=1,\;\|A-B\|_*=1.
\]
Левая часть: $1+1=2$. Правая: $2(1+1)=4$. $2\neq4$ --- не евклидова.

\answer $\|\cdot\|_*$ является нормой, но при $n=2$ не евклидовой.

%=============================================================
\task{19}
%=============================================================
\subsection*{Условие}
При каких $\lambda\in\R$ выражение $\langle x,y\rangle_\lambda=x_1y_1+\lambda x_1y_2+\lambda x_2y_1+4x_2y_2$ является скалярным произведением?

\begin{theory}
\textbf{Критерий Сильвестра:} симметричная матрица положительно определена тогда и только тогда, когда все угловые миноры положительны.
\end{theory}

\step{1} Матрица формы:
\[
G=\begin{pmatrix}1&\lambda\\\lambda&4\end{pmatrix}.
\]

\step{2} Условия Сильвестра:
\[
\Delta_1=1>0,\quad\Delta_2=\det G=4-\lambda^2>0.
\]

\step{3} Из $4-\lambda^2>0$ получаем $\lambda^2<4$, то есть $-2<\lambda<2$.

\answer Скалярное произведение определено при $\lambda\in(-2,2)$.

%=============================================================
\task{20}
%=============================================================
\subsection*{Условие}
Функция $p:\R^2\to\R$: $p(x)=(\sqrt{|x_1|}+\sqrt{|x_2|})^2$. Проверить аксиомы нормы.

\step{1} \textbf{Положительная определённость.} $p(x)\geq0$; $p(x)=0\Leftrightarrow\sqrt{|x_1|}+\sqrt{|x_2|}=0\Leftrightarrow x_1=x_2=0$. Выполнена.

\step{2} \textbf{Однородность.}
\[
p(\alpha x)=(\sqrt{|\alpha x_1|}+\sqrt{|\alpha x_2|})^2=(|\alpha|^{1/2}\sqrt{|x_1|}+|\alpha|^{1/2}\sqrt{|x_2|})^2
=|\alpha|(\sqrt{|x_1|}+\sqrt{|x_2|})^2=|\alpha|p(x).\quad\checkmark
\]

\step{3} \textbf{Неравенство треугольника.} Проверим для $x=(1,0)$, $y=(0,1)$:
\[
p(x)=1,\quad p(y)=1,\quad
p(x+y)=(\sqrt{1}+\sqrt{1})^2=4.
\]
\[
p(x)+p(y)=1+1=2<4=p(x+y).
\]
Неравенство треугольника \textbf{нарушено}.

\answer Положительная определённость и однородность выполнены, но неравенство треугольника \textbf{нарушено} --- $p$ не является нормой.

%=============================================================
\task{21}
%=============================================================
\subsection*{Условие}
В пространстве непрерывных комплекснозначных функций $C([0, 1], \C)$ задано выражение:
\[
\langle f, g\rangle = \int_0^1 f(t)\overline{g(t)}w(t)\,dt.
\]
Доказать, что если весовая функция $w(t)$ непрерывна и строго положительна на $[0, 1]$, то данное выражение удовлетворяет всем аксиомам унитарного скалярного произведения.

\begin{theory}
Для того чтобы выражение в комплексном пространстве $V$ задавало унитарное (комплексное) скалярное произведение, должны выполняться четыре аксиомы:
\begin{enumerate}
    \item \textbf{Эрмитова симметрия:} $\langle f, g\rangle = \overline{\langle g, f\rangle}$ для любых $f, g \in V$;
    \item \textbf{Линейность по первому аргументу:} $\langle \alpha f + \beta g, h\rangle = \alpha \langle f, h\rangle + \beta \langle g, h\rangle$ для любых $\alpha, \beta \in \C$ и $f, g, h \in V$;
    \item \textbf{Положительность:} $\langle f, f\rangle \geq 0$ для любого $f \in V$;
    \item \textbf{Невырожденность:} $\langle f, f\rangle = 0 \iff f \equiv 0$.
\end{enumerate}
\end{theory}

\step{1} \textbf{Эрмитова симметрия.} Проверим равенство $\langle f, g\rangle = \overline{\langle g, f\rangle}$:
\[
\overline{\langle g, f\rangle} = \overline{\int_0^1 g(t)\overline{f(t)}w(t)\,dt} = \int_0^1 \overline{g(t)\overline{f(t)}w(t)}\,dt.
\]
Поскольку весовая функция вещественна ($w(t) > 0$), её сопряжение равно ей самой ($\overline{w(t)} = w(t)$). Используя свойства комплексного сопряжения $\overline{z_1 z_2} = \overline{z_1} \cdot \overline{z_2}$ и $\overline{\overline{z}} = z$, получаем:
\[
\overline{\langle g, f\rangle} = \int_0^1 \overline{g(t)}f(t)w(t)\,dt = \int_0^1 f(t)\overline{g(t)}w(t)\,dt = \langle f, g\rangle. \quad\checkmark
\]

\step{2} \textbf{Линейность по первому аргументу.} Из линейности интеграла напрямую следует:
\[
\langle \alpha f + \beta g, h\rangle = \int_0^1 (\alpha f(t) + \beta g(t))\overline{h(t)}w(t)\,dt 
\]
\[
= \alpha \int_0^1 f(t)\overline{h(t)}w(t)\,dt + \beta \int_0^1 g(t)\overline{h(t)}w(t)\,dt = \alpha \langle f, h\rangle + \beta \langle g, h\rangle. \quad\checkmark
\]

\step{3} \textbf{Положительность.} Запишем скалярное произведение функции самой на себя:
\[
\langle f, f\rangle = \int_0^1 f(t)\overline{f(t)}w(t)\,dt = \int_0^1 |f(t)|^2 w(t)\,dt.
\]
Поскольку $|f(t)|^2 \geq 0$ и по условию $w(t) > 0$ на $[0, 1]$, подынтегральное выражение везде неотрицательно. Следовательно, интеграл от неотрицательной функции также неотрицателен: $\langle f, f\rangle \geq 0$. \quad\checkmark

\step{4} \textbf{Невырожденность.} Пусть $\langle f, f\rangle = 0$, то есть:
\[
\int_0^1 |f(t)|^2 w(t)\,dt = 0.
\]
Поскольку функции $f(t)$ и $w(t)$ непрерывны, функция $F(t) = |f(t)|^2 w(t)$ непрерывна и неотрицательна на $[0, 1]$. Из математического анализа известно, что интеграл от непрерывной неотрицательной функции равен нулю тогда и только тогда, когда сама функция тождественно равна нулю:
\[
|f(t)|^2 w(t) = 0 \quad \text{для всех } t \in [0, 1].
\]
Так как $w(t) > 0$ на всём отрезке, то $|f(t)|^2 = 0$, откуда $f(t) \equiv 0$ на $[0, 1]$. Обратное утверждение (если $f \equiv 0$, то $\langle f, f\rangle = 0$) очевидно. \quad\checkmark

\answer
Все аксиомы комплексного (унитарного) скалярного произведения выполнены. Вещественность и строгая положительность $w(t)$ гарантируют корректность аксиомы невырожденности.

%=============================================================
\task{22}
%=============================================================
\subsection*{Условие}
В евклидовом пространстве $\R^4$ со стандартным скалярным произведением дополнить систему векторов $a_1 = (1, 1, 1, 1)$, $a_2 = (1, 1, -1, -1)$ до ортогонального базиса и пронормировать его.

\begin{theory}
Два вектора $x, y \in \R^4$ ортогональны, если их стандартное скалярное произведение равно нулю: $\langle x, y\rangle = \sum x_i y_i = 0$. Чтобы дополнить ортогональную систему $\{a_1, a_2\}$ до ортогонального базиса в $\R^4$, нужно найти ненулевые векторы $a_3, a_4$, удовлетворяющие системе уравнений:
\[
\begin{cases}
\langle a_3, a_1\rangle = 0, \\
\langle a_3, a_2\rangle = 0
\end{cases}
\]
а затем выбрать их взаимно ортогональными ($\langle a_3, a_4\rangle = 0$). После построения ортогонального базиса каждый вектор нормируется: $e_i = a_i / \|a_i\|$.
\end{theory}

\step{1} Проверим ортогональность исходных векторов:
\[
\langle a_1, a_2\rangle = 1\cdot 1 + 1\cdot 1 + 1\cdot(-1) + 1\cdot(-1) = 1 + 1 - 1 - 1 = 0. \quad\checkmark
\]

\step{2} Найдём вектор $a_3 = (x_1, x_2, x_3, x_4)$, ортогональный $a_1$ и $a_2$:
\[
\begin{cases}
x_1 + x_2 + x_3 + x_4 = 0, \\
x_1 + x_2 - x_3 - x_4 = 0.
\end{cases}
\]
Складывая и вычитая эти уравнения, получаем эквивалентную систему:
\[
\begin{cases}
2x_1 + 2x_2 = 0 \iff x_2 = -x_1, \\
2x_3 + 2x_4 = 0 \iff x_4 = -x_3.
\end{cases}
\]
Выберем один простой ненулевой вектор, удовлетворяющий этим соотношениям. Например, положим $x_1 = 1$, тогда $x_2 = -1$. Положим $x_3 = 0$, тогда $x_4 = 0$. Получаем:
\[
a_3 = (1, -1, 0, 0).
\]

\step{3} Найдём вектор $a_4 = (y_1, y_2, y_3, y_4)$, ортогональный $a_1$, $a_2$ и $a_3$. Он должен удовлетворять условиям $y_2 = -y_1$, $y_4 = -y_3$, а также:
\[
\langle a_4, a_3\rangle = 0 \implies y_1 - y_2 = 0 \implies y_1 = y_2.
\]
Поскольку $y_2 = -y_1$ и $y_1 = y_2$, получаем $y_1 = y_2 = 0$.
Тогда из первого шага $y_4 = -y_3$. Положим $y_3 = 1$, тогда $y_4 = -1$. Получаем:
\[
a_4 = (0, 0, 1, -1).
\]
Легко проверить, что семейство $\{a_1, a_2, a_3, a_4\}$ взаимно ортогонально.

\step{4} Пронормируем полученные векторы. Вычислим их длины (нормы):
\[
\|a_1\| = \sqrt{1^2 + 1^2 + 1^2 + 1^2} = 2,
\]
\[
\|a_2\| = \sqrt{1^2 + 1^2 + (-1)^2 + (-1)^2} = 2,
\]
\[
\|a_3\| = \sqrt{1^2 + (-1)^2 + 0^2 + 0^2} = \sqrt{2},
\]
\[
\|a_4\| = \sqrt{0^2 + 0^2 + 1^2 + (-1)^2} = \sqrt{2}.
\]
Разделим каждый вектор на его норму:
\[
e_1 = \left(\frac{1}{2}, \frac{1}{2}, \frac{1}{2}, \frac{1}{2}\right), \qquad e_2 = \left(\frac{1}{2}, \frac{1}{2}, -\frac{1}{2}, -\frac{1}{2}\right),
\]
\[
e_3 = \left(\frac{1}{\sqrt{2}}, -\frac{1}{\sqrt{2}}, 0, 0\right), \qquad e_4 = \left(0, 0, \frac{1}{\sqrt{2}}, -\frac{1}{\sqrt{2}}\right).
\]

\answer
Ортонормированный базис:
\[
e_1 = \frac{1}{2}(1, 1, 1, 1), \quad e_2 = \frac{1}{2}(1, 1, -1, -1), \quad e_3 = \frac{1}{\sqrt{2}}(1, -1, 0, 0), \quad e_4 = \frac{1}{\sqrt{2}}(0, 0, 1, -1).
\]

%=============================================================
\task{23}
%=============================================================
\subsection*{Условие}
Доказать, что определитель Грама системы векторов $\{v_1, \ldots, v_k\}$ равен нулю тогда и только тогда, когда эти векторы линейно зависимы. Вычислить определитель Грама для векторов $v_1 = (1, i)$, $v_2 = (1, -i)$ в унитарном пространстве $\C^2$.

\begin{theory}
Матрица Грама системы векторов $\{v_1, \ldots, v_k\}$ определяется как $G(v_1, \ldots, v_k)_{ij} = \langle v_i, v_j\rangle$. Определитель матрицы Грама называется определителем Грама (грамианом) и обозначается $\Gamma(v_1, \ldots, v_k)$. Он всегда неотрицателен, а равенство нулю означает линейную зависимость системы.
\end{theory}

\subsubsection*{Часть 1. Доказательство критерия}

\step{1} \textbf{Необходимость ($\Rightarrow$).} Пусть $\Gamma(v_1, \ldots, v_k) = 0$. Это значит, что определитель матрицы Грама $G$ равен нулю, следовательно, её столбцы линейно зависимы. Существует ненулевой набор коэффициентов $c = (c_1, \ldots, c_k)^T \neq 0$, такой что $G c = 0$. Запишем это поэлементно:
\[
\sum_{j=1}^k \langle v_i, v_j\rangle c_j = 0 \quad \text{для всех } i = 1, \ldots, k.
\]
Используя линейность скалярного произведения по второму аргументу (или первому, в зависимости от конвенции сопряжения; для определённости будем считать линейность по первому аргументу, $\langle \sum c_j v_j, v_i \rangle = 0$):
\[
\left\langle v_i, \sum_{j=1}^k \overline{c_j} v_j \right\rangle = 0 \quad \text{для всех } i.
\]
Умножим каждое уравнение на $\overline{c_i}$ и просуммируем по $i$:
\[
\left\langle \sum_{i=1}^k c_i v_i, \sum_{j=1}^k c_j v_j \right\rangle = 0 \implies \left\| \sum_{i=1}^k c_i v_i \right\|^2 = 0.
\]
В силу положительной определённости нормы получаем $\sum_{i=1}^k c_i v_i = 0$. Так как не все $c_i$ равны нулю, система векторов $\{v_1, \ldots, v_k\}$ линейно зависима.

\step{2} \textbf{Достаточность ($\Leftarrow$).} Пусть векторы $\{v_1, \ldots, v_k\}$ линейно зависимы. Тогда существует нетривиальная линейная комбинация $\sum_{j=1}^k \alpha_j v_j = 0$, где не все $\alpha_j$ равны нулю.
Для любого вектора $v_i$ выполнено:
\[
\left\langle v_i, \sum_{j=1}^k \alpha_j v_j \right\rangle = \langle v_i, 0\rangle = 0 \implies \sum_{j=1}^k \overline{\alpha_j} \langle v_i, v_j\rangle = 0 \quad \text{для каждого } i = 1, \ldots, k.
\]
Это означает, что столбцы матрицы Грама линейно зависимы с коэффициентами $\overline{\alpha_j}$. Следовательно, матрица Грама вырождена, и её определитель $\Gamma(v_1, \ldots, v_k) = \det G = 0$. \quad$\blacksquare$

\subsubsection*{Часть 2. Вычисление грамиана в $C^2$}

\step{3} Будем использовать стандартное комплексное скалярное произведение в $C^2$: $\langle x, y\rangle = x_1\overline{y_1} + x_2\overline{y_2}$.
Вычислим элементы матрицы Грама для векторов $v_1 = (1, i)$ и $v_2 = (1, -i)$:
\[
\langle v_1, v_1\rangle = 1\cdot \overline{1} + i\cdot \overline{i} = 1 + i(-i) = 1 + 1 = 2,
\]
\[
\langle v_2, v_2\rangle = 1\cdot \overline{1} + (-i)\cdot \overline{-i} = 1 + (-i)(i) = 1 + 1 = 2,
\]
\[
\langle v_1, v_2\rangle = 1\cdot \overline{1} + i\cdot \overline{-i} = 1 + i(i) = 1 - 1 = 0,
\]
\[
\langle v_2, v_1\rangle = 1\cdot \overline{1} + (-i)\cdot \overline{i} = 1 + (-i)(-i) = 1 - 1 = 0.
\]
Матрица Грама имеет вид:
\[
G = \begin{pmatrix} 2 & 0 \\ 0 & 2 \end{pmatrix}.
\]

\step{4} Вычисляем определитель:
\[
\Gamma(v_1, v_2) = \det G = 2\cdot 2 - 0\cdot 0 = 4.
\]

\answer
Критерий равенства грамиана нулю доказан. Для векторов $v_1 = (1, i)$, $v_2 = (1, -i)$ определитель Грама равен $\Gamma(v_1, v_2) = 4$.

%=============================================================
\task{24}
%=============================================================
\subsection*{Условие}
В евклидовом пространстве $C[0, 1]$ со стандартным скалярным произведением $\langle f, g\rangle = \int_0^1 f(x)g(x)\,dx$ рассматривается система из трех функций: $f_1(x) = 1$, $f_2(x) = x$, $f_3(x) = x^2$.
\begin{enumerate}
    \item Составить матрицу Грама для этой системы функций, вычислив все её элементы.
    \item Вычислить определитель Грама полученной матрицы. Сделать заключение о линейной зависимости или независимости данной системы функций.
\end{enumerate}

\begin{theory}
Элементы матрицы Грама для системы функций $\{f_1, f_2, f_3\}$ вычисляются как интегралы $G_{ij} = \langle f_i, f_j\rangle = \int_0^1 f_i(x)f_j(x)\,dx$. Линейная независимость гарантируется неравенством $\det G > 0$.
\end{theory}

\step{1} Вычислим элементы матрицы Грама $G$:
\[
G_{11} = \langle 1, 1\rangle = \int_0^1 1\,dx = 1,
\]
\[
G_{12} = G_{21} = \langle 1, x\rangle = \int_0^1 x\,dx = \frac{1}{2},
\]
\[
G_{13} = G_{31} = G_{22} = \langle 1, x^2\rangle = \langle x, x\rangle = \int_0^1 x^2\,dx = \frac{1}{3},
\]
\[
G_{23} = G_{32} = \langle x, x^2\rangle = \int_0^1 x^3\,dx = \frac{1}{4},
\]
\[
G_{33} = \langle x^2, x^2\rangle = \int_0^1 x^4\,dx = \frac{1}{5}.
\]
Матрица Грама (известная как матрица Гильберта размера $3\times3$):
\[
G = \begin{pmatrix} 1 & 1/2 & 1/3 \\ 1/2 & 1/3 & 1/4 \\ 1/3 & 1/4 & 1/5 \end{pmatrix}.
\]

\step{2} Вычислим определитель матрицы $G$ разложением по первой строке:
\[
\det G = 1 \cdot \left( \frac{1}{3}\cdot\frac{1}{5} - \frac{1}{4}\cdot\frac{1}{4} \right) - \frac{1}{2} \cdot \left( \frac{1}{2}\cdot\frac{1}{5} - \frac{1}{4}\cdot\frac{1}{3} \right) + \frac{1}{3} \cdot \left( \frac{1}{2}\cdot\frac{1}{4} - \frac{1}{3}\cdot\frac{1}{3} \right)
\]
\[
= 1 \cdot \left( \frac{1}{15} - \frac{1}{16} \right) - \frac{1}{2} \cdot \left( \frac{1}{10} - \frac{1}{12} \right) + \frac{1}{3} \cdot \left( \frac{1}{8} - \frac{1}{9} \right)
\]
\[
= \frac{1}{240} - \frac{1}{2} \cdot \frac{1}{60} + \frac{1}{3} \cdot \frac{1}{72} = \frac{1}{240} - \frac{1}{120} + \frac{1}{216}.
\]
Приведём дроби к общему знаменателю. Наименьшее общее кратное чисел $240, 120, 216$ равно $2160$:
\[
\det G = \frac{9 - 18 + 10}{2160} = \frac{1}{2160}.
\]

\step{3} Так как определитель Грама $\Gamma(1, x, x^2) = \frac{1}{2160} \neq 0$ (и строго положителен), система функций $\{1, x, x^2\}$ является линейно независимой.

\answer
1. Матрица Грама: $G = \begin{pmatrix} 1 & 1/2 & 1/3 \\ 1/2 & 1/3 & 1/4 \\ 1/3 & 1/4 & 1/5 \end{pmatrix}$.\\
2. Определитель Грама: $\det G = \dfrac{1}{2160}$. Система функций линейно независима.

%=============================================================
\task{25}
%=============================================================
\subsection*{Условие}
В вещественном пространстве $\R^4$ со стандартным скалярным произведением заданы три вектора: $x_1 = (1, 1, 0, 0)$, $x_2 = (0, 1, 1, 0)$ и $y = (1, 1, 1, 1)$.
\begin{enumerate}
    \item Вычислить определитель Грама системы векторов $\{x_1, x_2\}$, определяющий площадь основания параллелепипеда. Затем вычислить определитель Грама системы трех векторов $\{x_1, x_2, y\}$, задающий квадрат трёхмерного объёма.
    \item Используя формулу геометрического смысла грамиана, найти высоту параллелепипеда, опущенную из вершины $y$ на плоскость линейной оболочки $\mathrm{Lin}(x_1, x_2)$.
\end{enumerate}

\begin{theory}
Геометрический смысл грамиана заключается в том, что определитель Грама $\Gamma(v_1, \ldots, v_k)$ равен квадрату $k$-мерного объёма параллелепипеда, построенного на этих векторах: $V_k^2 = \Gamma(v_1, \ldots, v_k)$. Высота параллелепипеда $h$, опущенная из вершины $y$ на основание $\mathrm{Lin}(x_1, x_2)$, находится по формуле:
\[
h = \frac{V_3}{S_{\text{осн}}} = \sqrt{\frac{\Gamma(x_1, x_2, y)}{\Gamma(x_1, x_2)}}.
\]
\end{theory}

\step{1} Вычислим попарные скалярные произведения векторов:
\[
\langle x_1, x_1\rangle = 1^2 + 1^2 + 0 + 0 = 2,
\]
\[
\langle x_2, x_2\rangle = 0 + 1^2 + 1^2 + 0 = 2,
\]
\[
\langle y, y\rangle = 1^2 + 1^2 + 1^2 + 1^2 = 4,
\]
\[
\langle x_1, x_2\rangle = 1\cdot 0 + 1\cdot 1 + 0\cdot 1 + 0\cdot 0 = 1,
\]
\[
\langle x_1, y\rangle = 1\cdot 1 + 1\cdot 1 + 0 + 0 = 2,
\]
\[
\langle x_2, y\rangle = 0\cdot 1 + 1\cdot 1 + 1\cdot 1 + 0 = 2.
\]

\step{2} Составим и вычислим определитель Грама для системы $\{x_1, x_2\}$:
\[
G_2 = \begin{pmatrix} \langle x_1, x_1\rangle & \langle x_1, x_2\rangle \\ \langle x_2, x_1\rangle & \langle x_2, x_2\rangle \end{pmatrix} = \begin{pmatrix} 2 & 1 \\ 1 & 2 \end{pmatrix}.
\]
\[
\Gamma(x_1, x_2) = \det G_2 = 2\cdot 2 - 1\cdot 1 = 3.
\]
Площадь основания параллелепипеда равна $S_{\text{осн}} = \sqrt{3}$.

\step{3} Составим и вычислим определитель Грама для системы трех векторов $\{x_1, x_2, y\}$:
\[
G_3 = \begin{pmatrix} 2 & 1 & 2 \\ 1 & 2 & 2 \\ 2 & 2 & 4 \end{pmatrix}.
\]
Для упрощения вычислений вычтем из третьей строки первую, умноженную на 2, а затем разложим:
\[
\det G_3 = \begin{vmatrix} 2 & 1 & 2 \\ 1 & 2 & 2 \\ -2 & 0 & 0 \end{vmatrix} = -2 \cdot \begin{vmatrix} 1 & 2 \\ 2 & 2 \end{vmatrix} = -2 \cdot (2 - 4) = 4.
\]
Квадрат трёхмерного объёма равен $V_3^2 = 4$, следовательно, объём параллелепипеда $V_3 = 2$.

\step{4} Найдем высоту $h$:
\[
h = \sqrt{\frac{\Gamma(x_1, x_2, y)}{\Gamma(x_1, x_2)}} = \sqrt{\frac{4}{3}} = \frac{2}{\sqrt{3}} = \frac{2\sqrt{3}}{3}.
\]

\answer
1. Определители Грама: $\Gamma(x_1, x_2) = 3$, $\Gamma(x_1, x_2, y) = 4$.\\
2. Высота параллелепипеда: $h = \dfrac{2\sqrt{3}}{3}$.
%=============================================================
\task{26}
%=============================================================
\subsection*{Условие}
В евклидовом пространстве многочленов степени не выше второй $\mathbb{P}_2[0, 1]$ со скалярным произведением $\displaystyle \langle f, g\rangle = \int_0^1 f(x)g(x)\,dx$ задан стандартный «косой» базис: $f_1(x) = 1$, $f_2(x) = x$, $f_3(x) = x^2$.
\begin{enumerate}
    \item Примените алгоритм Грама — Шмидта к этой системе функций и постройте ортогональный базис $\{g_1(x), g_2(x), g_3(x)\}$.
    \item Нормируйте полученные функции, чтобы получить ортонормированный базис (ОНБ) этого пространства.
\end{enumerate}

\begin{theory}
\textbf{Ортогонализация Грама — Шмидта.} По исходному базису $\{f_1, f_2, f_3\}$ строим ортогональную систему $\{g_1, g_2, g_3\}$ по формулам:
\[
g_1 = f_1, \qquad g_2 = f_2 - \frac{\langle f_2, g_1\rangle}{\|g_1\|^2}g_1, \qquad g_3 = f_3 - \frac{\langle f_3, g_1\rangle}{\|g_1\|^2}g_1 - \frac{\langle f_3, g_2\rangle}{\|g_2\|^2}g_2.
\]
Нормированные векторы: $e_i = g_i / \|g_i\|$, где $\|g_i\|^2 = \langle g_i, g_i\rangle$.
\end{theory}

\step{1} Найдём первый ортогональный вектор $g_1(x)$ и его квадрат нормы:
\[
g_1(x) = f_1(x) = 1, \qquad \|g_1\|^2 = \int_0^1 1\,dx = 1.
\]

\step{2} Найдём второй ортогональный вектор $g_2(x)$. Вычислим скалярное произведение:
\[
\langle f_2, g_1\rangle = \int_0^1 x \cdot 1\,dx = \frac{1}{2}.
\]
Тогда:
\[
g_2(x) = x - \frac{1/2}{1} \cdot 1 = x - \frac{1}{2}.
\]
Вычислим его квадрат нормы:
\[
\|g_2\|^2 = \int_0^1 \left(x - \frac{1}{2}\right)^2 dx = \int_0^1 \left(x^2 - x + \frac{1}{4}\right) dx = \frac{1}{3} - \frac{1}{2} + \frac{1}{4} = \frac{1}{12}.
\]

\step{3} Найдём третий ортогональный вектор $g_3(x)$. Вычислим скалярные произведения:
\[
\langle f_3, g_1\rangle = \int_0^1 x^2 \cdot 1\,dx = \frac{1}{3},
\]
\[
\langle f_3, g_2\rangle = \int_0^1 x^2 \left(x - \frac{1}{2}\right) dx = \int_0^1 \left(x^3 - \frac{1}{2}x^2\right) dx = \frac{1}{4} - \frac{1}{6} = \frac{1}{12}.
\]
Тогда:
\[
g_3(x) = x^2 - \frac{1/3}{1} \cdot 1 - \frac{1/12}{1/12} \left(x - \frac{1}{2}\right) = x^2 - \frac{1}{3} - \left(x - \frac{1}{2}\right) = x^2 - x + \frac{1}{6}.
\]
Вычислим квадрат нормы $g_3(x)$:
\[
\|g_3\|^2 = \int_0^1 \left(x^2 - x + \frac{1}{6}\right)^2 dx = \int_0^1 \left(x^4 - 2x^3 + \frac{4}{3}x^2 - \frac{1}{3}x + \frac{1}{36}\right) dx
\]
\[
= \frac{1}{5} - \frac{1}{2} + \frac{4}{9} - \frac{1}{6} + \frac{1}{36} = \frac{36 - 90 + 80 - 30 + 5}{180} = \frac{1}{180}.
\]

\step{4} Нормируем полученный базис:
\[
e_1(x) = \frac{g_1(x)}{\|g_1\|} = 1,
\]
\[
e_2(x) = \frac{g_2(x)}{\|g_2\|} = \sqrt{12}\left(x - \frac{1}{2}\right) = \sqrt{3}(2x - 1),
\]
\[
e_3(x) = \frac{g_3(x)}{\|g_3\|} = \sqrt{180}\left(x^2 - x + \frac{1}{6}\right) = \sqrt{5}(6x^2 - 6x + 1).
\]

\answer
Ортогональный базис:
\[
g_1(x) = 1, \qquad g_2(x) = x - \frac{1}{2}, \qquad g_3(x) = x^2 - x + \frac{1}{6}.
\]
Ортонормированный базис:
\[
e_1(x) = 1, \qquad e_2(x) = \sqrt{3}(2x - 1), \qquad e_3(x) = \sqrt{5}(6x^2 - 6x + 1).
\]
(Построенные многочлены представляют собой смещенные многочлены Лежандра).

%=============================================================
\task{27}
%=============================================================
\subsection*{Условие}
В евклидовом пространстве непрерывных функций $C[-1, 1]$ со стандартным скалярным произведением $\displaystyle \langle f, g\rangle = \int_{-1}^1 f(x)g(x)\,dx$ задано подпространство линейных функций $L = L(1, x)$. Рассмотрим функцию $f(x) = x^2$.
\begin{enumerate}
    \item Найдите ортогональную проекцию $f_L(x)$ и ортогональную составляющую $f_{L^\perp}(x)$ функции $f(x)$ относительно подпространства $L$.
    \item Вычислите евклидову норму ортогональной составляющей $\|f_{L^\perp}\|$. Какой глубокий геометрический смысл имеет эта величина в контексте аппроксимации функций?
\end{enumerate}

\begin{theory}
Если подпространство $L$ натянуто на ортогональный базис $\{\varphi_1, \varphi_2\}$, то ортогональная проекция вектора $f$ на $L$ вычисляется по формуле:
\[
f_L = \frac{\langle f, \varphi_1\rangle}{\|\varphi_1\|^2}\varphi_1 + \frac{\langle f, \varphi_2\rangle}{\|\varphi_2\|^2}\varphi_2.
\]
Ортогональная составляющая: $f_{L^\perp} = f - f_L$. Её норма $\|f_{L^\perp}\|$ выражает наименьшее расстояние от элемента $f$ до подпространства $L$.
\end{theory}

\step{1} В качестве базиса $L$ возьмём $\varphi_1(x) = 1$ и $\varphi_2(x) = x$. Они ортогональны на $[-1, 1]$, так как:
\[
\langle 1, x\rangle = \int_{-1}^1 x\,dx = 0.
\]
Вычислим квадраты их норм:
\[
\|\varphi_1\|^2 = \int_{-1}^1 1\,dx = 2, \qquad \|\varphi_2\|^2 = \int_{-1}^1 x^2\,dx = \frac{2}{3}.
\]

\step{2} Вычислим коэффициенты проекции для $f(x) = x^2$:
\[
\langle x^2, 1\rangle = \int_{-1}^1 x^2\,dx = \frac{2}{3}, \qquad \langle x^2, x\rangle = \int_{-1}^1 x^3\,dx = 0.
\]
Подставляем в формулу проекции:
\[
f_L(x) = \frac{2/3}{2} \cdot 1 + \frac{0}{2/3} \cdot x = \frac{1}{3}.
\]
Тогда ортогональная составляющая равна:
\[
f_{L^\perp}(x) = f(x) - f_L(x) = x^2 - \frac{1}{3}.
\]

\step{3} Вычислим норму ортогональной составляющей $f_{L^\perp}(x)$:
\[
\|f_{L^\perp}\|^2 = \int_{-1}^1 \left(x^2 - \frac{1}{3}\right)^2 dx = \int_{-1}^1 \left(x^4 - \frac{2}{3}x^2 + \frac{1}{9}\right) dx
\]
\[
= \left[ \frac{x^5}{5} - \frac{2}{9}x^3 + \frac{1}{9}x \right]_{-1}^1 = 2 \left( \frac{1}{5} - \frac{2}{9} + \frac{1}{9} \right) = 2 \left( \frac{1}{5} - \frac{1}{9} \right) = \frac{8}{45}.
\]
\[
\|f_{L^\perp}\| = \sqrt{\frac{8}{45}} = \frac{2\sqrt{2}}{3\sqrt{5}} = \frac{2\sqrt{10}}{15}.
\]

\step{4} \textbf{Геометрический смысл.} Величина $\|f_{L^\perp}\|$ представляет собой наименьшее расстояние в смысле среднего квадратичного отклонения от функции $f(x) = x^2$ до множества всех линейных функций на отрезке $[-1, 1]$. То есть константа $y = 1/3$ является наилучшим среднеквадратичным приближением (аппроксимацией) параболы линейной функцией на этом отрезке.

\answer
1. Проекция: $f_L(x) = \dfrac{1}{3}$, ортогональная составляющая: $f_{L^\perp}(x) = x^2 - \dfrac{1}{3}$.\\
2. Норма $\|f_{L^\perp}\| = \dfrac{2\sqrt{10}}{15}$. Физический/геометрический смысл: это минимальная среднеквадратичная ошибка при аппроксимации параболы линейными функциями на $[-1, 1]$.

%=============================================================
\task{28}
%=============================================================
\subsection*{Условие}
Ортогональное дополнение. Подпространство $L \subset \R^4$ задано системой линейных уравнений:
\[
\begin{cases}
x_1 + 2x_2 + x_3 = 0, \\
x_2 + x_3 + x_4 = 0.
\end{cases}
\]
Найти базис ортогонального дополнения $L^\perp$. Какова его размерность?

\begin{theory}
Подпространство $L$ можно трактовать как ядро (пространство решений) матрицы системы $A$:
\[
A = \begin{pmatrix} 1 & 2 & 1 & 0 \\ 0 & 1 & 1 & 1 \end{pmatrix}.
\]
В силу свойств евклидовых пространств, ортогональное дополнение к ядру матрицы совпадает с пространством, натянутым на её строки (транспонированным образом): $L^\perp = \mathrm{Im}(A^T)$. Базис $L^\perp$ образуют линейно независимые строки матрицы $A$. Размерности связаны соотношением: $\dim L + \dim L^\perp = n$.
\end{theory}

\step{1} Запишем вектор-строки, задающие коэффициенты уравнений связи:
\[
u_1 = (1, 2, 1, 0), \qquad u_2 = (0, 1, 1, 1).
\]
Любой вектор $x \in L$ ортогонален строкам матрицы системы, так как уравнения системы представляют собой в точности скалярные произведения $\langle x, u_1\rangle = 0$ и $\langle x, u_2\rangle = 0$.

\step{2} Следовательно, векторы $u_1$ и $u_2$ принадлежат ортогональному дополнению $L^\perp$. Поскольку строки матрицы $A$ линейно независимы (первая имеет ведущую единицу в 1-м столбце, вторая — во 2-м), они образуют базис пространства строк.

\step{3} Размерность пространства $L^\perp$ равна рангу матрицы коэффициентов $A$:
\[
\dim L^\perp = \mathrm{rg}(A) = 2.
\]
(При этом размерность исходного подпространства $L$ равна $\dim L = 4 - 2 = 2$).

\answer
Базис ортогонального дополнения $L^\perp$:
\[
u_1 = (1, 2, 1, 0), \qquad u_2 = (0, 1, 1, 1).
\]
Размерность ортогонального дополнения: $\dim L^\perp = 2$.

%=============================================================
\task{29}
%=============================================================
\subsection*{Условие}
В пространстве $C[0, 1]$ со скалярным произведением $(f, g) = \int_0^1 f(t)g(t)\,dt$ найти расстояние от функции $f(t) = t^2$ до подпространства $L = \mathrm{span}\{1, t\}$. Интерпретировать результат как среднеквадратичное приближение.

\begin{theory}
Расстояние от функции $f(t)$ до подпространства $L$ вычисляется как норма ортогональной составляющей $f_{L^\perp}(t) = f(t) - f_L(t)$, где $f_L(t)$ — наилучшая среднеквадратичная аппроксимация (проекция) функции $f(t)$ элементами из $L$.
\end{theory}

\step{1} Ранее (в Задаче 12) было получено, что ортогональной проекцией функции $f(t) = t^2$ на линейную оболочку $\mathrm{span}\{1, t\}$ является многочлен первой степени:
\[
p(t) = t - \frac{1}{6}.
\]

\step{2} Тогда ортогональная составляющая функции $f(t)$ имеет вид:
\[
f_{L^\perp}(t) = t^2 - t + \frac{1}{6}.
\]
Эта функция в точности совпадает со вторым ортогональным многочленом Лежандра $g_3(t)$, полученным в Задаче 26.

\step{3} Вычислим квадрат нормы (который также был рассчитан в шаге 3 задачи 26):
\[
\|f_{L^\perp}\|^2 = \int_0^1 \left(t^2 - t + \frac{1}{6}\right)^2 dt = \frac{1}{180}.
\]
Тогда расстояние от $f(t)$ до $L$ равно:
\[
d(f, L) = \|f_{L^\perp}\| = \sqrt{\frac{1}{180}} = \frac{1}{6\sqrt{5}} = \frac{\sqrt{5}}{30}.
\]

\step{4} \textbf{Интерпретация.} Данное расстояние является среднеквадратичной ошибкой наилучшего приближения параболы $t^2$ прямой линией на отрезке $[0, 1]$. Отрезок прямой $y = t - 1/6$ обеспечивает минимально возможное отклонение в смысле метрики $L^2[0, 1]$.

\answer
Расстояние равно $d(f, L) = \dfrac{\sqrt{5}}{30}$. Физически это величина минимально возможного среднеквадратичного уклонения при замене квадратичной зависимости линейной на интервале $[0,1]$.

%=============================================================
\task{30}
%=============================================================
\subsection*{Условие}
В евклидовом пространстве непрерывных функций $C[0, 1]$ со стандартным скалярным произведением $\langle f, g\rangle = \int_0^1 f(x)g(x)\,dx$ рассматривается функция $f(x) = x^3$ и подпространство многочленов первой степени $L = L(1, x)$.
\begin{enumerate}
    \item Найдите такие вещественные коэффициенты $a$ и $b$, при которых значение интеграла
    \[
    I(a, b) = \int_0^1 (x^3 - ax - b)^2 dx
    \]
    принимает минимально возможное значение.
    \item Объясните алгебраический смысл этой задачи. Вычислите минимальное значение интеграла $I_{\min}$, связав его с расстоянием от вектора до подпространства.
\end{enumerate}

\step{1} Нахождение минимума интеграла эквивалентно поиску ортогональной проекции функции $f(x) = x^3$ на подпространство $L = \mathrm{span}\{1, x\}$. Обозначим проекцию как $P_L(x^3) = ax + b$. По условию ортогональности, разность $x^3 - (ax + b)$ должна быть перпендикулярна базисным функциям $1$ и $x$:
\[
\begin{cases}
\langle x^3 - ax - b, 1\rangle = 0, \\
\langle x^3 - ax - b, x\rangle = 0.
\end{cases}
\]

\step{2} Запишем систему в виде интегралов:
\[
\begin{cases}
\displaystyle \int_0^1 (x^3 - ax - b)\,dx = 0 \iff \frac{1}{4} - \frac{a}{2} - b = 0, \\
\displaystyle \int_0^1 x(x^3 - ax - b)\,dx = 0 \iff \frac{1}{5} - \frac{a}{3} - b = 2 \cdot \frac{b}{2} = \frac{b}{2} = 0.
\end{cases}
\]
Получаем систему линейных уравнений для $a$ и $b$:
\[
\begin{cases}
2a + 4b = 1, \\
10a + 15b = 6.
\end{cases}
\]

\step{3} Решим систему. Из первого уравнения выразим $a = \frac{1 - 4b}{2} = 0.5 - 2b$. Подставим во второе:
\[
10(0.5 - 2b) + 15b = 6 \implies 5 - 20b + 15b = 6 \implies -5b = 1 \implies b = -0.2 = -\frac{1}{5}.
\]
Тогда:
\[
a = 0.5 - 2(-0.2) = 0.5 + 0.4 = 0.9 = \frac{9}{10}.
\]
Искомая проекция: $P_L(x^3) = 0.9x - 0.2$.

\step{4} \textbf{Алгебраический смысл.} Минимизация интеграла $I(a, b)$ — это задача нахождения наилучшего приближения (ортогональной проекции) вектора $x^3$ на подпространство линейных функций $L = \mathrm{span}\{1, x\}$. Минимальное значение интеграла равно квадрату расстояния от вектора до подпространства: $I_{\min} = d^2(f, L) = \|f_{L^\perp}\|^2$.

\step{5} Вычислим минимальное значение интеграла $I_{\min}$, используя теорему Пифагора: $\|f_{L^\perp}\|^2 = \|f\|^2 - \|f_L\|^2$:
\[
\|f\|^2 = \int_0^1 (x^3)^2 dx = \int_0^1 x^6 dx = \frac{1}{7},
\]
\[
\|f_L\|^2 = \int_0^1 (0.9x - 0.2)^2 dx = \int_0^1 (0.81x^2 - 0.36x + 0.04) dx 
\]
\[
= \frac{0.81}{3} - \frac{0.36}{2} + 0.04 = 0.27 - 0.18 + 0.04 = 0.13 = \frac{13}{100}.
\]
Тогда минимальное значение интеграла:
\[
I_{\min} = \frac{1}{7} - \frac{13}{100} = \frac{100 - 91}{700} = \frac{9}{700}.
\]

\answer
1. Минимизирующие коэффициенты: $a = \dfrac{9}{10} = 0.9$, $b = -\dfrac{1}{5} = -0.2$.\\
2. Алгебраический смысл задачи состоит в нахождении ортогональной проекции вектора на подпространство. Минимальное значение интеграла: $I_{\min} = \dfrac{9}{700}$ (квадрат расстояния).
%=============================================================
\task{31}
%=============================================================
\subsection*{Условие}
Пусть $P$ — оператор ортогонального проектирования на подпространство $L$. Доказать, что для любого вектора $x$ выполняется неравенство $\|Px\| \leq \|x\|$. В каком случае достигается равенство?

\begin{theory}
\textbf{Теорема о разложении пространства.} Любой вектор $x$ евклидова пространства может быть единственным образом представлен в виде прямой суммы:
\[
x = Px + (x - Px), \qquad \text{где } Px \in L, \quad (x - Px) \in L^\perp.
\]
Поскольку проекция $Px$ и ортогональная составляющая $x - Px$ взаимно ортогональны, для них выполняется теорема Пифагора:
\[
\|x\|^2 = \|Px\|^2 + \|x - Px\|^2.
\]
\end{theory}

\step{1} Запишем равенство, следующее из теоремы Пифагора:
\[
\|Px\|^2 = \|x\|^2 - \|x - Px\|^2.
\]

\step{2} Поскольку норма любого вектора неотрицательна, квадрат нормы ортогональной составляющей $\|x - Px\|^2 \geq 0$. Следовательно:
\[
\|Px\|^2 \leq \|x\|^2 \implies \|Px\| \leq \|x\|. \quad\checkmark
\]

\step{3} Исследуем случай достижения равенства $\|Px\| = \|x\|$. Это возможно тогда и только тогда, когда:
\[
\|x - Px\|^2 = 0 \iff x - Px = 0 \iff x = Px \iff x \in L.
\]

\answer
Неравенство доказано. Равенство достигается тогда и только тогда, когда вектор $x$ изначально лежит в подпространстве проектирования $L$ (то есть $x \in L$).

%=============================================================
\task{32}
%=============================================================
\subsection*{Условие}
Найти объем четырехмерного параллелепипеда в $\R^5$, построенного на векторах $v_1, v_2, v_3, v_4$, если известна их матрица Грама $G$. Как изменится объем, если один из векторов заменить на его ортогональную составляющую относительно остальных?

\begin{theory}
Объем $k$-мерного параллелепипеда, построенного на линейно независимой системе векторов $\{v_1, \ldots, v_k\}$, равен квадратному корню из определителя их матрицы Грама:
\[
V_k = \sqrt{\det G(v_1, \ldots, v_k)}.
\]
Ортогональная составляющая вектора $v_4$ относительно подпространства $L = \mathrm{span}\{v_1, v_2, v_3\}$ имеет вид $h_4 = v_4 - P_L(v_4)$, где $P_L(v_4) \in L$.
\end{theory}

\step{1} По определению геометрического смысла грамиана, объем четырехмерного параллелепипеда равен:
\[
V_4 = \sqrt{\det G}.
\]

\step{2} Рассмотрим замену вектора $v_4$ на его ортогональную составляющую $h_4 = v_4 - P_L(v_4)$. Поскольку проекция $P_L(v_4)$ является линейной комбинацией остальных векторов $v_1, v_2, v_3$, данное преобразование эквивалентно прибавлению к одному из столбцов (и одновременно к строке в матрице Грама) линейной комбинации остальных.

\step{3} Из свойств определителя и элементарных преобразований известно, что такие сдвиги (сдвиговые деформации/транс shear) не меняют определитель матрицы Грама. 
Геометрически объем параллелепипеда равен произведению «площади» трехмерного основания на высоту:
\[
V_4 = V_3(v_1, v_2, v_3) \cdot \|h_4\|.
\]
Если мы заменим $v_4$ на $h_4$, то новое основание останется прежним, а новый вектор будет уже перпендикулярен основанию, то есть его норма в точности совпадет с высотой $\|h_4\|$. Таким образом, объем не изменится.

\answer
Объем параллелепипеда равен $V_4 = \sqrt{\det G}$. При замене любого из векторов на его ортогональную составляющую относительно остальных векторов объем параллелепипеда \textbf{не изменится}.

%=============================================================
\task{33}
%=============================================================
\subsection*{Условие}
Дан вектор $x = (1, 2, 3, 4)$. Найти его ортогональную проекцию на гиперплоскость $x_1 - x_2 + x_3 - x_4 = 0$ и соответствующую ортогональную составляющую. Проверить выполнение теоремы Пифагора для полученных векторов.

\begin{theory}
Гиперплоскость $L$ задана уравнением $\langle x, n\rangle = 0$, где $n = (1, -1, 1, -1)$ — вектор нормали. Ортогональная составляющая вектора $x$ относительно $L$ — это его проекция на одномерное подпространство $L^\perp$, натянутое на нормаль $n$:
\[
x_{L^\perp} = \frac{\langle x, n\rangle}{\|n\|^2} n.
\]
Тогда проекция на саму гиперплоскость $L$ равна: $x_L = x - x_{L^\perp}$.
\end{theory}

\step{1} Вычислим скалярное произведение $\langle x, n\rangle$ и квадрат нормы вектора $n$:
\[
\langle x, n\rangle = 1\cdot 1 + 2\cdot(-1) + 3\cdot 1 + 4\cdot(-1) = 1 - 2 + 3 - 4 = -2,
\]
\[
\|n\|^2 = 1^2 + (-1)^2 + 1^2 + (-1)^2 = 4.
\]

\step{2} Найдем ортогональную составляющую $x_{L^\perp}$:
\[
x_{L^\perp} = \frac{-2}{4} (1, -1, 1, -1) = \left(-\frac{1}{2}, \frac{1}{2}, -\frac{1}{2}, \frac{1}{2}\right).
\]

\step{3} Вычислим ортогональную проекцию $x_L$ на плоскость:
\[
x_L = x - x_{L^\perp} = (1, 2, 3, 4) - \left(-\frac{1}{2}, \frac{1}{2}, -\frac{1}{2}, \frac{1}{2}\right) = \left(\frac{3}{2}, \frac{3}{2}, \frac{7}{2}, \frac{7}{2}\right).
\]
Проверим принадлежность $x_L$ плоскости $L$: $\frac{3}{2} - \frac{3}{2} + \frac{7}{2} - \frac{7}{2} = 0$. Условие выполнено.

\step{4} Проверим теорему Пифагора $\|x\|^2 = \|x_L\|^2 + \|x_{L^\perp}\|^2$:
\[
\|x\|^2 = 1^2 + 2^2 + 3^2 + 4^2 = 30,
\]
\[
\|x_{L^\perp}\|^2 = \left(-\frac{1}{2}\right)^2 + \left(\frac{1}{2}\right)^2 + \left(-\frac{1}{2}\right)^2 + \left(\frac{1}{2}\right)^2 = \frac{1}{4} \cdot 4 = 1,
\]
\[
\|x_L\|^2 = \left(\frac{3}{2}\right)^2 + \left(\frac{3}{2}\right)^2 + \left(\frac{7}{2}\right)^2 + \left(\frac{7}{2}\right)^2 = \frac{9 + 9 + 49 + 49}{4} = \frac{116}{4} = 29.
\]
Поскольку $29 + 1 = 30$, теорема Пифагора выполняется. \quad$\checkmark$

\answer
Ортогональная проекция: $x_L = \left(\dfrac{3}{2}, \dfrac{3}{2}, \dfrac{7}{2}, \dfrac{7}{2}\right)$. \\
Ортогональная составляющая: $x_{L^\perp} = \left(-\dfrac{1}{2}, \dfrac{1}{2}, -\dfrac{1}{2}, \dfrac{1}{2}\right)$. \\
Теорема Пифагора проверена: $\|x\|^2 = 30$, $\|x_L\|^2 = 29$, $\|x_{L^\perp}\|^2 = 1$.

%=============================================================
\task{34}
%=============================================================
\subsection*{Условие}
Линейный оператор $\mathcal{A} : \mathbb{R}^2 \to \mathbb{R}^3$ переводит векторы $e_1 = (1, 0)$ и $e_2 = (0, 1)$ в векторы $a_1 = (1, 2, 0)$ и $a_2 = (0, 1, 1)$ соответственно. Найти образ вектора $x = (3, -2)$ и составить матрицу оператора в стандартных базисах.

\begin{theory}
Столбцами матрицы линейного оператора в стандартных базисах являются координаты образов базисных векторов стандартного базиса. Если $\mathcal{A}(e_1) = a_1$ и $\mathcal{A}(e_2) = a_2$, то матрица $A$ имеет вид $A = \begin{pmatrix} a_1 & a_2 \end{pmatrix}$. Образ любого вектора $x$ вычисляется как произведение матрицы на вектор-столбец: $y = A \cdot x$.
\end{theory}

\step{1} Составим матрицу оператора $A$. Поскольку $e_1, e_2$ — стандартный базис пространства $\R^2$, запишем образы $a_1, a_2$ в виде столбцов:
\[
A = \begin{pmatrix} 1 & 0 \\ 2 & 1 \\ 0 & 1 \end{pmatrix}.
\]

\step{2} Найдем образ вектора $x = (3, -2)^T$, используя линейность оператора:
\[
\mathcal{A}(x) = 3\mathcal{A}(e_1) - 2\mathcal{A}(e_2) = 3(1, 2, 0)^T - 2(0, 1, 1)^T = (3, 6, 0)^T - (0, 2, 2)^T = (3, 4, -2)^T.
\]

\step{3} Проверим результат матричным умножением:
\[
A \cdot x = \begin{pmatrix} 1 & 0 \\ 2 & 1 \\ 0 & 1 \end{pmatrix} \begin{pmatrix} 3 \\ -2 \end{pmatrix} = \begin{pmatrix} 1\cdot 3 + 0\cdot(-2) \\ 2\cdot 3 + 1\cdot(-2) \\ 0\cdot 3 + 1\cdot(-2) \end{pmatrix} = \begin{pmatrix} 3 \\ 4 \\ -2 \end{pmatrix}. \quad\checkmark
\]

\answer
Матрица оператора: $A = \begin{pmatrix} 1 & 0 \\ 2 & 1 \\ 0 & 1 \end{pmatrix}$. \\
Образ вектора: $\mathcal{A}(x) = (3, 4, -2)$.

%=============================================================
\task{35}
%=============================================================
\subsection*{Условие}
В линейном пространстве $\mathbb{R}^3$ задано отображение $A : \mathbb{R}^3 \to \mathbb{R}^3$, которое переводит произвольный вектор $x = (x_1, x_2, x_3)$ в вектор $y = A(x)$ по закону:
\[
A(x) = (2x_1 - x_2,\, x_2 + 3x_3,\, x_1 \cdot x_3).
\]
\begin{enumerate}
    \item Проверить выполнение аксиом аддитивности и однородности для данного отображения. Является ли оно линейным оператором?
    \item Рассмотреть другое отображение $B(x) = (2x_1 - x_2,\, x_2 + 3x_3,\, x_1 + 1)$. Выполняется ли для него условие $B(\vec{0}) = \vec{0}$? Является ли оно линейным оператором? Обосновать ответ.
\end{enumerate}

\begin{theory}
Отображение называется линейным оператором, если выполнены две аксиомы:
\begin{enumerate}
    \item \textbf{Аддитивность:} $A(u + v) = A(u) + A(v)$ для любых $u, v$.
    \item \textbf{Однородность:} $A(\alpha u) = \alpha A(u)$ для любого числа $\alpha$ и вектора $u$.
\end{enumerate}
Любой линейный оператор обязательно переводит нулевой вектор в нулевой: $A(\vec{0}) = \vec{0}$.
\end{theory}

\subsubsection*{Часть 1. Отображение $A(x)$}

\step{1} Проверим аддитивность для векторов $u = (u_1, u_2, u_3)$ и $v = (v_1, v_2, v_3)$:
\[
A(u+v) = \Big(2(u_1+v_1) - (u_2+v_2),\, (u_2+v_2) + 3(u_3+v_3),\, (u_1+v_1)(u_3+v_3)\Big).
\]
Третья координата суммы: $(u_1+v_1)(u_3+v_3) = u_1 u_3 + v_1 v_3 + u_1 v_3 + v_1 u_3$. \\
В то же время, третья координата вектора $A(u) + A(v)$ равна $u_1 u_3 + v_1 v_3$. \\
Поскольку $u_1 v_3 + v_1 u_3 \neq 0$ в общем случае, аксиома аддитивности \textbf{не выполняется} (например, при $u=(1,0,1)$ и $v=(1,0,1)$).

\step{2} Проверим однородность:
\[
A(\alpha x) = \Big(2(\alpha x_1) - \alpha x_2,\, \alpha x_2 + 3(\alpha x_3),\, (\alpha x_1)(\alpha x_3)\Big) = \Big(\alpha(2x_1 - x_2),\, \alpha(x_2 + 3x_3),\, \alpha^2(x_1 x_3)\Big).
\]
Для третьей координаты $\alpha^2(x_1 x_3) \neq \alpha(x_1 x_3)$ при $\alpha \neq 1, 0$. Однородность \textbf{не выполняется}. Отображение $A$ не является линейным.

\subsubsection*{Часть 2. Отображение $B(x)$}

\step{3} Проверим условие сохранения нуля для $B(x)$:
\[
B(\vec{0}) = B(0, 0, 0) = (2\cdot 0 - 0,\, 0 + 3\cdot 0,\, 0 + 1) = (0, 0, 1) \neq (0, 0, 0).
\]

\step{4} Так как $B(\vec{0}) \neq \vec{0}$, отображение $B$ заведомо \textbf{не является} линейным оператором (нарушается свойство однородности при $\alpha = 0$, поскольку $B(0 \cdot x) = B(\vec{0}) = (0, 0, 1)$, но $0 \cdot B(x) = (0, 0, 0)$).

\answer
1. Для отображения $A(x)$ аксиомы аддитивности и однородности не выполняются. Оно \textbf{не является} линейным оператором.\\
2. Для отображения $B(x)$ условие $B(\vec{0}) = \vec{0}$ \textbf{не выполняется} (так как $B(\vec{0}) = (0, 0, 1)$). Оно также \textbf{не является} линейным оператором.

%=============================================================
\task{36}
%=============================================================
\subsection*{Условие}
В пространстве многочленов $\R[x]_3$ задан оператор $\mathcal{A}p(x) = p''(x) + p'(x)$. Найти базисы ядра и образа этого оператора, а также его ранг и дефект. Проверить выполнение теоремы о сумме размерностей.

\begin{theory}
Пусть $V$ — линейное пространство размерности $n$. 
\begin{itemize}
    \item \textbf{Ядро оператора} $\mathrm{ker}\,\mathcal{A} = \{p \in V : \mathcal{A}p = 0\}$, дефект $\mathrm{def}\,\mathcal{A} = \dim \mathrm{ker}\,\mathcal{A}$.
    \item \textbf{Образ оператора} $\mathrm{im}\,\mathcal{A} = \{q \in V : \exists p \in V, \mathcal{A}p = q\}$, ранг $\mathrm{rg}\,\mathcal{A} = \dim \mathrm{im}\,\mathcal{A}$.
    \item \textbf{Теорема о сумме размерностей (rank-nullity theorem):} $\mathrm{rg}\,\mathcal{A} + \mathrm{def}\,\mathcal{A} = n$.
\end{itemize}
В пространстве $\R[x]_3$ базисом является $\{1, x, x^2, x^3\}$, следовательно, $n = 4$.
\end{theory}

\step{1} Найдём ядро оператора. Пусть $p(x) = a + bx + cx^2 + dx^3$. Вычислим производные:
\[
p'(x) = b + 2cx + 3dx^2, \qquad p''(x) = 2c + 6dx.
\]
Уравнение $\mathcal{A}p(x) = 0$ принимает вид:
\[
(2c + b) + (6d + 2c)x + 3dx^2 = 0 \cdot x^2 + 0 \cdot x + 0.
\]
Приравнивая коэффициенты при степенях $x$ к нулю, получаем систему:
\[
\begin{cases}
3d = 0 \implies d = 0, \\
6d + 2c = 0 \implies c = 0, \\
2c + b = 0 \implies b = 0.
\end{cases}
\]
Коэффициент $a$ свободный. Таким образом, $p(x) = a$ — константа. \\
Базис ядра $\mathrm{ker}\,\mathcal{A}$ состоит из одного элемента $\{1\}$. Дефект равен $\mathrm{def}\,\mathcal{A} = 1$.

\step{2} Найдём образ оператора. Образ натянут на систему образов базисных векторов:
\[
\mathcal{A}(1) = 0,
\]
\[
\mathcal{A}(x) = 1,
\]
\[
\mathcal{A}(x^2) = 2x + 2,
\]
\[
\mathcal{A}(x^3) = 3x^2 + 6x.
\]
Система $\{1, 2x+2, 3x^2+6x\}$ эквивалентна стандартному базису пространства $\R[x]_2$, то есть $\{1, x, x^2\}$. \\
Следовательно, базис образа $\mathrm{im}\,\mathcal{A}$ — это $\{1, x, x^2\}$. Ранг равен $\mathrm{rg}\,\mathcal{A} = 3$.

\step{3} Проверим теорему о сумме размерностей для $n = \dim \R[x]_3 = 4$:
\[
\mathrm{rg}\,\mathcal{A} + \mathrm{def}\,\mathcal{A} = 3 + 1 = 4. \quad\checkmark
\]

\answer
Базис ядра: $\{1\}$; дефект $\mathrm{def}\,\mathcal{A} = 1$.\\
Базис образа: $\{1, x, x^2\}$; ранг $\mathrm{rg}\,\mathcal{A} = 3$.\\
Теорема о сумме размерностей подтверждена ($3 + 1 = 4$).

%=============================================================
\task{37}
%=============================================================
\subsection*{Условие}
Найти матрицу линейного оператора $\mathcal{A}p(x) = p(x + 1)$ в базисе $\{1, x, x^2\}$ пространства $\R[x]_2$. Является ли этот оператор обратимым?

\begin{theory}
Матрица оператора составляется по столбцам: $j$-й столбец матрицы — это координаты образа $j$-го базисного вектора в том же базисе. Оператор обратим тогда и только тогда, когда его матрица невырождена ($\det A \neq 0$).
\end{theory}

\step{1} Найдём образы базисных многочленов $\varphi_1 = 1$, $\varphi_2 = x$, $\varphi_3 = x^2$:
\[
\mathcal{A}(1) = 1,
\]
\[
\mathcal{A}(x) = x + 1 = 1\cdot 1 + 1\cdot x + 0\cdot x^2,
\]
\[
\mathcal{A}(x^2) = (x + 1)^2 = x^2 + 2x + 1 = 1\cdot 1 + 2\cdot x + 1\cdot x^2.
\]

\step{2} Запишем полученные коэффициенты в столбцы матрицы $A$:
\[
A = \begin{pmatrix} 1 & 1 & 1 \\ 0 & 1 & 2 \\ 0 & 0 & 1 \end{pmatrix}.
\]

\step{3} Проверим обратимость оператора. Определитель верхнетреугольной матрицы равен произведению её диагональных элементов:
\[
\det A = 1 \cdot 1 \cdot 1 = 1 \neq 0.
\]
Поскольку определитель отличен от нуля, оператор $\mathcal{A}$ обратим (обратным является оператор сдвига назад: $\mathcal{A}^{-1}p(x) = p(x-1)$).

\answer
Матрица оператора: $A = \begin{pmatrix} 1 & 1 & 1 \\ 0 & 1 & 2 \\ 0 & 0 & 1 \end{pmatrix}$. Оператор обратим.

%=============================================================
\task{38}
%=============================================================
\subsection*{Условие}
Оператор $\mathcal{A} : \R^3 \to \R^3$ задан как ортогональное проектирование на прямую $L$, проходящую через начало координат и точку $(1, 1, 1)$. Найти матрицу этого оператора в стандартном базисе и вычислить её ранг.

\begin{theory}
Пусть прямая задана направляющим вектором $u = (1, 1, 1)^T$. Ортогональная проекция произвольного вектора $x$ на прямую $L$ вычисляется по формуле:
\[
\mathcal{A}(x) = \frac{\langle x, u\rangle}{\|u\|^2} u.
\]
Матрицу оператора в стандартном базисе можно найти, подавая на вход базисные векторы $e_1 = (1, 0, 0)^T$, $e_2 = (0, 1, 0)^T$, $e_3 = (0, 0, 1)^T$.
\end{theory}

\step{1} Вычислим квадрат длины направляющего вектора: $\|u\|^2 = 1^2 + 1^2 + 1^2 = 3$.

\step{2} Найдём образы базисных векторов:
\[
\mathcal{A}(e_1) = \frac{\langle e_1, u\rangle}{3} u = \frac{1}{3} \begin{pmatrix} 1 \\ 1 \\ 1 \end{pmatrix} = \begin{pmatrix} 1/3 \\ 1/3 \\ 1/3 \end{pmatrix},
\]
\[
\mathcal{A}(e_2) = \frac{\langle e_2, u\rangle}{3} u = \frac{1}{3} \begin{pmatrix} 1 \\ 1 \\ 1 \end{pmatrix} = \begin{pmatrix} 1/3 \\ 1/3 \\ 1/3 \end{pmatrix},
\]
\[
\mathcal{A}(e_3) = \frac{\langle e_3, u\rangle}{3} u = \frac{1}{3} \begin{pmatrix} 1 \\ 1 \\ 1 \end{pmatrix} = \begin{pmatrix} 1/3 \\ 1/3 \\ 1/3 \end{pmatrix}.
\]

\step{3} Составим матрицу $A$:
\[
A = \frac{1}{3} \begin{pmatrix} 1 & 1 & 1 \\ 1 & 1 & 1 \\ 1 & 1 & 1 \end{pmatrix}.
\]

\step{4} Вычислим ранг матрицы $A$. Так как все строки (и столбцы) пропорциональны (равны), максимальное число линейно независимых строк равно 1. Следовательно, $\mathrm{rg}\,\mathcal{A} = 1$. (Геометрически это очевидно, так как образ оператора — одномерная прямая).

\answer
Матрица оператора: $A = \dfrac{1}{3} \begin{pmatrix} 1 & 1 & 1 \\ 1 & 1 & 1 \\ 1 & 1 & 1 \end{pmatrix}$. Ранг оператора: $\mathrm{rg}\,\mathcal{A} = 1$.

%=============================================================
\task{39}
%=============================================================
\subsection*{Условие}
Доказать, что если $\mathcal{A}$ — линейный оператор в $n$-мерном пространстве, такой что $\mathrm{im}\,\mathcal{A} \subseteq \mathrm{ker}\,\mathcal{A}$, то $\mathrm{rg}\,\mathcal{A} \leq n/2$. Привести пример такого оператора в $\R^2$.

\begin{theory}
Размерности ядра и образа любого линейного оператора в $n$-мерном пространстве связаны теоремой о сумме размерностей:
\[
\mathrm{rg}\,\mathcal{A} + \mathrm{def}\,\mathcal{A} = n, \qquad \text{где } \mathrm{rg}\,\mathcal{A} = \dim \mathrm{im}\,\mathcal{A}, \ \mathrm{def}\,\mathcal{A} = \dim \mathrm{ker}\,\mathcal{A}.
\]
Для подпространств верно свойство: если $U \subseteq W$, то $\dim U \leq \dim W$.
\end{theory}

\step{1} Из условия $\mathrm{im}\,\mathcal{A} \subseteq \mathrm{ker}\,\mathcal{A}$ получаем неравенство для их размерностей:
\[
\dim \mathrm{im}\,\mathcal{A} \leq \dim \mathrm{ker}\,\mathcal{A} \implies \mathrm{rg}\,\mathcal{A} \leq \mathrm{def}\,\mathcal{A}.
\]

\step{2} Подставим это соотношение в уравнение теоремы о сумме размерностей:
\[
n = \mathrm{rg}\,\mathcal{A} + \mathrm{def}\,\mathcal{A} \geq \mathrm{rg}\,\mathcal{A} + \mathrm{rg}\,\mathcal{A} = 2\,\mathrm{rg}\,\mathcal{A}.
\]
Отсюда напрямую следует требуемое неравенство:
\[
2\,\mathrm{rg}\,\mathcal{A} \leq n \implies \mathrm{rg}\,\mathcal{A} \leq \frac{n}{2}. \quad\blacksquare
\]

\step{3} \textbf{Пример в $\R^2$ ($n = 2$):}
Нам необходим неструктурированный линейный оператор с $\mathrm{rg}\,\mathcal{A} \leq 1$. Зададим его матрицей $A$:
\[
A = \begin{pmatrix} 0 & 1 \\ 0 & 0 \end{pmatrix}.
\]
Образ оператора: $\mathrm{im}\,\mathcal{A} = \mathrm{span}\{(1, 0)^T\}$, так как $A \begin{pmatrix} x_1 \\ x_2 \end{pmatrix} = \begin{pmatrix} x_2 \\ 0 \end{pmatrix}$. \\
Ядро оператора: $\mathrm{ker}\,\mathcal{A} = \left\{ \begin{pmatrix} x_1 \\ x_2 \end{pmatrix} : x_2 = 0 \right\} = \mathrm{span}\{(1, 0)^T\}$. \\
Здесь $\mathrm{im}\,\mathcal{A} = \mathrm{ker}\,\mathcal{A}$ (значит, выполнено включение $\mathrm{im}\,\mathcal{A} \subseteq \mathrm{ker}\,\mathcal{A}$), а ранг $\mathrm{rg}\,\mathcal{A} = 1 = 2/2$.

\answer
Утверждение доказано. Пример оператора в $\R^2$: $\mathcal{A}$ с матрицей $A = \begin{pmatrix} 0 & 1 \\ 0 & 0 \end{pmatrix}$.

%=============================================================
\task{40}
%=============================================================
\subsection*{Условие}
Найти базис ядра оператора $\mathcal{A} : \R^4 \to \R^3$, заданного матрицей:
\[
A = \begin{pmatrix} 1 & -2 & 3 & 1 \\ 2 & 1 & 1 & -3 \\ 1 & 8 & -7 & -9 \end{pmatrix}
\]
Какова размерность образа этого оператора?

\begin{theory}
Ядро оператора — это пространство решений однородной системы линейных уравнений $A x = 0$. Базис ядра представляет собой фундаментальную систему решений (ФСР) этой системы. Размерность образа равна рангу матрицы: $\dim \mathrm{im}\,\mathcal{A} = \mathrm{rg}(A) = 4 - \dim \mathrm{ker}\,\mathcal{A}$.
\end{theory}

\step{1} Приведём матрицу $A$ к ступенчатому виду с помощью метода Гаусса.
Вычтем из второй строки первую, умноженную на 2, а из третьей строки — первую строку:
\[
\begin{pmatrix} 1 & -2 & 3 & 1 \\ 0 & 5 & -5 & -5 \\ 0 & 10 & -10 & -10 \end{pmatrix}.
\]
Разделим вторую строку на 5:
\[
\begin{pmatrix} 1 & -2 & 3 & 1 \\ 0 & 1 & -1 & -1 \\ 0 & 10 & -10 & -10 \end{pmatrix}.
\]
Вычтем из третьей строки вторую, умноженную на 10:
\[
\begin{pmatrix} 1 & -2 & 3 & 1 \\ 0 & 1 & -1 & -1 \\ 0 & 0 & 0 & 0 \end{pmatrix}.
\]

\step{2} Ранг матрицы (число ненулевых строк) равен $\mathrm{rg}(A) = 2$.
Следовательно, размерность образа:
\[
\dim \mathrm{im}\,\mathcal{A} = \mathrm{rg}(A) = 2.
\]

\step{3} Запишем эквивалентную систему уравнений для нахождения ФСР:
\[
\begin{cases}
x_1 - 2x_2 + 3x_3 + x_4 = 0, \\
x_2 - x_3 - x_4 = 0.
\end{cases}
\]
Базисными переменными являются $x_1, x_2$, свободными — $x_3, x_4$. Выразим базисные переменные через свободные:
\[
x_2 = x_3 + x_4,
\]
\[
x_1 = 2x_2 - 3x_3 - x_4 = 2(x_3 + x_4) - 3x_3 - x_4 = -x_3 + x_4.
\]

\step{4} Составим фундаментальную систему решений (ФСР), поочередно полагая свободные переменные равными $(1, 0)$ и $(0, 1)$:
\begin{itemize}
    \item При $x_3 = 1, x_4 = 0$: \ $u_1 = (-1, 1, 1, 0)^T$.
    \item При $x_3 = 0, x_4 = 1$: \ $u_2 = (1, 1, 0, 1)^T$.
\end{itemize}
Векторы $\{u_1, u_2\}$ образуют базис ядра оператора.

\answer
Базис ядра оператора: $u_1 = \begin{pmatrix} -1 \\ 1 \\ 1 \\ 0 \end{pmatrix}$, $u_2 = \begin{pmatrix} 1 \\ 1 \\ 0 \\ 1 \end{pmatrix}$.\\
Размерность образа: $\dim \mathrm{im}\,\mathcal{A} = 2$.

%=============================================================
\task{41}
%=============================================================
\subsection*{Условие}
Найти собственные значения и собственные векторы линейного оператора, заданного в стандартном базисе матрицей:
\[
A = \begin{pmatrix} 4 & -1 & -2 \\ 1 & 1 & -1 \\ 2 & -1 & 0 \end{pmatrix}
\]
Указать кратность каждого собственного значения.

\begin{theory}
Собственные значения $\lambda$ находятся как корни характеристического уравнения:
\[
\det(A - \lambda I) = 0.
\]
Кратность корня в характеристическом многочлене называется его \textbf{алгебраической кратностью}. Собственный вектор $x \neq 0$, отвечающий значению $\lambda$, находится из однородной системы линейных уравнений:
\[
(A - \lambda I)x = 0.
\]
\end{theory}

\step{1} Составим характеристический многочлен $\chi(\lambda) = \det(A - \lambda I)$:
\[
\det \begin{pmatrix} 4-\lambda & -1 & -2 \\ 1 & 1-\lambda & -1 \\ 2 & -1 & -\lambda \end{pmatrix} = 0.
\]
Раскроем определитель по первой строке:
\[
(4-\lambda) \Big( (1-\lambda)(-\lambda) - 1 \Big) - (-1) \Big( 1(-\lambda) - 2(-1) \Big) + (-2) \Big( 1(-1) - 2(1-\lambda) \Big) = 0,
\]
\[
(4-\lambda)(\lambda^2 - \lambda - 1) + (2 - \lambda) - 2(2\lambda - 3) = 0,
\]
\[
(4\lambda^2 - 4\lambda - 4 - \lambda^3 + \lambda^2 + \lambda) + (2 - \lambda) - (4\lambda - 6) = 0,
\]
\[
-\lambda^3 + 5\lambda^2 - 8\lambda + 4 = 0.
\]
Перепишем уравнение в виде: $\lambda^3 - 5\lambda^2 + 8\lambda - 4 = 0$.

\step{2} Найдём корни этого уравнения. Легко заметить, что $\lambda = 1$ является корнем: $1 - 5 + 8 - 4 = 0$. Разделим многочлен на $(\lambda - 1)$:
\[
\lambda^3 - 5\lambda^2 + 8\lambda - 4 = (\lambda - 1)(\lambda^2 - 4\lambda + 4) = (\lambda - 1)(\lambda - 2)^2 = 0.
\]
Следовательно, собственные значения:
\begin{itemize}
    \item $\lambda_1 = 1$ с алгебраической кратностью $k_1 = 1$;
    \item $\lambda_2 = 2$ с алгебраической кратностью $k_2 = 2$.
\end{itemize}

\step{3} Найдём собственный вектор для $\lambda_1 = 1$. Решим систему $(A - I)x = 0$:
\[
A - I = \begin{pmatrix} 3 & -1 & -2 \\ 1 & 0 & -1 \\ 2 & -1 & -1 \end{pmatrix} \sim \begin{pmatrix} 1 & 0 & -1 \\ 0 & -1 & 1 \\ 0 & 0 & 0 \end{pmatrix}.
\]
Отсюда $x_1 = x_3$, $x_2 = x_3$. Положим $x_3 = 1$. Собственный вектор: $v_1 = (1, 1, 1)^T$.

\step{4} Найдём собственный вектор для $\lambda_2 = 2$. Решим систему $(A - 2I)x = 0$:
\[
A - 2I = \begin{pmatrix} 2 & -1 & -2 \\ 1 & -1 & -1 \\ 2 & -1 & -2 \end{pmatrix} \sim \begin{pmatrix} 1 & -1 & -1 \\ 0 & 1 & 0 \\ 0 & 0 & 0 \end{pmatrix}.
\]
Отсюда $x_2 = 0$, $x_1 = x_3$. Положим $x_3 = 1$. Собственный вектор: $v_2 = (1, 0, 1)^T$.

\answer
Собственные значения: \\
$\lambda_1 = 1$ (алгебраическая кратность 1), собственный вектор $v_1 = (1, 1, 1)^T$; \\
$\lambda_2 = 2$ (алгебраическая кратность 2), собственный вектор $v_2 = (1, 0, 1)^T$ (геометрическая кратность равна 1).

%=============================================================
\task{42}
%=============================================================
\subsection*{Условие}
Найти базис собственного подпространства $V_\lambda$, отвечающего собственному значению $\lambda = 2$ для матрицы:
\[
B = \begin{pmatrix} 3 & -1 & 1 \\ -1 & 3 & -1 \\ 1 & -1 & 3 \end{pmatrix}
\]
Какова геометрическая кратность этого собственного значения?

\begin{theory}
Собственное подпространство $V_\lambda$ — это ядро оператора $(B - \lambda I)$, то есть пространство решений однородной системы $(B - \lambda I)x = 0$. \textbf{Геометрическая кратность} собственного значения — это размерность собственного подпространства, то есть $\dim V_\lambda = n - \mathrm{rg}(B - \lambda I)$.
\end{theory}

\step{1} Составим систему уравнений для $\lambda = 2$, записав матрицу $B - 2I$:
\[
B - 2I = \begin{pmatrix} 1 & -1 & 1 \\ -1 & 1 & -1 \\ 1 & -1 & 1 \end{pmatrix}.
\]

\step{2} Приведём её к ступенчатому виду. Очевидно, что вторая строка равна первой, умноженной на $-1$, а третья строка просто равна первой. Система сводится к одному линейному уравнению:
\[
x_1 - x_2 + x_3 = 0 \implies x_1 = x_2 - x_3.
\]

\step{3} Свободными переменными являются $x_2$ и $x_3$. Найдём ФСР:
\begin{itemize}
    \item При $x_2 = 1, x_3 = 0$ получаем вектор: $u_1 = (1, 1, 0)^T$;
    \item При $x_2 = 0, x_3 = 1$ получаем вектор: $u_2 = (-1, 0, 1)^T$.
\end{itemize}
Система векторов $\{u_1, u_2\}$ образует базис собственного подпространства $V_2$.

\step{4} Геометрическая кратность собственного значения $\lambda = 2$ равна размерности построенного подпространства:
\[
g = \dim V_2 = 2.
\]

\answer
Базис собственного подпространства $V_2$: $u_1 = \begin{pmatrix} 1 \\ 1 \\ 0 \end{pmatrix}$, $u_2 = \begin{pmatrix} -1 \\ 0 \\ 1 \end{pmatrix}$. Геометрическая кратность собственного значения $\lambda = 2$ равна 2.

%=============================================================
\task{43}
%=============================================================
\subsection*{Условие}
Выяснить, можно ли привести матрицу оператора к диагональному виду путем перехода к новому базису:
\[
C = \begin{pmatrix} 1 & 1 & 0 \\ 0 & 1 & 0 \\ 0 & 0 & 2 \end{pmatrix}
\]
Ответ обосновать, проверив совпадение алгебраической и геометрической кратностей собственных чисел.

\begin{theory}
Матрица оператора диагонализуема тогда и только тогда, когда для каждого её собственного значения $\lambda_i$ геометрическая кратность $g_i$ равна алгебраической кратности $a_i$. 
Для треугольной или блочно-треугольной матрицы собственные значения стоят на главной диагонали.
\end{theory}

\step{1} Выпишем собственные значения матрицы $C$, которая является верхнетреугольной:
\begin{itemize}
    \item $\lambda_1 = 1$ с алгебраической кратностью $a_1 = 2$;
    \item $\lambda_2 = 2$ с алгебраической кратностью $a_2 = 1$.
\end{itemize}

\step{2} Для $\lambda_2 = 2$ (кратности 1) геометрическая кратность заведомо равна алгебраической: $g_2 = a_2 = 1$.

\step{3} Вычислим геометрическую кратность для $\lambda_1 = 1$. Найдём ранг матрицы $C - I$:
\[
C - I = \begin{pmatrix} 0 & 1 & 0 \\ 0 & 0 & 0 \\ 0 & 0 & 1 \end{pmatrix}.
\]
В этой матрице две строки линейно независимы (первая и третья), значит $\mathrm{rg}(C - I) = 2$. \\
Тогда геометрическая кратность $g_1$ равна:
\[
g_1 = 3 - \mathrm{rg}(C - I) = 3 - 2 = 1.
\]

\step{4} Сравним кратности для собственного числа $\lambda_1 = 1$:
\[
g_1 = 1 < a_1 = 2.
\]
Поскольку геометрическая кратность строго меньше алгебраической, построить базис из собственных векторов невозможно.

\answer
Матрицу $C$ \textbf{нельзя} привести к диагональному виду, так как для собственного числа $\lambda_1 = 1$ геометрическая кратность ($g_1 = 1$) не совпадает с его алгебраической кратностью ($a_1 = 2$).

%=============================================================
\task{44}
%=============================================================
\subsection*{Условие}
Доказать, что если $\lambda$ является собственным значением невырожденного оператора $\mathcal{A}$, то $1/\lambda$ является собственным значением оператора $\mathcal{A}^{-1}$.

\begin{theory}
Вектор $x \neq 0$ называется собственным вектором оператора $\mathcal{A}$, отвечающим собственному значению $\lambda$, если выполнено равенство:
\[
\mathcal{A}x = \lambda x.
\]
Невырожденный оператор $\mathcal{A}$ обладает обратным оператором $\mathcal{A}^{-1}$, при этом все его собственные значения отличны от нуля ($\lambda \neq 0$).
\end{theory}

\step{1} Запишем определяющее уравнение для собственного значения $\lambda$ и соответствующего собственного вектора $x \neq 0$:
\[
\mathcal{A}x = \lambda x.
\]

\step{2} Поскольку оператор $\mathcal{A}$ невырожден, применим обратный оператор $\mathcal{A}^{-1}$ к обеим частям равенства:
\[
\mathcal{A}^{-1}(\mathcal{A}x) = \mathcal{A}^{-1}(\lambda x).
\]

\step{3} По определению обратного оператора левая часть равна $x$. В правой части, используя линейность (однородность) оператора $\mathcal{A}^{-1}$, вынесем скаляр $\lambda$:
\[
x = \lambda \mathcal{A}^{-1}x.
\]

\step{4} Так как оператор невырожден, $\lambda \neq 0$. Разделим обе части полученного равенства на $\lambda$:
\[
\mathcal{A}^{-1}x = \frac{1}{\lambda} x.
\]
Это равенство означает, что ненулевой вектор $x$ является собственным вектором оператора $\mathcal{A}^{-1}$, отвечающим собственному значению $1/\lambda$. \quad$\blacksquare$

\answer
Утверждение доказано: из $\mathcal{A}x = \lambda x$ следует $\mathcal{A}^{-1}x = \dfrac{1}{\lambda}x$.

%=============================================================
\task{45}
%=============================================================
\subsection*{Условие}
Пусть $\mathcal{P}$ — оператор проектирования ($\mathcal{P}^2 = \mathcal{P}$). Доказать, что его собственными значениями могут быть только числа 0 и 1. Описать соответствующие собственные подпространства.

\begin{theory}
Собственный вектор $x \neq 0$ удовлетворяет уравнению $\mathcal{P}x = \lambda x$. Подпространство, отвечающее собственному числу $\lambda$, состоит из всех векторов, удовлетворяющих этому уравнению.
\end{theory}

\step{1} Запишем уравнение на собственные значения для оператора $\mathcal{P}$:
\[
\mathcal{P}x = \lambda x \quad (x \neq 0).
\]

\step{2} Применим оператор $\mathcal{P}$ к обеим частям равенства ещё раз:
\[
\mathcal{P}^2 x = \mathcal{P}(\lambda x) = \lambda \mathcal{P}x.
\]

\step{3} Используя тождество $\mathcal{P}^2 = \mathcal{P}$ для левой части и подставляя $\mathcal{P}x = \lambda x$ в правую часть, получаем:
\[
\mathcal{P}x = \lambda (\lambda x) \implies \lambda x = \lambda^2 x \implies (\lambda^2 - \lambda)x = 0.
\]

\step{4} Поскольку собственный вектор по определению ненулевой ($x \neq 0$), должно выполняться числовое равенство:
\[
\lambda^2 - \lambda = 0 \implies \lambda(\lambda - 1) = 0 \implies \lambda \in \{0, 1\}.
\]
Таким образом, другими собственные значения быть не могут.

\step{5} Опишем структуру собственных подпространств:
\begin{itemize}
    \item Собственное подпространство $V_0$, отвечающее $\lambda = 0$:
    \[
    V_0 = \{x : \mathcal{P}x = 0\} = \mathrm{ker}\,\mathcal{P}.
    \]
    Это ядро проектора — множество векторов, проецирующихся в нуль.
    \item Собственное подпространство $V_1$, отвечающее $\lambda = 1$:
    \[
    V_1 = \{x : \mathcal{P}x = x\} = \mathrm{im}\,\mathcal{P}.
    \]
    Это образ проектора — множество векторов, остающихся неизменными при проектировании (само подпространство проектирования).
\end{itemize}

\answer
Доказано, что $\lambda \in \{0, 1\}$. Собственное подпространство $V_0 = \mathrm{ker}\,\mathcal{P}$ (ядро проектора), а собственное подпространство $V_1 = \mathrm{im}\,\mathcal{P}$ (подпространство, на которое проектируют).

%=============================================================
\task{46}
%=============================================================
\subsection*{Условие}
Для матрицы $A = \begin{pmatrix} 1 & 2 \\ 3 & 4 \end{pmatrix}$ составить характеристический многочлен $\chi(\lambda)$ и прямой подстановкой убедиться, что $\chi(A) = 0$.

\begin{theory}
\textbf{Теорема Гамильтона — Кэли:} любая квадратная матрица $A$ удовлетворяет своему характеристическому уравнению, то есть если $\chi(\lambda) = \det(A - \lambda I) = \lambda^n - c_1 \lambda^{n-1} + \dots + (-1)^n c_n$, то $\chi(A) = A^n - c_1 A^{n-1} + \dots + (-1)^n c_n I = 0$.
Для матрицы порядка $2\times2$ характеристический многочлен имеет вид:
\[
\chi(\lambda) = \lambda^2 - \tr(A)\lambda + \det(A).
\]
\end{theory}

\step{1} Найдём след и определитель матрицы $A$:
\[
\tr(A) = 1 + 4 = 5,
\]
\[
\det(A) = 1\cdot 4 - 2\cdot 3 = 4 - 6 = -2.
\]
Составим характеристический многочлен:
\[
\chi(\lambda) = \lambda^2 - 5\lambda - 2.
\]

\step{2} Для подстановки матрицы вычислим её квадрат $A^2$:
\[
A^2 = \begin{pmatrix} 1 & 2 \\ 3 & 4 \end{pmatrix} \begin{pmatrix} 1 & 2 \\ 3 & 4 \end{pmatrix} = \begin{pmatrix} 1\cdot1 + 2\cdot3 & 1\cdot2 + 2\cdot4 \\ 3\cdot1 + 4\cdot3 & 3\cdot2 + 4\cdot4 \end{pmatrix} = \begin{pmatrix} 7 & 10 \\ 15 & 22 \end{pmatrix}.
\]

\step{3} Выполним прямую подстановку в многочлен $\chi(A) = A^2 - 5A - 2I$:
\[
\chi(A) = \begin{pmatrix} 7 & 10 \\ 15 & 22 \end{pmatrix} - 5 \begin{pmatrix} 1 & 2 \\ 3 & 4 \end{pmatrix} - 2 \begin{pmatrix} 1 & 0 \\ 0 & 1 \end{pmatrix}
\]
\[
= \begin{pmatrix} 7 & 10 \\ 15 & 22 \end{pmatrix} - \begin{pmatrix} 5 & 10 \\ 15 & 20 \end{pmatrix} - \begin{pmatrix} 2 & 0 \\ 0 & 2 \end{pmatrix}
\]
\[
= \begin{pmatrix} 7 - 5 - 2 & 10 - 10 - 0 \\ 15 - 15 - 0 & 22 - 20 - 2 \end{pmatrix} = \begin{pmatrix} 0 & 0 \\ 0 & 0 \end{pmatrix}. \quad\checkmark
\]

\answer
Характеристический многочлен: $\chi(\lambda) = \lambda^2 - 5\lambda - 2$. \\
Прямой подстановкой подтверждено, что $\chi(A) = A^2 - 5A - 2I = \begin{pmatrix} 0 & 0 \\ 0 & 0 \end{pmatrix}$.

%=============================================================
\task{47}
%=============================================================
\subsection*{Условие}
Доказать, что для любого линейного оператора $\mathcal{A}$ в евклидовом пространстве выполняется равенство $(\mathrm{Im}\,\mathcal{A})^\perp = \mathrm{Ker}\,\mathcal{A}^*$. Дать геометрическую интерпретацию этого факта.

\begin{theory}
По определению сопряжённого оператора $\mathcal{A}^*$, для любых векторов $x, y$ выполняется равенство:
\[
\langle \mathcal{A}x, y\rangle = \langle x, \mathcal{A}^* y\rangle.
\]
Подпространство $U^\perp$ состоит из векторов, ортогональных всем векторам из $U$.
\end{theory}

\step{1} Докажем включение $(\mathrm{Im}\,\mathcal{A})^\perp \subseteq \mathrm{Ker}\,\mathcal{A}^*$. \\
Пусть вектор $y \in (\mathrm{Im}\,\mathcal{A})^\perp$. Это означает, что для любого вектора $u \in \mathrm{Im}\,\mathcal{A}$ выполнено $\langle u, y\rangle = 0$. Любой вектор из образа можно записать в виде $u = \mathcal{A}x$, следовательно:
\[
\langle \mathcal{A}x, y\rangle = 0 \quad \text{для всех } x.
\]
Используя свойство сопряжённого оператора, перепишем равенство:
\[
\langle x, \mathcal{A}^* y\rangle = 0 \quad \text{для всех } x.
\]
Поскольку вектор $\mathcal{A}^* y$ ортогонален абсолютно любому вектору пространства $x$, он обязан быть нулевым:
\[
\mathcal{A}^* y = 0 \implies y \in \mathrm{Ker}\,\mathcal{A}^*.
\]

\step{2} Докажем обратное включение $\mathrm{Ker}\,\mathcal{A}^* \subseteq (\mathrm{Im}\,\mathcal{A})^\perp$. \\
Пусть вектор $y \in \mathrm{Ker}\,\mathcal{A}^*$, то есть $\mathcal{A}^* y = 0$. 
Для любого вектора образа $u = \mathcal{A}x$ рассмотрим скалярное произведение:
\[
\langle u, y\rangle = \langle \mathcal{A}x, y\rangle = \langle x, \mathcal{A}^* y\rangle = \langle x, 0\rangle = 0.
\]
Следовательно, вектор $y$ ортогонален любому элементу из образа, откуда $y \in (\mathrm{Im}\,\mathcal{A})^\perp$.

\step{3} Объединяя шаги 1 и 2, получаем тождество:
\[
(\mathrm{Im}\,\mathcal{A})^\perp = \mathrm{Ker}\,\mathcal{A}^*. \quad\blacksquare
\]

\step{4} \textbf{Геометрическая интерпретация.} Данное соотношение показывает двойственность между областью значений оператора и его ядром. Ортогональное дополнение к пространству всех достижимых состояний (образу) оператора — это в точности те направления, которые переводятся сопряжённым оператором в ноль.

\answer
Равенство $(\mathrm{Im}\,\mathcal{A})^\perp = \mathrm{Ker}\,\mathcal{A}^*$ строго доказано. Геометрически это означает, что направления, ортогональные образу исходного оператора, образуют ядро сопряжённого оператора.

%=============================================================
\task{48}
%=============================================================
\subsection*{Условие}
Пусть в базисе $\{e_i\}$ оператор $\mathcal{A}$ имеет матрицу $A$, а метрика пространства задана матрицей Грама $G$. Доказать, что матрица сопряженного оператора $A^*$ вычисляется по формуле $A^* = G^{-1}A^TG$.

\begin{theory}
В базисе $\{e_i\}$ скалярное произведение векторов с векторами-столбцами координат $X$ и $Y$ вычисляется как:
\[
\langle x, y\rangle = X^T G Y.
\]
Сопряжённый оператор $\mathcal{A}^*$ по определению удовлетворяет равенству $\langle \mathcal{A}x, y\rangle = \langle x, \mathcal{A}^* y\rangle$ для всех векторов.
\end{theory}

\step{1} Пусть вектор $x$ имеет координаты $X$, а вектор $y$ — координаты $Y$. Образ $\mathcal{A}x$ имеет координаты $A X$, а образ $\mathcal{A}^* y$ имеет координаты $A^* Y$.

\step{2} Запишем скалярное произведение $\langle \mathcal{A}x, y\rangle$ в матричном виде:
\[
\langle \mathcal{A}x, y\rangle = (A X)^T G Y = X^T A^T G Y.
\]

\step{3} Запишем скалярное произведение $\langle x, \mathcal{A}^* y\rangle$ аналогичным образом:
\[
\langle x, \mathcal{A}^* y\rangle = X^T G (A^* Y).
\]

\step{4} Приравняем полученные выражения, так как равенство должно выполняться для любых координатных столбцов $X$ и $Y$:
\[
X^T A^T G Y = X^T G A^* Y \quad \forall X, Y \implies A^T G = G A^*.
\]

\step{5} Поскольку матрица Грама $G$ положительно определена, она невырождена. Домножим полученное соотношение слева на $G^{-1}$:
\[
G^{-1} A^T G = A^* \iff A^* = G^{-1} A^T G. \quad\blacksquare
\]

\answer
Утверждение доказано. Формула $A^* = G^{-1} A^T G$ является обобщением сопряжения для неортонормированных базисов.

%=============================================================
\task{49}
%=============================================================
\subsection*{Условие}
Пусть $L$ — инвариантное подпространство самосопряженного оператора $\mathcal{A}$. Доказать, что его ортогональное дополнение $L^\perp$ также является инвариантным подпространством для $\mathcal{A}$.

\begin{theory}
\begin{itemize}
    \item Подпространство $L$ инвариантно относительно $\mathcal{A}$, если для любого $x \in L$ выполняется $\mathcal{A}x \in L$.
    \item Оператор $\mathcal{A}$ самосопряжён, если $\langle \mathcal{A}x, y\rangle = \langle x, \mathcal{A}y\rangle$ для всех $x, y$.
\end{itemize}
\end{theory}

\step{1} Нам необходимо доказать, что для любого $y \in L^\perp$ вектор $\mathcal{A}y$ также принадлежит $L^\perp$. По определению ортогонального дополнения, это означает, что:
\[
\langle x, \mathcal{A}y\rangle = 0 \quad \text{для всех } x \in L.
\]

\step{2} Воспользуемся самосопряжённостью оператора $\mathcal{A}$ для переноса действия оператора на первый аргумент:
\[
\langle x, \mathcal{A}y\rangle = \langle \mathcal{A}x, y\rangle.
\]

\step{3} Поскольку $L$ инвариантно относительно $\mathcal{A}$ и $x \in L$, образ этого вектора также принадлежит $L$:
\[
\mathcal{A}x = u \in L.
\]

\step{4} Из условия $y \in L^\perp$ следует, что вектор $y$ ортогонален любому вектору из подпространства $L$. Так как $u \in L$, имеем:
\[
\langle \mathcal{A}x, y\rangle = \langle u, y\rangle = 0.
\]

\step{5} Объединяя результаты шагов 2 и 4, получаем:
\[
\langle x, \mathcal{A}y\rangle = 0 \quad \text{для всех } x \in L \implies \mathcal{A}y \in L^\perp.
\]
Это доказывает инвариантность подпространства $L^\perp$. \quad$\blacksquare$

\answer
Утверждение доказано: инвариантность подпространства $L$ для самосопряжённого оператора автоматически влечёт инвариантность его ортогонального дополнения $L^\perp$.

%=============================================================
\task{50}
%=============================================================
\subsection*{Условие}
Пусть $\lambda_{\min}$ и $\lambda_{\max}$ — минимальное и максимальное собственные значения самосопряженного оператора $\mathcal{A}$. Доказать, что для любого единичного вектора $\|x\| = 1$ выполняется неравенство: $\lambda_{\min} \leq \langle \mathcal{A}x, x\rangle \leq \lambda_{\max}$.

\begin{theory}
В силу спектральной теоремы, для любого самосопряжённого оператора в конечномерном пространстве существует ортонормированный базис из собственных векторов $\{e_1, \ldots, e_n\}$, отвечающих вещественным собственным значениям $\lambda_1, \ldots, \lambda_n$.
\end{theory}

\step{1} Пусть собственные значения упорядочены по возрастанию:
\[
\lambda_{\min} = \lambda_1 \leq \lambda_2 \leq \dots \leq \lambda_n = \lambda_{\max}.
\]
Разложим произвольный единичный вектор $x$ по ортонормированному базису собственных векторов:
\[
x = \sum_{i=1}^n c_i e_i.
\]

\step{2} Так как вектор $x$ единичный ($\|x\| = 1$), сумма квадратов координат равна 1:
\[
\|x\|^2 = \sum_{i=1}^n |c_i|^2 = 1.
\]

\step{3} Вычислим действие оператора на вектор $x$:
\[
\mathcal{A}x = \sum_{i=1}^n c_i \mathcal{A}e_i = \sum_{i=1}^n c_i \lambda_i e_i.
\]

\step{4} Найдём скалярное произведение $\langle \mathcal{A}x, x\rangle$. В силу ортонормированности базиса $\{e_i\}$:
\[
\langle \mathcal{A}x, x\rangle = \sum_{i=1}^n \lambda_i |c_i|^2.
\]

\step{5} Оценим сумму снизу, заменив каждое собственное значение $\lambda_i$ на минимальное $\lambda_{\min}$:
\[
\langle \mathcal{A}x, x\rangle = \sum_{i=1}^n \lambda_i |c_i|^2 \geq \sum_{i=1}^n \lambda_{\min} |c_i|^2 = \lambda_{\min} \sum_{i=1}^n |c_i|^2 = \lambda_{\min} \cdot 1 = \lambda_{\min}.
\]

\step{6} Аналогично оценим сумму сверху, заменив собственные значения на максимальное $\lambda_{\max}$:
\[
\langle \mathcal{A}x, x\rangle = \sum_{i=1}^n \lambda_i |c_i|^2 \leq \sum_{i=1}^n \lambda_{\max} |c_i|^2 = \lambda_{\max} \sum_{i=1}^n |c_i|^2 = \lambda_{\max} \cdot 1 = \lambda_{\max}.
\]
Объединяя оценки, получаем требуемое неравенство: $\lambda_{\min} \leq \langle \mathcal{A}x, x\rangle \leq \lambda_{\max}$. \quad$\blacksquare$

\answer
Двойное неравенство (известное как отношение Рэлея) строго доказано с помощью разложения по ортонормированному собственному базису.

%=============================================================
\task{51}
%=============================================================
\subsection*{Условие}
Доказать, что оператор проектирования $\mathcal{P}$ ($\mathcal{P}^2 = \mathcal{P}$) является самосопряженным тогда и только тогда, когда это оператор ортогонального проектирования.

\begin{theory}
\begin{itemize}
    \item Оператор $\mathcal{P}$ является самосопряжённым, если $\mathcal{P}^* = \mathcal{P}$, то есть $\langle \mathcal{P}x, y\rangle = \langle x, \mathcal{P}y\rangle$ для всех $x, y$.
    \item Проектор $\mathcal{P}$ является ортогональным проектором, если его образ $\mathrm{Im}\,\mathcal{P}$ и ядро $\mathrm{Ker}\,\mathcal{P}$ ортогональны: $\mathrm{Im}\,\mathcal{P} \perp \mathrm{Ker}\,\mathcal{P}$.
\end{itemize}
\end{theory}

\step{1} \textbf{Необходимость ($\Rightarrow$).} Пусть $\mathcal{P}$ — самосопряжённый проектор. 
Возьмём произвольные векторы $u \in \mathrm{Im}\,\mathcal{P}$ и $v \in \mathrm{Ker}\,\mathcal{P}$. По свойствам проектора выполнено $\mathcal{P}u = u$ и $\mathcal{P}v = 0$.
Вычислим скалярное произведение:
\[
\langle u, v\rangle = \langle \mathcal{P}u, v\rangle = \langle u, \mathcal{P}^* v\rangle = \langle u, \mathcal{P}v\rangle = \langle u, 0\rangle = 0.
\]
Так как любой вектор из образа ортогонален любому вектору из ядра, проектор $\mathcal{P}$ ортогональный.

\step{2} \textbf{Достаточность ($\Leftarrow$).} Пусть $\mathcal{P}$ — ортогональный проектор, то есть $\mathrm{Im}\,\mathcal{P} \perp \mathrm{Ker}\,\mathcal{P}$.
Любые два вектора $x, y$ пространства можно единственным образом разложить в прямую сумму:
\[
x = x_1 + x_2, \qquad y = y_1 + y_2, \qquad \text{где } x_1, y_1 \in \mathrm{Im}\,\mathcal{P}, \ \ x_2, y_2 \in \mathrm{Ker}\,\mathcal{P}.
\]
При этом $\mathcal{P}x = x_1$ и $\mathcal{P}y = y_1$.

\step{3} Вычислим левую часть равенства самосопряжённости $\langle \mathcal{P}x, y\rangle$:
\[
\langle \mathcal{P}x, y\rangle = \langle x_1, y_1 + y_2\rangle = \langle x_1, y_1\rangle + \langle x_1, y_2\rangle.
\]
Так как $x_1 \in \mathrm{Im}\,\mathcal{P}$ и $y_2 \in \mathrm{Ker}\,\mathcal{P}$, в силу ортогональности проектора их произведение равно нулю: $\langle x_1, y_2\rangle = 0$. Значит:
\[
\langle \mathcal{P}x, y\rangle = \langle x_1, y_1\rangle.
\]

\step{4} Аналогично вычислим правую часть равенства $\langle x, \mathcal{P}y\rangle$:
\[
\langle x, \mathcal{P}y\rangle = \langle x_1 + x_2, y_1\rangle = \langle x_1, y_1\rangle + \langle x_2, y_1\rangle.
\]
Поскольку $x_2 \in \mathrm{Ker}\,\mathcal{P}$ и $y_1 \in \mathrm{Im}\,\mathcal{P}$, то $\langle x_2, y_1\rangle = 0$. Получаем:
\[
\langle x, \mathcal{P}y\rangle = \langle x_1, y_1\rangle.
\]
Следовательно, $\langle \mathcal{P}x, y\rangle = \langle x, \mathcal{P}y\rangle$ для любых $x, y$, что доказывает равенство $\mathcal{P}^* = \mathcal{P}$. \quad$\blacksquare$

\answer
Утверждение доказано: самосопряжённость проектора эквивалентна ортогональности его образа и ядра.

%=============================================================
\task{52}
%=============================================================
\subsection*{Условие}
Проверить, является ли ортогональной матрица
\[
Q = \frac{1}{9} \begin{pmatrix} 8 & 1 & -4 \\ -4 & 4 & -7 \\ 1 & 8 & 4 \end{pmatrix}
\]
В случае утвердительного ответа выяснить, сохраняет ли она ориентацию пространства.

\begin{theory}
Матрица $Q$ называется ортогональной, если $Q^T Q = I$ (где $I$ — единичная матрица). Ортогональный оператор сохраняет ориентацию пространства, если определитель его матрицы равен $+1$, и меняет ориентацию, если определитель равен $-1$.
\end{theory}

\step{1} Найдём транспонированную матрицу $Q^T$:
\[
Q^T = \frac{1}{9} \begin{pmatrix} 8 & -4 & 1 \\ 1 & 4 & 8 \\ -4 & -7 & 4 \end{pmatrix}.
\]

\step{2} Вычислим произведение $Q^T Q$:
\[
Q^T Q = \frac{1}{81} \begin{pmatrix} 8 & -4 & 1 \\ 1 & 4 & 8 \\ -4 & -7 & 4 \end{pmatrix} \begin{pmatrix} 8 & 1 & -4 \\ -4 & 4 & -7 \\ 1 & 8 & 4 \end{pmatrix}.
\]
Перемножим строки первой матрицы на столбцы второй:
\begin{itemize}
    \item Элемент $(1, 1)$: $8\cdot 8 + (-4)\cdot(-4) + 1\cdot 1 = 64 + 16 + 1 = 81$;
    \item Элемент $(1, 2)$: $8\cdot 1 + (-4)\cdot 4 + 1\cdot 8 = 8 - 16 + 8 = 0$;
    \item Элемент $(1, 3)$: $8\cdot(-4) + (-4)\cdot(-7) + 1\cdot 4 = -32 + 28 + 4 = 0$;
    \item Элемент $(2, 2)$: $1\cdot 1 + 4\cdot 4 + 8\cdot 8 = 1 + 16 + 64 = 81$;
    \item Элемент $(2, 3)$: $1\cdot(-4) + 4\cdot(-7) + 8\cdot 4 = -4 - 28 + 32 = 0$;
    \item Элемент $(3, 3)$: $(-4)\cdot(-4) + (-7)\cdot(-7) + 4\cdot 4 = 16 + 49 + 16 = 81$.
\end{itemize}
Остальные недиагональные элементы равны нулю в силу симметрии. Итого:
\[
Q^T Q = \frac{1}{81} \begin{pmatrix} 81 & 0 & 0 \\ 0 & 81 & 0 \\ 0 & 0 & 81 \end{pmatrix} = \begin{pmatrix} 1 & 0 & 0 \\ 0 & 1 & 0 \\ 0 & 0 & 1 \end{pmatrix} = I.
\]
Матрица действительно ортогональная. \quad$\checkmark$

\step{3} Вычислим определитель матрицы $Q$. Вынесем общий множитель $1/9$:
\[
\det(Q) = \left(\frac{1}{9}\right)^3 \det \begin{pmatrix} 8 & 1 & -4 \\ -4 & 4 & -7 \\ 1 & 8 & 4 \end{pmatrix} = \frac{1}{729} \det \begin{pmatrix} 8 & 1 & -4 \\ -4 & 4 & -7 \\ 1 & 8 & 4 \end{pmatrix}.
\]
Вычислим определитель внутренней матрицы разложением по первой строке:
\[
\det \begin{pmatrix} 8 & 1 & -4 \\ -4 & 4 & -7 \\ 1 & 8 & 4 \end{pmatrix} = 8 \cdot (16 - (-56)) - 1 \cdot (-16 - (-7)) + (-4) \cdot (-32 - 4)
\]
\[
= 8 \cdot 72 - 1 \cdot (-9) - 4 \cdot (-36) = 576 + 9 + 144 = 729.
\]
Тогда:
\[
\det(Q) = \frac{729}{729} = 1.
\]
Так как определитель равен $+1$, оператор сохраняет ориентацию пространства.

\answer
Матрица $Q$ является ортогональной. Она сохраняет ориентацию пространства (так как $\det(Q) = 1$).

%=============================================================
\task{53}
%=============================================================
\subsection*{Условие}
Доказать, что произведение двух ортогональных операторов является ортогональным оператором. Верно ли, что сумма двух ортогональных операторов всегда ортогональна?

\begin{theory}
Оператор $\mathcal{Q}$ ортогонален тогда и только тогда, когда для его матрицы выполняется равенство $Q^* Q = I$. Ортогональные операторы сохраняют скалярное произведение и длины векторов.
\end{theory}

\step{1} Пусть $\mathcal{Q}_1$ и $\mathcal{Q}_2$ — два ортогональных оператора. Тогда для их матриц выполнено:
\[
Q_1^* Q_1 = I, \qquad Q_2^* Q_2 = I.
\]

\step{2} Рассмотрим произведение операторов $\mathcal{Q} = \mathcal{Q}_1 \mathcal{Q}_2$ (и матрицу $Q = Q_1 Q_2$). Вычислим сопряжение произведения:
\[
Q^* Q = (Q_1 Q_2)^* (Q_1 Q_2) = (Q_2^* Q_1^*) (Q_1 Q_2).
\]
Используя ассоциативность матричного умножения:
\[
Q^* Q = Q_2^* (Q_1^* Q_1) Q_2 = Q_2^* I Q_2 = Q_2^* Q_2 = I.
\]
Следовательно, произведение двух ортогональных операторов также ортогонально. \quad$\blacksquare$

\step{3} \textbf{Опровержение для суммы.} Сумма двух ортогональных операторов в общем случае \textbf{не является} ортогональным оператором. Приведём контрпример:
Пусть в одномерном пространстве $\R^1$ заданы два ортогональных оператора с матрицами $Q_1 = (1)$ и $Q_2 = (1)$ (поскольку $1\cdot1 = 1$, они ортогональны).
Их сумма равна:
\[
Q = Q_1 + Q_2 = (2).
\]
Но оператор $Q = (2)$ не является ортогональным, поскольку $2 \cdot 2 = 4 \neq 1$.

\answer
Первое утверждение доказано. Второе утверждение \textbf{неверно} (сумма двух ортогональных операторов не обязана быть ортогональной; простейший контрпример: $I + I = 2I$, который не является изометрией).

%=============================================================
\task{54}
%=============================================================
\subsection*{Условие}
Доказать, что если $L$ — инвариантное подпространство ортогонального оператора $\mathcal{Q}$, то его ортогональное дополнение $L^\perp$ также инвариантно относительно $\mathcal{Q}$.

\begin{theory}
\begin{itemize}
    \item Подпространство $L$ инвариантно относительно $\mathcal{Q}$, если $\mathcal{Q}(L) \subseteq L$.
    \item Поскольку ортогональный оператор взаимно однозначен (обратим), в конечномерном пространстве ограничение $\mathcal{Q}$ на инвариантное подпространство $L$ является биекцией, то есть $\mathcal{Q}(L) = L$.
    \item Ортогональный оператор сохраняет скалярное произведение: $\langle \mathcal{Q}u, \mathcal{Q}v\rangle = \langle u, v\rangle$.
\end{itemize}
\end{theory}

\step{1} Нам нужно показать, что для любого вектора $y \in L^\perp$ его образ $\mathcal{Q}y$ также принадлежит $L^\perp$. Это равносильно тому, что:
\[
\langle \mathcal{Q}y, x\rangle = 0 \quad \text{для любого } x \in L.
\]

\step{2} Так как подпространство $L$ инвариантно относительно обратимого оператора $\mathcal{Q}$, то любой вектор $x \in L$ может быть представлен в виде $x = \mathcal{Q}x'$ для некоторого вектора $x' \in L$.

\step{3} Запишем скалярное произведение, заменив $x$ на $\mathcal{Q}x'$:
\[
\langle \mathcal{Q}y, x\rangle = \langle \mathcal{Q}y, \mathcal{Q}x'\rangle.
\]

\step{4} В силу ортогональности оператора $\mathcal{Q}$, снимем действие оператора:
\[
\langle \mathcal{Q}y, \mathcal{Q}x'\rangle = \langle y, x'\rangle.
\]

\step{5} Поскольку $y \in L^\perp$ по условию, а вектор $x'$ принадлежит инвариантному подпространству $L$, их скалярное произведение равно нулю:
\[
\langle y, x'\rangle = 0.
\]
Следовательно, $\langle \mathcal{Q}y, x\rangle = 0$ для всех $x \in L$, то есть $\mathcal{Q}y \in L^\perp$. Инвариантность $L^\perp$ доказана. \quad$\blacksquare$

\answer
Утверждение доказано с использованием свойства обратимости ортогонального оператора на конечномерном инвариантном подпространстве и сохранения им скалярного произведения.

%=============================================================
\task{55}
%=============================================================
\subsection*{Условие}
Дана квадратичная форма в $\mathbb{R}^3$:
\[
Q(x) = x_1^2 - 4x_1 x_2 + 6x_2 x_3.
\]
Записать её симметричную матрицу $A$ и найти соответствующую ей симметричную билинейную форму $B(x, y)$.

\begin{theory}
Любая квадратичная форма $Q(x)$ однозначно порождается симметричной билинейной формой $B(x, y)$ по правилу $Q(x) = B(x, x)$. Матрица квадратичной формы $A$ составляется по правилу:
\[
a_{ii} = \text{коэффициент при } x_i^2, \qquad a_{ij} = a_{ji} = \frac{1}{2} \cdot \left(\text{коэффициент при } x_i x_j\right) \quad (i \neq j).
\]
Соответствующая симметричная билинейная форма имеет вид $B(x, y) = x^T A y$.
\end{theory}

\step{1} Найдём элементы симметричной матрицы квадратичной формы $A$:
\begin{itemize}
    \item Коэффициент при $x_1^2$ равен 1, поэтому $a_{11} = 1$;
    \item Члены $x_2^2$ и $x_3^2$ отсутствуют, поэтому $a_{22} = 0$, $a_{33} = 0$;
    \item Коэффициент при произведении $x_1 x_2$ равен $-4$, делим пополам: $a_{12} = a_{21} = -2$;
    \item Коэффициент при $x_2 x_3$ равен $6$, делим пополам: $a_{23} = a_{32} = 3$;
    \item Член с произведением $x_1 x_3$ отсутствует, поэтому $a_{13} = a_{31} = 0$.
\end{itemize}
Запишем матрицу $A$:
\[
A = \begin{pmatrix} 1 & -2 & 0 \\ -2 & 0 & 3 \\ 0 & 3 & 0 \end{pmatrix}.
\]

\step{2} Запишем симметричную билинейную форму $B(x, y) = x^T A y$:
\[
B(x, y) = \begin{pmatrix} x_1 & x_2 & x_3 \end{pmatrix} \begin{pmatrix} 1 & -2 & 0 \\ -2 & 0 & 3 \\ 0 & 3 & 0 \end{pmatrix} \begin{pmatrix} y_1 \\ y_2 \\ y_3 \end{pmatrix}
\]
\[
= x_1 y_1 - 2x_1 y_2 - 2x_2 y_1 + 3x_2 y_3 + 3x_3 y_2.
\]

\step{3} Объединим слагаемые с симметричными коэффициентами:
\[
B(x, y) = x_1 y_1 - 2(x_1 y_2 + x_2 y_1) + 3(x_2 y_3 + x_3 y_2).
\]
Легко проверить, что при подстановке $y = x$ получается исходное выражение: $B(x, x) = x_1^2 - 4x_1 x_2 + 6x_2 x_3 = Q(x)$.

\answer
Матрица квадратичной формы: $A = \begin{pmatrix} 1 & -2 & 0 \\ -2 & 0 & 3 \\ 0 & 3 & 0 \end{pmatrix}$.\\
Симметричная билинейная форма: $B(x, y) = x_1 y_1 - 2(x_1 y_2 + x_2 y_1) + 3(x_2 y_3 + x_3 y_2)$.

%=============================================================
\task{56}
%=============================================================
\subsection*{Условие}
Квадратичная форма $Q(x)$ в некотором базисе приведена к виду
\[
x_1^2 + 2x_2^2 - x_3^2.
\]
В другом базисе она же приняла вид
\[
5y_1^2 - 3y_2^2 - y_3^2.
\]
Возможно ли это? Ответ обоснуйте, вычислив сигнатуру в обоих случаях.

\begin{theory}
\textbf{Закон инерции квадратичных форм:} при любом невырожденном линейном преобразовании квадратичной формы к каноническому виду число положительных коэффициентов (положительный индекс инерции $p$) и число отрицательных коэффициентов (отрицательный индекс инерции $q$) остаются неизменными. Пара $(p, q)$ называется сигнатурой квадратичной формы и является её инвариантом.
\end{theory}

\step{1} Найдём сигнатуру квадратичной формы в первом базисе. Канонический вид:
\[
Q(x) = 1\cdot x_1^2 + 2\cdot x_2^2 - 1\cdot x_3^2.
\]
Положительные коэффициенты: $1, 2$ \ $\implies p_1 = 2$. \\
Отрицательные коэффициенты: $-1$ \ $\implies q_1 = 1$. \\
Сигнатура формы в первом базисе: $(p_1, q_1) = (2, 1)$.

\step{2} Найдём сигнатуру квадратичной формы во втором базисе. Канонический вид:
\[
Q(y) = 5\cdot y_1^2 - 3\cdot y_2^2 - 1\cdot y_3^2.
\]
Положительные коэффициенты: $5$ \ $\implies p_2 = 1$. \\
Отрицательные коэффициенты: $-3, -1$ \ $\implies q_2 = 2$. \\
Сигнатура формы во втором базисе: $(p_2, q_2) = (1, 2)$.

\step{3} Сравним полученные сигнатуры. Поскольку сигнатура является инвариантом квадратичной формы относительно любых невырожденных замен переменных, а в данном случае:
\[
(2, 1) \neq (1, 2),
\]
то одна и та же квадратичная форма не может иметь указанные канонические виды в разных базисах.

\answer
Это \textbf{невозможно}. Сигнатура первой формы равна $(2, 1)$, а второй — $(1, 2)$. В силу закона инерции квадратичных форм сигнатура является инвариантом, поэтому данные выражения описывают разные квадратичные формы.

%=============================================================
\task{57}
%=============================================================
\subsection*{Условие}
Используя критерий Сильвестра, исследовать на определенность квадратичную форму в зависимости от параметра $\lambda$:
\[
Q(x) = x_1^2 + 4x_1 x_2 + \lambda x_2^2 + 2x_3^2.
\]
При каких значениях $\lambda$ форма будет положительно определена?

\begin{theory}
\textbf{Критерий Сильвестра:} симметричная матрица $A$ (и соответствующая ей квадратичная форма) положительно определена тогда и только тогда, когда все её угловые (главные ведущие) миноры строго положительны:
\[
\Delta_1 > 0, \qquad \Delta_2 > 0, \qquad \Delta_3 > 0.
\]
\end{theory}

\step{1} Запишем симметричную матрицу квадратичной формы $A$:
\[
A = \begin{pmatrix} 1 & 2 & 0 \\ 2 & \lambda & 0 \\ 0 & 0 & 2 \end{pmatrix}.
\]

\step{2} Вычислим главные ведущие миноры матрицы $A$:
\begin{itemize}
    \item Минор первого порядка:
    \[
    \Delta_1 = a_{11} = 1 > 0. \quad\checkmark
    \]
    \item Минор второго порядка:
    \[
    \Delta_2 = \begin{vmatrix} 1 & 2 \\ 2 & \lambda \end{vmatrix} = \lambda - 4.
    \]
    Для положительной определённости должно выполняться: $\lambda - 4 > 0 \implies \lambda > 4$.
    \item Минор третьего порядка:
    \[
    \Delta_3 = \det A = 2 \cdot \Delta_2 = 2(\lambda - 4).
    \]
    Для положительной определённости должно выполняться: $2(\lambda - 4) > 0 \implies \lambda > 4$.
\end{itemize}

\step{3} Объединяя полученные условия, находим область параметров $\lambda$, при которых все три минора строго положительны: $\lambda > 4$.

\answer
Квадратичная форма будет положительно определена при $\lambda \in (4, +\infty)$.

%=============================================================
\task{58}
%=============================================================
\subsection*{Условие}
Пусть $P(t)$ обозначает популяцию некоторых бактерий на чашке Петри. Мы предполагаем, что еды достаточно и места тоже достаточно. Тогда скорость роста бактерий пропорциональна численности населения — чем больше популяция, тем выше. Составьте дифференциальное уравнение, моделирующее процесс роста популяции.
Предположим, что в чашке Петри было 100 бактерий в момент времени 0 и 200 бактерий через 10 секунд. Сколько бактерий будет через 1 минуту от времени 0 (через 60 секунд)?

\begin{theory}
Процесс неограниченного роста популяции (мальтузианская модель) описывается линейным однородным дифференциальным уравнением первого порядка:
\[
\frac{dP}{dt} = k P,
\]
где $k > 0$ — коэффициент прироста. Решением этого уравнения с начальным условием $P(0) = P_0$ является экспоненциальная функция: $P(t) = P_0 e^{kt}$.
\end{theory}

\step{1} Составим дифференциальное уравнение по условию задачи (скорость изменения популяции $dP/dt$ прямо пропорциональна численности $P$):
\[
\frac{dP}{dt} = k P \implies P(t) = P_0 e^{kt}.
\]

\step{2} Найдём константы $P_0$ и $k$ из начальных условий:
\begin{itemize}
    \item В момент $t = 0$ популяция составляла $P(0) = 100 \implies P_0 = 100$.
    \item Через $t = 10$ секунд популяция составляла $P(10) = 200$:
    \[
    100 e^{10k} = 200 \implies e^{10k} = 2 \implies 10k = \ln 2 \implies k = \frac{\ln 2}{10}.
    \]
\end{itemize}
Таким образом, закон изменения численности популяции имеет вид:
\[
P(t) = 100 \cdot e^{\frac{\ln 2}{10} t} = 100 \cdot 2^{t/10}.
\]

\step{3} Найдём численность бактерий через 1 минуту ($t = 60$ секунд):
\[
P(60) = 100 \cdot 2^{60/10} = 100 \cdot 2^6 = 100 \cdot 64 = 6400.
\]

\answer
Дифференциальное уравнение модели: $\dfrac{dP}{dt} = k P$.\\
Через 1 минуту (60 секунд) на чашке Петри будет 6400 бактерий.

%=============================================================
\task{59}
%=============================================================
\subsection*{Условие}
Пусть в резервуаре имеется $a$ кг водного раствора соли, в котором содержится $b$ кг соли. В определенный момент включается устройство, непрерывно подающее в резервуар $c$ кг чистой воды в секунду и одновременно удаляющее из него ежесекундно $c$ кг раствора. При этом в самом резервуаре жидкость непрерывно перемешивается. Как изменяется количество соли в резервуаре?

\begin{theory}
Пусть $y(t)$ — количество соли (в кг) в резервуаре в момент времени $t$. 
Скорость изменения количества соли в резервуаре равна разности скоростей её поступления и убывания:
\[
\frac{dy}{dt} = v_{\text{вх}} - v_{\text{вых}}.
\]
Поскольку подаётся чистая вода, соль не поступает ($v_{\text{вх}} = 0$). Скорость убывания соли равна произведению расхода выходящего раствора $c$ на его концентрацию $y(t)/a$.
\end{theory}

\step{1} Составим дифференциальное уравнение, описывающее скорость изменения количества соли $y(t)$:
\[
\frac{dy}{dt} = 0 - c \cdot \frac{y(t)}{a} \implies \frac{dy}{dt} = -\frac{c}{a} y.
\]

\step{2} Решим это уравнение методом разделения переменных:
\[
\frac{dy}{y} = -\frac{c}{a} dt \implies \int \frac{dy}{y} = -\frac{c}{a} \int dt \implies \ln y = -\frac{c}{a} t + C \implies y(t) = C_1 e^{-\frac{c}{a} t}.
\]

\step{3} Используем начальное условие: в момент времени $t = 0$ в резервуаре содержалось $b$ кг соли, то есть $y(0) = b$. Отсюда получаем $C_1 = b$.
Итоговый закон изменения количества соли со временем:
\[
y(t) = b e^{-\frac{c}{a} t}.
\]

\answer
Количество соли в резервуаре убывает по экспоненциальному закону: $y(t) = b e^{-\frac{c}{a} t}$.

%=============================================================
\task{60}
%=============================================================
\subsection*{Условие}
Требуется узнать, за какое время упадет на поверхность Луны камень с высоты $h$. Пусть $x(t)$ – высота камня над поверхностью в момент времени $t$. Согласно закону свободного падения, все тела падают в поле силы тяжести с постоянным ускорением $a$ (для поверхности Луны $a = 1/6 g$).

\begin{theory}
Ускорение тела является второй производной от его координаты высоты по времени: $\ddot{x}(t) = -a$. Свободное падение с высоты $h$ без начальной скорости описывается начальными условиями: $x(0) = h$, $\dot{x}(0) = 0$.
\end{theory}

\step{1} Запишем дифференциальное уравнение движения камня и проинтегрируем его один раз по времени:
\[
\frac{d^2 x}{dt^2} = -a \implies \frac{dx}{dt} = -at + C_1.
\]
Поскольку начальная скорость равна нулю, то $\dot{x}(0) = 0 \implies C_1 = 0$.

\step{2} Проинтегрируем уравнение скорости ещё раз:
\[
x(t) = -\frac{a t^2}{2} + C_2.
\]
Поскольку в момент времени $t = 0$ высота равна $h$, получаем $C_2 = h$. \\
Закон движения камня:
\[
x(t) = h - \frac{a t^2}{2}.
\]

\step{3} Камень упадёт на поверхность Луны в момент времени $T$, когда его высота станет равной нулю ($x(T) = 0$):
\[
h - \frac{a T^2}{2} = 0 \implies T^2 = \frac{2h}{a} \implies T = \sqrt{\frac{2h}{a}}.
\]

\step{4} Подставим лунное ускорение $a = \frac{1}{6} g$:
\[
T = \sqrt{\frac{2h}{\frac{1}{6}g}} = \sqrt{\frac{12h}{g}}.
\]

\answer
Дифференциальное уравнение движения: $\ddot{x}(t) = -\dfrac{1}{6}g$.\\
Камень упадёт на поверхность Луны за время $T = \sqrt{\dfrac{12h}{g}}$.

%=============================================================
\task{61}
%=============================================================
\subsection*{Условие}
Рассмотреть уравнение $(y')^2 = 4y$.
\begin{enumerate}
    \item Найти общее решение.
    \item Найти особое решение.
\end{enumerate}

\begin{theory}
\begin{itemize}
    \item \textbf{Общим решением} ОДУ первого порядка называется семейство функций $y = \varphi(x, C)$, зависящее от произвольной постоянной $C$ и удовлетворяющее уравнению.
    \item \textbf{Особым решением} называется решение, во всех точках которого нарушается свойство единственности решения задачи Коши (то есть через каждую точку этого решения проходит по крайней мере ещё одна интегральная кривая с тем же наклоном касательной).
\end{itemize}
\end{theory}

\step{1} Найдём общее решение. Перепишем уравнение в виде:
\[
y' = \pm 2\sqrt{y} \quad (\text{при } y \geq 0).
\]
Разделим переменные (предполагая $y > 0$):
\[
\frac{dy}{2\sqrt{y}} = \pm dx.
\]
Интегрируя обе части, получаем:
\[
\sqrt{y} = \pm x + C.
\]
Поскольку левая часть $\sqrt{y}$ неотрицательна, правая часть также должна быть неотрицательной. Возводя в квадрат, запишем семейство решений в виде:
\[
y = (x + C)^2 \quad (\text{при } x + C \geq 0) \qquad \text{и} \qquad y = (C - x)^2 \quad (\text{при } C - x \geq 0).
\]
Объединяя эти ветви, мы получаем семейство парабол:
\[
y = (x - C)^2 \quad (\text{для ветвей парабол с вершиной на оси } Ox).
\]

\step{2} Найдём особые решения. 
При делении на $\sqrt{y}$ мы могли потерять решение $y \equiv 0$. Прямой подстановкой убеждаемся, что функция:
\[
y(x) \equiv 0
\]
является решением исходного уравнения ($(0)^2 = 4\cdot 0$).

\step{3} Покажем, что $y(x) \equiv 0$ — особое решение.
Возьмём любую точку $(x_0, 0)$ на оси абсцисс. Через неё проходит решение $y_1(x) = 0$. Однако через эту же точку проходит и решение из общего семейства $y_2(x) = (x - x_0)^2$. При этом их производные в точке $x_0$ совпадают:
\[
y_1'(x_0) = 0, \qquad y_2'(x_0) = \left. 2(x - x_0) \right|_{x=x_0} = 0.
\]
Таким образом, в любой точке $(x_0, 0)$ нарушается единственность решения задачи Коши. Следовательно, $y \equiv 0$ — особое решение.

\answer
1. Общее решение: $y = (x - C)^2$ (для полуветвей парабол при $x \geq C$ или $x \leq C$).\\
2. Особое решение: $y(x) \equiv 0$.

%=============================================================
\task{62}
%=============================================================
\subsection*{Условие}
Проверьте выполнение условия Липшица для уравнений:
\begin{enumerate}
    \item $\begin{cases} y' = |y| \\ y(0) = 0 \end{cases}$
    \item $\begin{cases} y' = \sqrt{|y|} \\ y(0) = 0 \end{cases}$
\end{enumerate}

\begin{theory}
Функция $f(t, y)$ удовлетворяет условию Липшица по переменной $y$ в области $D$, если существует постоянная $L > 0$ (константа Липшица), такая что для любых $(t, y_1), (t, y_2) \in D$ выполнено неравенство:
\[
|f(t, y_1) - f(t, y_2)| \leq L |y_1 - y_2|.
\]
Если условие Липшица выполнено в окрестности начальной точки, то теорема Пикара гарантирует существование и единственность решения задачи Коши.
\end{theory}

\step{1} Исследуем первое уравнение: $f(y) = |y|$.
Проверим липшицевость в любой окрестности нуля:
\[
|f(y_1) - f(y_2)| = ||y_1| - |y_2|| \leq |y_1 - y_2|.
\]
Неравенство треугольника выполнено с константой $L = 1$. Условие Липшица \textbf{выполняется}. Следовательно, решение данной задачи Коши существует и единственно. Это тривиальное решение $y(t) \equiv 0$.

\step{2} Исследуем второе уравнение: $f(y) = \sqrt{|y|}$.
Рассмотрим отношение при $y_2 = 0$:
\[
\frac{|f(y_1) - f(0)|}{|y_1 - 0|} = \frac{\sqrt{|y_1|}}{|y_1|} = \frac{1}{\sqrt{|y_1|}}.
\]
При приближении $y_1 \to 0$ данное отношение стремится к бесконечности:
\[
\lim_{y_1 \to 0} \frac{1}{\sqrt{|y_1|}} = +\infty.
\]
Следовательно, не существует конечной константы $L$, которая могла бы ограничить это отношение в любой окрестности нуля. Условие Липшица в окрестности точки $y = 0$ \textbf{не выполняется}. Из-за этого единственность решения задачи Коши нарушается (помимо тривиального решения $y \equiv 0$, существуют решения вида $y = \pm \frac{1}{4}t^2$).

\answer
1. Для уравнения $y' = |y|$ условие Липшица \textbf{выполнено} (с константой $L = 1$), решение единственно.\\
2. Для уравнения $y' = \sqrt{|y|}$ условие Липшица в окрестности нуля \textbf{не выполнено}, единственность решения нарушается.

%=============================================================
\task{63}
%=============================================================
\subsection*{Условие}
Маша приготовила чашку кофе. Маша любит пить кофе только когда его температура достигает ровно 60 градусов по Цельсию. Сначала в момент времени $t = 0$ минут Маша измерила температуру кофе, и она оказалась 89 градусов. Через минуту Маша снова измерила температуру кофе, и она оказалась 85 градусов. Температура помещения (температура окружающей среды) составляет 22 градуса. Когда Маша может начать пить кофе?

\begin{theory}
\textbf{Закон охлаждения Ньютона:} скорость остывания тела пропорциональна разности температур тела $T(t)$ и окружающей среды $T_{\text{ср}}$:
\[
\frac{dT}{dt} = -k (T - T_{\text{ср}}) \implies T(t) = T_{\text{ср}} + (T_0 - T_{\text{ср}}) e^{-kt},
\]
где $T_0$ — начальная температура тела в момент времени $t = 0$.
\end{theory}

\step{1} Подставим известные данные задачи: $T_{\text{ср}} = 22^\circ\text{C}$ и начальную температуру $T_0 = 89^\circ\text{C}$:
\[
T(t) = 22 + (89 - 22) e^{-kt} = 22 + 67 e^{-kt}.
\]

\step{2} Используем второе измерение: в момент $t = 1$ минута температура составляла $T(1) = 85^\circ\text{C}$:
\[
85 = 22 + 67 e^{-k} \implies 63 = 67 e^{-k} \implies e^{-k} = \frac{63}{67} \approx 0.9403.
\]
Тогда закон остывания кофе принимает вид:
\[
T(t) = 22 + 67 \left(\frac{63}{67}\right)^t.
\]

\step{3} Найдём время $t$, при котором температура кофе опустится до $60^\circ\text{C}$:
\[
60 = 22 + 67 \left(\frac{63}{67}\right)^t \implies 38 = 67 \left(\frac{63}{67}\right)^t \implies \left(\frac{63}{67}\right)^t = \frac{38}{67}.
\]
Прологарифмируем обе части:
\[
t \ln\left(\frac{63}{67}\right) = \ln\left(\frac{38}{67}\right) \implies t = \frac{\ln(38/67)}{\ln(63/67)} = \frac{\ln(67/38)}{\ln(67/63)}.
\]

\step{4} Проведём численный расчёт:
\[
\ln(67/38) \approx 0.5671, \qquad \ln(67/63) \approx 0.0616,
\]
\[
t \approx \frac{0.5671}{0.0616} \approx 9.21 \text{ минут}.
\]
Переведём дробную часть в секунды: $0.21 \text{ мин} \approx 13$ секунд.

\answer
Маша может начать пить кофе примерно через $9.21$ минут (или 9 минут 13 секунд) после первого измерения.

%=============================================================
\task{64}
%=============================================================
\subsection*{Условие}
Найдите кривую $y(x)$, проходящую через точку $(1, 1)$, у которой отрезок любой касательной, заключенный между осями координат, делится точкой касания пополам.

\begin{theory}
Уравнение касательной к кривой $y = y(x)$ в точке $(x_0, y_0)$ имеет вид:
\[
Y - y_0 = y'(x_0)(X - x_0).
\]
Точки пересечения этой касательной с осями координат $X = 0$ (ось $Oy$) и $Y = 0$ (ось $Ox$) обозначим как $A$ и $B$. Условие деления отрезка пополам означает, что точка касания является серединой отрезка $AB$.
\end{theory}

\step{1} Найдём координаты точек пересечения касательной с осями:
\begin{itemize}
    \item Пересечение с осью $Oy$ ($X = 0$): \ $Y_A = y_0 - x_0 y'(x_0) \implies A = (0, \ y_0 - x_0 y'(x_0))$.
    \item Пересечение с осью $Ox$ ($Y = 0$): \ $0 - y_0 = y'(x_0)(X_B - x_0) \implies B = \left(x_0 - \frac{y_0}{y'(x_0)}, \ 0\right)$.
\end{itemize}

\step{2} Точка касания $M(x_0, y_0)$ является серединой отрезка $AB$. По формуле координат середины отрезка (для абсциссы):
\[
x_0 = \frac{X_A + X_B}{2} \implies x_0 = \frac{0 + x_0 - \frac{y_0}{y'(x_0)}}{2} \implies 2x_0 = x_0 - \frac{y_0}{y'(x_0)}.
\]
Упростим уравнение:
\[
x_0 = -\frac{y_0}{y'(x_0)} \implies y'(x_0) = -\frac{y_0}{x_0}.
\]

\step{3} Поскольку это свойство выполнено в любой точке кривой, запишем дифференциальное уравнение:
\[
\frac{dy}{dx} = -\frac{y}{x}.
\]

\step{4} Разделим переменные и проинтегрируем:
\[
\frac{dy}{y} = -\frac{dx}{x} \implies \ln|y| = -\ln|x| + \ln C \implies y = \frac{C}{x}.
\]

\step{5} Используем начальное условие: кривая проходит через точку $(1, 1)$, следовательно:
\[
1 = \frac{C}{1} \implies C = 1.
\]
Получаем уравнение кривой: $y = \frac{1}{x}$ (при $x > 0$).

\answer
Искомая кривая — ветвь гиперболы $y = \dfrac{1}{x}$ (для $x > 0$).

%=============================================================
\task{65}
%=============================================================
\subsection*{Условие}
Четыре черепашки находятся в вершинах квадрата со стороной $a$. Каждая начинает двигаться с постоянной скоростью $v$ в направлении своей соседки по часовой стрелке. Найти траекторию движения.

\begin{theory}
В силу симметрии задачи, в любой момент времени черепашки будут находиться в вершинах квадрата, центр которого неподвижен и совпадает с центром исходного квадрата. Удобно перейти в полярную систему координат $(r, \theta)$ с началом в центре квадрата.
\end{theory}

\step{1} В полярных координатах радиальный вектор направлен от центра к черепашке. Траектория соседки всегда повернута на угол $90^\circ$ относительно данного радиального вектора. 
Поскольку черепашка бежит точно на свою соседку, угол между вектором её скорости $\vec{v}$ и направлением к центру (радиальным направлением внутрь) всегда постоянен и равен $45^\circ$ ($\pi/4$).

\step{2} Запишем проекции скорости на полярные оси:
\begin{itemize}
    \item Радиальная скорость (направлена к центру, поэтому знак минус):
    \[
    \frac{dr}{dt} = -v \cos\left(\frac{\pi}{4}\right) = -\frac{v}{\sqrt{2}}.
    \]
    \item Трансверсальная (угловая) скорость:
    \[
    r \frac{d\theta}{dt} = v \sin\left(\frac{\pi}{4}\right) = \frac{v}{\sqrt{2}}.
    \]
\end{itemize}

\step{3} Разделим радиальное уравнение на угловое, чтобы исключить параметр времени $t$:
\[
\frac{dr}{r d\theta} = \frac{-v/\sqrt{2}}{v/\sqrt{2}} = -1 \implies \frac{dr}{r} = -d\theta.
\]

\step{4} Проинтегрируем полученное дифференциальное уравнение:
\[
\ln r = -\theta + C_1 \implies r(\theta) = C e^{-\theta}.
\]

\step{5} Найдём константу $C$. В начальный момент времени $\theta = 0$ расстояние черепашки до центра квадрата равно половине его диагонали:
\[
r(0) = \frac{a}{\sqrt{2}} \implies C = \frac{a}{\sqrt{2}}.
\]
Получаем уравнение траектории в полярных координатах:
\[
r(\theta) = \frac{a}{\sqrt{2}} e^{-\theta}.
\]
Это уравнение описывает логарифмическую спираль, скручивающуюся к центру квадрата.

\answer
Траектория движения черепашек в полярных координатах описывается логарифмической спиралью: $r(\theta) = \dfrac{a}{\sqrt{2}} e^{-\theta}$ (где центр координат совмещен с центром квадрата).
%=============================================================
\task{66}
%=============================================================
\subsection*{Условие}
Парашютист массы $m$ совершает прыжок с большой высоты. Сила сопротивления воздуха пропорциональна скорости. Составьте и решите уравнение движения методом Бернулли, используя второй закон Ньютона.

\begin{theory}
На парашютиста действуют две силы: сила тяжести $F_g = mg$ (направлена вниз) и сила сопротивления воздуха $F_d = -kv$ (направлена вверх, пропорциональна скорости $v(t)$, где $k > 0$ — коэффициент сопротивления).
По второму закону Ньютона $m a = \sum F \implies m \dfrac{dv}{dt} = mg - kv$.
Метод Бернулли для решения линейного уравнения первого порядка заключается в замене $v(t) = u(t)w(t)$.
\end{theory}

\step{1} Запишем дифференциальное уравнение движения:
\[
\frac{dv}{dt} + \frac{k}{m} v = g.
\]

\step{2} Применим подстановку Бернулли $v = uw$, тогда $v' = u'w + uw'$. Подставим в уравнение:
\[
u'w + uw' + \frac{k}{m} uw = g \implies u'w + u\left(w' + \frac{k}{m}w\right) = g.
\]

\step{3} Положим выражение в скобках равным нулю для нахождения вспомогательной функции $w(t)$:
\[
w' + \frac{k}{m}w = 0 \implies \frac{dw}{w} = -\frac{k}{m} dt \implies \ln w = -\frac{k}{m}t \implies w(t) = e^{-\frac{k}{m}t}.
\]

\step{4} Подставим $w(t)$ обратно в уравнение для нахождения $u(t)$:
\[
u' e^{-\frac{k}{m}t} = g \implies u' = g e^{\frac{k}{m}t}.
\]
Интегрируем:
\[
u(t) = g \int e^{\frac{k}{m}t}\,dt = \frac{mg}{k} e^{\frac{k}{m}t} + C.
\]

\step{5} Запишем общее решение для скорости $v(t) = u(t)w(t)$:
\[
v(t) = \left(\frac{mg}{k} e^{\frac{k}{m}t} + C\right) e^{-\frac{k}{m}t} = \frac{mg}{k} + C e^{-\frac{k}{m}t}.
\]
Если принять начальную скорость $v(0) = v_0$, то $v_0 = \frac{mg}{k} + C \implies C = v_0 - \frac{mg}{k}$. Получаем:
\[
v(t) = \frac{mg}{k} + \left(v_0 - \frac{mg}{k}\right) e^{-\frac{k}{m}t}.
\]
(При $t \to +\infty$ скорость стремится к установившемуся пределу $v_{\text{пред}} = mg/k$).

\answer
Дифференциальное уравнение: $\dfrac{dv}{dt} + \dfrac{k}{m}v = g$. \\
Его решение методом Бернулли: $v(t) = \dfrac{mg}{k} + C e^{-\frac{k}{m}t}$.

%=============================================================
\task{67}
%=============================================================
\subsection*{Условие}
Дана цепь с последовательно соединенными индуктивностью $L$ и сопротивлением $R$. К цепи приложено переменное напряжение $E(t) = E_0 \sin(\omega t)$. Составьте и решите уравнение для зависимости тока $I(t)$.

\begin{theory}
Согласно второму закону Кирхгофа, падение напряжения на катушке индуктивности $U_L = L \dfrac{dI}{dt}$ и на резисторе $U_R = RI$ в сумме равны ЭДС источника $E(t)$:
\[
L \frac{dI}{dt} + R I = E_0 \sin(\omega t) \implies \frac{dI}{dt} + \frac{R}{L} I = \frac{E_0}{L} \sin(\omega t).
\]
\end{theory}

\step{1} Это линейное уравнение первого порядка. Решим его с помощью интегрирующего множителя $\mu(t) = e^{\int \frac{R}{L} dt} = e^{\frac{R}{L}t}$:
\[
\frac{d}{dt}\left( I(t) e^{\frac{R}{L}t} \right) = \frac{E_0}{L} e^{\frac{R}{L}t} \sin(\omega t).
\]

\step{2} Интегрируем обе части:
\[
I(t) e^{\frac{R}{L}t} = \frac{E_0}{L} \int e^{\frac{R}{L}t} \sin(\omega t)\,dt + C.
\]

\step{3} Вычислим интеграл методом интегрирования по частям (дважды) или используя табличную формулу $\int e^{at}\sin(bt)dt = \frac{e^{at}}{a^2+b^2}(a\sin(bt) - b\cos(bt))$:
\[
\int e^{\frac{R}{L}t} \sin(\omega t)\,dt = \frac{e^{\frac{R}{L}t}}{\left(\frac{R}{L}\right)^2 + \omega^2} \left( \frac{R}{L} \sin(\omega t) - \omega \cos(\omega t) \right).
\]

\step{4} Подставим интеграл в выражение для тока:
\[
I(t) e^{\frac{R}{L}t} = \frac{E_0}{L} \cdot \frac{e^{\frac{R}{L}t}}{\frac{R^2 + \omega^2 L^2}{L^2}} \left( \frac{R}{L} \sin(\omega t) - \omega \cos(\omega t) \right) + C
\]
\[
= \frac{E_0 e^{\frac{R}{L}t}}{R^2 + \omega^2 L^2} \Big( R \sin(\omega t) - \omega L \cos(\omega t) \Big) + C.
\]

\step{5} Разделим обе части на $e^{\frac{R}{L}t}$:
\[
I(t) = \frac{E_0}{R^2 + \omega^2 L^2} \Big( R \sin(\omega t) - \omega L \cos(\omega t) \Big) + C e^{-\frac{R}{L}t}.
\]
Используя тригонометрическое преобразование, введём фазовый сдвиг $\varphi = \mathrm{arctg}\left(\frac{\omega L}{R}\right)$:
\[
I(t) = \frac{E_0}{\sqrt{R^2 + \omega^2 L^2}} \sin(\omega t - \varphi) + C e^{-\frac{R}{L}t}.
\]

\answer
Зависимость тока от времени: $I(t) = \dfrac{E_0}{\sqrt{R^2 + \omega^2 L^2}} \sin(\omega t - \varphi) + C e^{-\frac{R}{L}t}$, где $\varphi = \mathrm{arctg}\left(\dfrac{\omega L}{R}\right)$.

%=============================================================
\task{68}
%=============================================================
\subsection*{Условие}
Популяция изолированного вида растет со скоростью, пропорциональной их количеству $y$, но при этом ограничена ресурсами. Это описывается логистическим уравнением:
\[
y' = ry - ky^2, \quad (r, k > 0).
\]
Решите уравнение.

\begin{theory}
Уравнение $y' = ry - ky^2$ является уравнением Бернулли (при $n=2$) или уравнением с разделяющимися переменными. Будем решать его как уравнение с разделяющимися переменными.
\end{theory}

\step{1} Разделим переменные в уравнении (при $y \neq 0$ и $y \neq r/k$):
\[
\frac{dy}{y(r - ky)} = dt.
\]

\step{2} Разложим левую часть на простейшие дроби:
\[
\frac{1}{y(r - ky)} = \frac{1/r}{y} + \frac{k/r}{r - ky}.
\]
Тогда интеграл принимает вид:
\[
\frac{1}{r} \int \left( \frac{1}{y} + \frac{k}{r - ky} \right) dy = \int dt \implies \frac{1}{r} \Big( \ln|y| - \ln|r - ky| \Big) = t + C_0.
\]

\step{3} Умножим на $r$ и потенцируем:
\[
\ln\left| \frac{y}{r - ky} \right| = rt + C_1 \implies \frac{y}{r - ky} = C e^{rt}.
\]

\step{4} Выразим $y$ из полученного алгебраического уравнения:
\[
y = C e^{rt} (r - ky) \implies y(1 + k C e^{rt}) = r C e^{rt} \implies y(t) = \frac{r C e^{rt}}{1 + k C e^{rt}}.
\]
Разделим числитель и знаменатель на $C e^{rt}$ и введём новую константу $C_2 = 1/C$:
\[
y(t) = \frac{r}{k + C_2 e^{-rt}}.
\]
Если задана начальная численность популяции $y(0) = y_0$, то $y_0 = \frac{r}{k+C_2} \implies C_2 = \frac{r}{y_0} - k$. Получаем:
\[
y(t) = \frac{r y_0}{k y_0 + (r - k y_0)e^{-rt}}.
\]

\answer
Общее решение логистического уравнения: $y(t) = \dfrac{r y_0}{k y_0 + (r - k y_0)e^{-rt}}$ (при $t \to +\infty$ численность стабилизируется на уровне $y_{\text{пред}} = r/k$).

%=============================================================
\task{69}
%=============================================================
\subsection*{Условие}
Решите уравнение $(y + xy^2)dx - xdy = 0$, найдя интегрирующий множитель вида $\mu(y)$.

\begin{theory}
Для уравнения $P(x,y)dx + Q(x,y)dy = 0$ функция $\mu(y)$ является интегрирующим множителем, если после умножения на неё уравнение становится уравнением в полных дифференциалах. Условие липшицевости/тождественности смешанных производных принимает вид:
\[
\frac{\partial (\mu P)}{\partial y} = \frac{\partial (\mu Q)}{\partial x} \implies \mu' P + \mu \frac{\partial P}{\partial y} = \mu \frac{\partial Q}{\partial x}.
\]
\end{theory}

\step{1} Выпишем коэффициенты уравнения: $P = y + xy^2$, $Q = -x$.
Их частные производные: $\frac{\partial P}{\partial y} = 1 + 2xy$, $\frac{\partial Q}{\partial x} = -1$.
Подставим в уравнение для $\mu(y)$:
\[
\mu' (y + xy^2) + \mu (1 + 2xy) = \mu (-1) \implies \mu' y(1 + xy) = -2\mu (1 + xy).
\]

\step{2} Разделим обе части на $(1 + xy)$ (при условии $1 + xy \neq 0$):
\[
\mu' y = -2\mu \implies \frac{d\mu}{\mu} = -2 \frac{dy}{y}.
\]
Интегрируя, находим интегрирующий множитель:
\[
\ln \mu = -2 \ln y \implies \mu(y) = \frac{1}{y^2}.
\]

\step{3} Умножим исходное уравнение на $\mu(y) = \frac{1}{y^2}$:
\[
\left( \frac{1}{y} + x \right) dx - \frac{x}{y^2} dy = 0.
\]
Это уравнение в полных дифференциалах. Найдём его потенциал $u(x, y)$:
\[
\frac{\partial u}{\partial x} = \frac{1}{y} + x \implies u(x, y) = \frac{x}{y} + \frac{x^2}{2} + \psi(y).
\]
Дифференцируем по $y$ и приравниваем к коэффициенту при $dy$:
\[
\frac{\partial u}{\partial y} = -\frac{x}{y^2} + \psi'(y) = -\frac{x}{y^2} \implies \psi'(y) = 0 \implies \psi(y) = C_0.
\]

\step{4} Запишем общий интеграл уравнения $u(x, y) = C$:
\[
\frac{x}{y} + \frac{x^2}{2} = C \implies y(x) = \frac{2x}{2C - x^2}.
\]

\answer
Интегрирующий множитель: $\mu(y) = \dfrac{1}{y^2}$. \\
Общее решение: $\dfrac{x}{y} + \dfrac{x^2}{2} = C$ или $y(x) = \dfrac{2x}{C - x^2}$.

%=============================================================
\task{70}
%=============================================================
\subsection*{Условие}
Гибкая однородная нить подвешена за два конца. Найти уравнение кривой, по которой расположится нить под действием собственного веса (уравнение цепной линии).

\begin{theory}
Рассмотрим элемент нити длиной $ds = \sqrt{1 + (y')^2} dx$. На него действуют три силы: сила тяжести $dF_g = \rho_0 g ds$ (где $\rho_0$ — линейная плотность) и силы натяжения на концах элемента.
Условия статического равновесия элемента приводят к дифференциальному уравнению:
\[
y'' = \frac{1}{a} \sqrt{1 + (y')^2}, \qquad \text{где } a = \frac{T_0}{\rho_0 g} \text{ (параметр цепной линии)}.
\]
\end{theory}

\step{1} Для решения уравнения $y'' = \frac{1}{a} \sqrt{1 + (y')^2}$ понизим порядок, сделав замену $p(x) = y'(x)$:
\[
\frac{dp}{dx} = \frac{1}{a} \sqrt{1 + p^2}.
\]

\step{2} Разделим переменные и проинтегрируем:
\[
\frac{dp}{\sqrt{1 + p^2}} = \frac{dx}{a} \implies \ln\left(p + \sqrt{1 + p^2}\right) = \frac{x}{a} + C_1 \implies \mathrm{arsinh}(p) = \frac{x}{a} + C_1.
\]
Выберем систему координат так, чтобы в низшей точке нити ($x = 0$) касательная была горизонтальной ($y'(0) = 0 \implies p(0) = 0$). Отсюда $C_1 = 0$:
\[
p = \sinh\left(\frac{x}{a}\right).
\]

\step{3} Вернёмся к исходной переменной $y(x)$:
\[
\frac{dy}{dx} = \sinh\left(\frac{x}{a}\right) \implies y(x) = a \cosh\left(\frac{x}{a}\right) + C_2.
\]

\step{4} Поместим начало координат на расстояние $a$ ниже самой нижней точки цепи, то есть потребуем $y(0) = a$. Тогда:
\[
a = a \cosh(0) + C_2 \implies a = a + C_2 \implies C_2 = 0.
\]
Получаем каноническое уравнение цепной линии:
\[
y(x) = a \cosh\left(\frac{x}{a}\right).
\]

\answer
Уравнение цепной линии имеет вид гиперболического косинуса: $y(x) = a \cosh\left(\dfrac{x}{a}\right)$, где $a = \dfrac{T_0}{\rho_0 g}$.

%=============================================================
\task{71}
%=============================================================
\subsection*{Условие}
Проверьте, гарантирует ли теорема единственность решения задачи Коши для уравнения $y'' = \sqrt{y'}$ с условиями $y(0) = 0$, $y'(0) = 0$. Найдите два различных решения.

\begin{theory}
Для системы дифференциальных уравнений первого порядка теорема Коши — Липшица (Пикара) гарантирует существование и единственность решения, если правая часть удовлетворяет условию Липшица по зависимым переменным в окрестности начальной точки.
Для уравнения $y'' = f(x, y, y')$ сделаем замену $y' = z$. Получаем систему:
\[
\begin{cases}
y' = z, \\
z' = \sqrt{z}.
\end{cases}
\]
Правая часть второго уравнения $g(y, z) = \sqrt{z}$ не является дифференцируемой в точке $z = 0$, так как $\frac{\partial g}{\partial z} = \frac{1}{2\sqrt{z}} \to \infty$ при $z \to 0^+$. Соответственно, в окрестности начальной точки $z(0) = 0$ условие Липшица не выполняется.
\end{theory}

\step{1} Поскольку условие Липшица для функции $g(z) = \sqrt{z}$ в точке $z = 0$ не выполнено, классическая теорема единственности решения задачи Коши \textbf{не гарантирует} единственность решения данной задачи.

\step{2} Найдем первое (тривиальное) решение. Очевидно, что тождественный нуль:
\[
y_1(x) \equiv 0
\]
удовлетворяет уравнению $y'' = \sqrt{y'}$ ($0 = \sqrt{0}$) и начальным условиям $y_1(0) = 0$, $y_1'(0) = 0$.

\step{3} Найдем второе решение. Сделаем замену $y' = z$, тогда $z' = \sqrt{z}$. Разделим переменные (для $z > 0$):
\[
\frac{dz}{\sqrt{z}} = dx \implies 2\sqrt{z} = x + C_1.
\]
Используя начальное условие $y'(0) = 0 \implies z(0) = 0$, находим $C_1 = 0$. Получаем:
\[
2\sqrt{z} = x \implies z = \frac{x^2}{4} \quad (x \geq 0).
\]

\step{4} Интегрируем $y' = z = \frac{x^2}{4}$:
\[
y(x) = \int \frac{x^2}{4}\,dx = \frac{x^3}{12} + C_2.
\]
Из условия $y(0) = 0$ находим $C_2 = 0$. Получаем второе решение:
\[
y_2(x) = \frac{x^3}{12} \quad (x \geq 0).
\]
Легко проверить, что $y_2''(x) = \frac{x}{2}$ и $\sqrt{y_2'(x)} = \sqrt{\frac{x^2}{4}} = \frac{x}{2}$, то есть равенство выполняется, и условия $y_2(0) = 0, y_2'(0) = 0$ соблюдены.

\answer
Теорема Коши — Липшица единственность решения \textbf{не гарантирует}. \\
Два различных решения: $y_1(x) = 0$ \ и \ $y_2(x) = \dfrac{x^3}{12}$ (при $x \geq 0$).

%=============================================================
\task{72}
%=============================================================
\subsection*{Условие}
Решите задачу о второй космической скорости: найти с какой наименьшей скоростью $v_0$ надо бросить тело вертикально вверх, чтобы оно не вернулось на Землю.

\begin{theory}
По закону всемирного тяготения Ньютона на тело массой $m$ на расстоянии $r$ от центра Земли (масса Земли $M$, радиус Земли $R$) действует сила:
\[
F = -G \frac{Mm}{r^2}.
\]
Согласно второму закону Ньютона:
\[
m \frac{d^2 r}{dt^2} = -G \frac{Mm}{r^2} \implies \frac{dv}{dt} = -G \frac{M}{r^2}.
\]
Используя связь $v = \frac{dr}{dt}$, перепишем ускорение как $\frac{dv}{dt} = \frac{dv}{dr} \frac{dr}{dt} = v \frac{dv}{dr}$.
\end{theory}

\step{1} Запишем дифференциальное уравнение для скорости $v(r)$ как функции от расстояния $r$:
\[
v \frac{dv}{dr} = -G \frac{M}{r^2}.
\]

\step{2} Разделим переменные и проинтегрируем от поверхности Земли ($r = R$, $v = v_0$) до некоторого расстояния $r$:
\[
\int_{v_0}^v v\,dv = -GM \int_R^r \frac{dr}{r^2} \implies \frac{v^2}{2} - \frac{v_0^2}{2} = GM \left( \frac{1}{r} - \frac{1}{R} \right).
\]
Отсюда получаем закон изменения скорости:
\[
v^2 = v_0^2 + 2GM \left( \frac{1}{r} - \frac{1}{R} \right).
\]

\step{3} Чтобы тело не вернулось на Землю, его скорость должна оставаться вещественной ($v^2 \geq 0$) при сколь угодно больших удалениях, то есть при $r \to \infty$. На бесконечности член $\frac{1}{r} \to 0$:
\[
v_{\infty}^2 = v_0^2 - \frac{2GM}{R} \geq 0 \implies v_0 \geq \sqrt{\frac{2GM}{R}}.
\]

\step{4} Наименьшая скорость $v_0$ (вторая космическая скорость) равна:
\[
v_0 = \sqrt{\frac{2GM}{R}}.
\]
Учитывая, что ускорение свободного падения на поверхности Земли равно $g = \frac{GM}{R^2}$, формулу можно записать как:
\[
v_0 = \sqrt{2gR}.
\]
Для Земли ($g \approx 9.81 \text{ м/с}^2$, $R \approx 6371 \text{ км}$) получаем $v_0 \approx 11.2 \text{ км/с}$.

\answer
Вторая космическая скорость вычисляется по формуле: $v_0 = \sqrt{\dfrac{2GM}{R}} = \sqrt{2gR}$ (для Земли $v_0 \approx 11.2 \text{ км/с}$).

%=============================================================
\task{73}
%=============================================================
\subsection*{Условие}
Рассмотрим объект массой $m$, смещение которого от положения равновесия равно $x(t)$. На него действуют две основные силы:
\begin{enumerate}
    \item Сила упругости (возвращающая): $F_s = -kx$, где $k$ — жесткость пружины.
    \item Сила вязкого трения (тормозящая): $F_d = -c\dot{x}$, где $c$ — коэффициент трения.
\end{enumerate}
Составьте дифференциальное уравнение, моделирующее движение этого объекта.

\begin{theory}
Второй закон Ньютона утверждает, что произведение массы тела на его ускорение равно сумме всех действующих на него сил:
\[
m a = \sum F \implies m \ddot{x} = F_s + F_d,
\]
где ускорение $\ddot{x} = \frac{d^2 x}{dt^2}$, а скорость $\dot{x} = \frac{dx}{dt}$.
\end{theory}

\step{1} Подставим выражения для сил упругости и вязкого трения во второй закон Ньютона:
\[
m \ddot{x} = -kx - c\dot{x}.
\]

\step{2} Перенесем все члены уравнения в левую часть:
\[
m \ddot{x} + c\dot{x} + kx = 0.
\]

\step{3} Обычно это уравнение делят на $m$ и представляют в канонической форме колебательного процесса:
\[
\ddot{x} + 2\beta \dot{x} + \omega_0^2 x = 0,
\]
где $\beta = \frac{c}{2m}$ — коэффициент затухания, а $\omega_0 = \sqrt{\frac{k}{m}}$ — собственная частота недемпфированных колебаний.

\answer
Дифференциальное уравнение затухающих колебаний имеет вид: $m \frac{d^2 x}{dt^2} + c \frac{dx}{dt} + kx = 0$.

%=============================================================
\task{74}
%=============================================================
\subsection*{Условие}
Докажите, что функции $y_1 = x$ и $y_2 = x^2$ линейно независимы, но не могут быть решениями одного и того же линейного однородного дифференциального уравнения (ЛОУ) второго порядка с непрерывными коэффициентами на интервале, содержащем точку $x = 0$.

\begin{theory}
\begin{itemize}
    \item Функции $y_1(x), y_2(x)$ линейно независимы на интервале, если равенство $c_1 y_1(x) + c_2 y_2(x) = 0$ выполняется для всех $x$ только при $c_1 = c_2 = 0$.
    \item \textbf{Теорема Лиувилля:} если $y_1(x)$ и $y_2(x)$ — решения ЛОУ второго порядка $y'' + p(x)y' + q(x)y = 0$ с непрерывными на интервале $I$ коэффициентами, то их определитель Вронского $W(x)$ либо тождественно равен нулю на $I$, либо ни в одной точке интервала $I$ в нуль не обращается.
\end{itemize}
\end{theory}

\step{1} Докажем линейную независимость функций $y_1 = x$ и $y_2 = x^2$. Составим линейную комбинацию:
\[
c_1 x + c_2 x^2 = 0 \quad \forall x \in I.
\]
Подставим $x = 1$: $c_1 + c_2 = 0 \implies c_1 = -c_2$. \\
Подставим $x = 2$: $2c_1 + 4c_2 = 0 \implies -2c_2 + 4c_2 = 0 \implies 2c_2 = 0 \implies c_2 = 0$. \\
Отсюда $c_1 = 0$. Поскольку оба коэффициента обязаны быть нулевыми, функции $y_1, y_2$ линейно независимы.

\step{2} Вычислим определитель Вронского (Вронскиан) для системы функций $\{x, x^2\}$:
\[
W(x) = \begin{vmatrix} y_1 & y_2 \\ y_1' & y_2' \end{vmatrix} = \begin{vmatrix} x & x^2 \\ 1 & 2x \end{vmatrix} = x(2x) - 1(x^2) = x^2.
\]

\step{3} Проанализируем поведение Вронскиана на интервале $I$, содержащем $x = 0$:
\begin{itemize}
    \item В точке $x = 0$: \ $W(0) = 0^2 = 0$.
    \item В любой точке $x \neq 0$: \ $W(x) = x^2 \neq 0$.
\end{itemize}
Таким образом, Вронскиан обращается в нуль в одной точке ($x=0$), но не равен нулю тождественно на всём интервале.

\step{4} В силу теоремы Лиувилля, для решений ЛОУ второго порядка с непрерывными коэффициентами такое поведение Вронскиана невозможно. Следовательно, эти функции не могут являться решениями одного ЛОУ второго порядка на интервале, содержащем $x = 0$. \quad$\blacksquare$

\answer
Линейная независимость доказана. Функции не могут быть решениями одного ЛОУ второго порядка, так как их Вронскиан $W(x) = x^2$ равен нулю в $x=0$, но отличен от нуля при $x \neq 0$, что противоречит теореме Лиувилля.

%=============================================================
\task{75}
%=============================================================
\subsection*{Условие}
(Метод Абеля) Используя формулу Остроградского–Лиувилля, найдите второе решение $y_2$ уравнения $y'' + p(x)y' + q(x)y = 0$, если известно одно частное решение $y_1$:
\[
y_2(x) = y_1(x) \int \frac{e^{-\int p(x)dx}}{y_1^2(x)}\,dx.
\]

\begin{theory}
Формула Остроградского–Лиувилля связывает определитель Вронского двух решений ЛОУ второго порядка с коэффициентом при первой производной:
\[
W(x) = y_1 y_2' - y_1' y_2 = C_0 e^{-\int p(x)\,dx}.
\]
Для нахождения одного конкретного линейно независимого решения можно положить константу $C_0 = 1$.
\end{theory}

\step{1} Запишем уравнение Остроградского–Лиувилля для Вронскиана при $C_0 = 1$:
\[
y_1 y_2' - y_1' y_2 = e^{-\int p(x)\,dx}.
\]

\step{2} Поделим обе части уравнения на $y_1^2(x)$ (предполагая, что на рассматриваемом интервале $y_1(x) \neq 0$):
\[
\frac{y_1 y_2' - y_1' y_2}{y_1^2(x)} = \frac{e^{-\int p(x)\,dx}}{y_1^2(x)}.
\]

\step{3} Заметим, что левая часть представляет собой производную частного двух функций:
\[
\frac{d}{dx} \left( \frac{y_2}{y_1} \right) = \frac{e^{-\int p(x)\,dx}}{y_1^2(x)}.
\]

\step{4} Интегрируем полученное равенство по переменной $x$:
\[
\frac{y_2}{y_1} = \int \frac{e^{-\int p(x)\,dx}}{y_1^2(x)}\,dx + C.
\]
Положим постоянную интегрирования $C = 0$.

\step{5} Умножим обе части на $y_1(x)$, чтобы получить явное выражение для второго решения:
\[
y_2(x) = y_1(x) \int \frac{e^{-\int p(x)\,dx}}{y_1^2(x)}\,dx. \quad\blacksquare
\]

\answer
Формула Абеля выведена напрямую интегрированием соотношения Остроградского–Лиувилля: $y_2(x) = y_1(x) \int \dfrac{e^{-\int p(x)dx}}{y_1^2(x)}\,dx$.

%=============================================================
\task{76}
%=============================================================
\subsection*{Условие}
Найдите общее решение системы матричным методом Эйлера:
\[
\begin{cases}
\dot{x} = 2x + y, \\
\dot{y} = 3x + 4y.
\end{cases}
\]

\begin{theory}
Линейная однородная система ОДУ записывается в матричном виде $\dot{X} = A X$. Метод Эйлера состоит в поиске частных решений вида $X = v e^{\lambda t}$. При этом $\lambda$ является собственным значением матрицы $A$, а $v$ — отвечающим ему собственным вектором. Если собственные значения вещественны и различны ($\lambda_1 \neq \lambda_2$), то общее решение имеет вид:
\[
X(t) = C_1 v^{(1)} e^{\lambda_1 t} + C_2 v^{(2)} e^{\lambda_2 t}.
\]
\end{theory}

\step{1} Запишем матрицу системы:
\[
A = \begin{pmatrix} 2 & 1 \\ 3 & 4 \end{pmatrix}.
\]
Найдём собственные значения матрицы, составив характеристическое уравнение:
\[
\det(A - \lambda I) = \det \begin{pmatrix} 2-\lambda & 1 \\ 3 & 4-\lambda \end{pmatrix} = (2-\lambda)(4-\lambda) - 3 = \lambda^2 - 6\lambda + 5 = 0.
\]
Корни этого уравнения:
\[
(\lambda - 1)(\lambda - 5) = 0 \implies \lambda_1 = 1, \qquad \lambda_2 = 5.
\]

\step{2} Найдём собственный вектор $v^{(1)}$ для $\lambda_1 = 1$. Решим систему $(A - I)v = 0$:
\[
A - I = \begin{pmatrix} 1 & 1 \\ 3 & 3 \end{pmatrix} \sim \begin{pmatrix} 1 & 1 \\ 0 & 0 \end{pmatrix} \implies v_1 + v_2 = 0.
\]
Положим $v_1 = 1$, тогда $v_2 = -1$. Первый собственный вектор: $v^{(1)} = \begin{pmatrix} 1 \\ -1 \end{pmatrix}$.

\step{3} Найдём собственный вектор $v^{(2)}$ для $\lambda_2 = 5$. Решим систему $(A - 5I)v = 0$:
\[
A - 5I = \begin{pmatrix} -3 & 1 \\ 3 & -1 \end{pmatrix} \sim \begin{pmatrix} -3 & 1 \\ 0 & 0 \end{pmatrix} \implies -3v_1 + v_2 = 0.
\]
Положим $v_1 = 1$, тогда $v_2 = 3$. Второй собственный вектор: $v^{(2)} = \begin{pmatrix} 1 \\ 3 \end{pmatrix}$.

\step{4} Запишем общее решение системы в координатной форме:
\[
\begin{pmatrix} x(t) \\ y(t) \end{pmatrix} = C_1 \begin{pmatrix} 1 \\ -1 \end{pmatrix} e^t + C_2 \begin{pmatrix} 1 \\ 3 \end{pmatrix} e^{5t}.
\]

\answer
Общее решение системы: \\
$x(t) = C_1 e^t + C_2 e^{5t}$, \\
$y(t) = -C_1 e^t + 3C_2 e^{5t}$.

%=============================================================
\task{77}
%=============================================================
\subsection*{Условие}
Рассмотрим груз на пружине без трения ($m = 1, k = 1$). Его уравнение $y'' + y = 0$. Перейдите к системе и решите ее методом Эйлера, введя переменные $x = y$ (положение) и $y = v = y'$ (скорость).

\begin{theory}
Пусть $x$ обозначает положение тела, а $y$ — его скорость. Тогда $y = \dot{x}$, а ускорение равно $\dot{y} = \ddot{x} = -x$. Получаем линейную систему ОДУ первого порядка:
\[
\begin{cases}
\dot{x} = y, \\
\dot{y} = -x.
\end{cases}
\]
Для комплексного собственного значения $\lambda = \alpha + i\beta$ с комплексным собственным вектором $v = a + ib$ общее вещественное решение строится как линейная комбинация вещественной и мнимой частей комплексного решения $v e^{\lambda t}$.
\end{theory}

\step{1} Запишем систему в матричной форме с матрицей $A$:
\[
A = \begin{pmatrix} 0 & 1 \\ -1 & 0 \end{pmatrix}.
\]
Составим характеристическое уравнение:
\[
\det(A - \lambda I) = \det \begin{pmatrix} -\lambda & 1 \\ -1 & -\lambda \end{pmatrix} = \lambda^2 + 1 = 0 \implies \lambda = \pm i.
\]

\step{2} Найдём комплексный собственный вектор $v$ для $\lambda_1 = i$:
\[
A - iI = \begin{pmatrix} -i & 1 \\ -1 & -i \end{pmatrix} \implies -i v_1 + v_2 = 0 \implies v_2 = i v_1.
\]
Положим $v_1 = 1$, тогда $v_2 = i$. Получаем собственный вектор: $v = \begin{pmatrix} 1 \\ i \end{pmatrix}$.

\step{3} Составим комплексное решение $Z(t) = v e^{it}$ и выделим его вещественную и мнимую части:
\[
Z(t) = \begin{pmatrix} 1 \\ i \end{pmatrix} (\cos t + i\sin t) = \begin{pmatrix} \cos t + i\sin t \\ -\sin t + i\cos t \end{pmatrix} = \begin{pmatrix} \cos t \\ -\sin t \end{pmatrix} + i \begin{pmatrix} \sin t \\ \cos t \end{pmatrix}.
\]
Получаем два фундаментальных вещественных решения:
\[
X^{(1)}(t) = \begin{pmatrix} \cos t \\ -\sin t \end{pmatrix}, \qquad X^{(2)}(t) = \begin{pmatrix} \sin t \\ \cos t \end{pmatrix}.
\]

\step{4} Запишем общее решение в виде линейной комбинации фундаментальных решений:
\[
\begin{pmatrix} x(t) \\ y(t) \end{pmatrix} = C_1 \begin{pmatrix} \cos t \\ -\sin t \end{pmatrix} + C_2 \begin{pmatrix} \sin t \\ \cos t \end{pmatrix}.
\]

\answer
Общее решение системы (описывающее гармонические колебания): \\
$x(t) = C_1 \cos t + C_2 \sin t$, \\
$y(t) = -C_1 \sin t + C_2 \cos t$.

%=============================================================
\task{78}
%=============================================================
\subsection*{Условие}
Решите систему, описываемую уравнениями:
\[
\begin{cases}
\dot{x} = -x + 4y, \\
\dot{y} = -4x - y.
\end{cases}
\]

\begin{theory}
Матрица системы имеет комплексные сопряжённые собственные значения $\lambda = \alpha \pm i\beta$, что указывает на фокус (устойчивый, если $\alpha < 0$, или неустойчивый, если $\alpha > 0$). Метод Эйлера в этом случае использует вещественные и мнимые части комплексного векторного решения.
\end{theory}

\step{1} Запишем матрицу системы:
\[
A = \begin{pmatrix} -1 & 4 \\ -4 & -1 \end{pmatrix}.
\]
Составим характеристическое уравнение:
\[
\det(A - \lambda I) = \det \begin{pmatrix} -1-\lambda & 4 \\ -4 & -1-\lambda \end{pmatrix} = (-1-\lambda)^2 + 16 = \lambda^2 + 2\lambda + 17 = 0.
\]
Корни уравнения:
\[
\lambda = \frac{-2 \pm \sqrt{4 - 68}}{2} = -1 \pm 4i.
\]

\step{2} Найдем комплексный собственный вектор $v$ для собственного значения $\lambda_1 = -1 + 4i$:
\[
A - \lambda_1 I = \begin{pmatrix} -4i & 4 \\ -4 & -4i \end{pmatrix} \implies -4i v_1 + 4v_2 = 0 \implies v_2 = i v_1.
\]
Положим $v_1 = 1 \implies v_2 = i$. Комплексный собственный вектор: $v = \begin{pmatrix} 1 \\ i \end{pmatrix}$.

\step{3} Составим комплексное векторное решение $Z(t) = v e^{\lambda_1 t}$ и выделим его действительную и мнимую части:
\[
Z(t) = \begin{pmatrix} 1 \\ i \end{pmatrix} e^{-t} \big( \cos(4t) + i\sin(4t) \big) = e^{-t} \begin{pmatrix} \cos(4t) + i\sin(4t) \\ -\sin(4t) + i\cos(4t) \end{pmatrix}
\]
\[
= e^{-t} \begin{pmatrix} \cos(4t) \\ -\sin(4t) \end{pmatrix} + i e^{-t} \begin{pmatrix} \sin(4t) \\ \cos(4t) \end{pmatrix}.
\]

\step{4} Запишем общее вещественное решение как линейную комбинацию этих векторов:
\[
\begin{pmatrix} x(t) \\ y(t) \end{pmatrix} = C_1 e^{-t} \begin{pmatrix} \cos(4t) \\ -\sin(4t) \end{pmatrix} + C_2 e^{-t} \begin{pmatrix} \sin(4t) \\ \cos(4t) \end{pmatrix}.
\]

\answer
Общее решение системы (траектории представляют собой устойчивый фокус): \\
$x(t) = e^{-t} \big( C_1 \cos(4t) + C_2 \sin(4t) \big)$, \\
$y(t) = e^{-t} \big( -C_1 \sin(4t) + C_2 \cos(4t) \big)$.

%=============================================================
\task{79}
%=============================================================
\subsection*{Условие}
Найдите общее решение системы с кратным корнем:
\[
\begin{cases}
\dot{x} = 3x + y, \\
\dot{y} = -x + 5y.
\end{cases}
\]

\begin{theory}
Если характеристическое уравнение линейной системы имеет кратный корень $\lambda_1 = \lambda_2 = \lambda$ с геометрической кратностью 1, то общее решение ищется методом неопределённых коэффициентов в виде:
\[
\begin{pmatrix} x(t) \\ y(t) \end{pmatrix} = \begin{pmatrix} a t + b \\ c t + d \end{pmatrix} e^{\lambda t}.
\]
Подстановка этого вида в исходную систему позволяет связать коэффициенты $a, b, c, d$.
\end{theory}

\step{1} Составим матрицу системы:
\[
A = \begin{pmatrix} 3 & 1 \\ -1 & 5 \end{pmatrix}.
\]
Характеристическое уравнение:
\[
\det(A - \lambda I) = \det \begin{pmatrix} 3-\lambda & 1 \\ -1 & 5-\lambda \end{pmatrix} = (3-\lambda)(5-\lambda) + 1 = \lambda^2 - 8\lambda + 16 = 0.
\]
Получаем кратный корень $\lambda = 4$ алгебраической кратности 2.

\step{2} Запишем решение в виде $x(t) = (at + b)e^{4t}$, $y(t) = (ct + d)e^{4t}$. Дифференцируем обе функции:
\[
\dot{x}(t) = \big( 4at + 4b + a \big)e^{4t}, \qquad \dot{y}(t) = \big( 4ct + 4d + c \big)e^{4t}.
\]

\step{3} Подставим эти производные в исходную систему ОДУ:
\[
\begin{cases}
4at + 4b + a = 3(at + b) + (ct + d) = (3a + c)t + (3b + d), \\
4ct + 4d + c = -(at + b) + 5(ct + d) = (-a + 5c)t + (-b + 5d).
\end{cases}
\]

\step{4} Приравняем коэффициенты при одинаковых степенях $t$:
\begin{itemize}
    \item Для коэффициентов при $t$:
    \[
    \begin{cases}
    4a = 3a + c \implies c = a, \\
    4c = -a + 5c \implies c = a.
    \end{cases}
    \]
    Обозначим $a = c = C_1$ (первая свободная константа).
    \item Для свободных коэффициентов:
    \[
    \begin{cases}
    4b + a = 3b + d \implies d = b + a \implies d = b + C_1, \\
    4d + c = -b + 5d \implies d = b + c \implies d = b + C_1.
    \end{cases}
    \]
    Обозначим $b = C_2$ (вторая свободная константа). Тогда $d = C_1 + C_2$.
\end{itemize}

\step{5} Подставим найденные коэффициенты в исходный вид решения:
\[
x(t) = (C_1 t + C_2)e^{4t}, \qquad y(t) = (C_1 t + C_1 + C_2)e^{4t}.
\]

\answer
Общее решение системы ОДУ с кратным корнем: \\
$x(t) = (C_1 t + C_2) e^{4t}$, \\
$y(t) = (C_1 t + C_1 + C_2) e^{4t}$.

\end{document}